\section{Computational interfaces and composition}
\label{sec:computation}

The first-message result in Corollary~\ref{cor:live-core-next-rekey}
leaves a further choice to account for: an observer can see earlier
message keys and use them to decide which later message to test.
Theorem~\ref{thm:native-online-message} allows that choice while tying
the selected send to the records and original draws that supplied its
root. The confirmation events are evaluated for that same root in
Theorem~\ref{thm:native-selected-authentication}. We connect these three
results by following when each choice becomes fixed in the native
execution and what information remains available afterwards.

\paragraph{The experiment and its selection boundaries.}
Fix public lifetime budgets $G,C\geq1$ for original generation and
encapsulation draws, and $1\le L\le2^{64}-1$ for the selected message chain.
Write $J=[G]\times[C]$. An invocation receives its slot when its entropy is
reserved, including attempts that later fail; replaying or copying retains
that slot. Unused slots remain unused. These lifetime budgets count work
across the execution, including operations no longer resident in storage.

Table~\ref{tab:online-selection-boundaries} locates the online choices.
The target is one accepted logical send under the next rekey of an established
session, with the eligibility predicate in~\eqref{eq:eligibility}.
The live-core interface retains write failures, process loss, exact resumption
and asynchronous confirmation, including sends after local FinalReady.
Complete record acquisitions keep their original scope and save time.

\begin{table}[t]
\caption{Selection boundaries in the native online-message experiment.
The terminal source-pair and cut selection is evaluated against earlier
allocation decisions.}
\label{tab:online-selection-boundaries}
\small
\begin{tabularx}{\textwidth}{@{}L{0.26\textwidth}Y@{}}
\toprule
Object selected & Boundary and retained identity \\
\midrule
Effective contribution tuple & The fixed candidate's original encapsulation
binding identifies the tuple before its first private ContextBound hash at
any protected alias (Section~\ref{sec:contextbound-connection}). \\
Traffic entry and direction & Designate the complete entry before its first
block allocation; equal-input aliases share its material handle
(Section~\ref{sec:joint-continuation}). \\
Message step & Continue or Mark immediately before the step's first private
computation at either endpoint. Earlier keys may influence this choice;
receive-side skips and failed attempts count
(Section~\ref{sec:native-online-chain}). \\
Accepted send, source pair and cut & The terminal selector identifies the
same accepted challenge send, its original reservation handles and a cut.
Eligibility and usable attribution are evaluated on the complete history
(Section~\ref{sec:security-model}). \\
\bottomrule
\end{tabularx}
\end{table}

The target step's paired successor seed may be disclosed after derivation,
and the suffix then uses the real chain computation. The target plaintexts
and payload-bearing delivery remain private to the application interface;
the network has the declared plaintext-free Resume and protocol calls.
Theorem~\ref{thm:native-online-message} bounds the eligible-message advantage
and the fresh authenticated-installation event in this experiment. Its two
terminal tests use the same source and record connection with their own
complete-continuation costs.

\paragraph{Proof dependencies and advantage accounting.}
Table~\ref{tab:computational-dependencies} identifies the interface supplied
by each connection. Each fixed-candidate comparison retains the original
label $j=(g,a)$ and evaluates the same terminal code on its own execution.
Eligibility, accepted-send provenance and acquisition history are recomputed
from that world's ledgers, with the allocation choices in
Table~\ref{tab:online-selection-boundaries}.

\begin{table}[t]
\caption{Shared dependencies of the native computational results.
Each row identifies an interface and its establishing result.}
\label{tab:computational-dependencies}
\small
\begin{tabularx}{\textwidth}{@{}L{0.25\textwidth}Y@{}}
\toprule
Connection & Supplied interface and principal result \\
\midrule
Record exposure & Original sources, immutable versions and acquisition
history; guarded normalization
(Theorem~\ref{thm:guarded-normalization}). \\
Private retained state & Real-to-opaque image and intent correspondence
(Corollary~\ref{cor:real-to-opaque}); candidate-local stopping and
original-score localization (Lemma~\ref{lem:native-local-score}). \\
Hybrid contribution & Original-operation component adapter and
complete-tuple comparison; native single-root use
(Lemma~\ref{lem:cb-native-root-use}); ordinary KEM instantiation
(Proposition~\ref{prop:kem-deferred-challenge}). \\
Application continuation & Joint confirmation and traffic schedule
(Theorem~\ref{thm:joint-message-leaf}); online prefix continuation
(Theorem~\ref{thm:online-chain-leaf}). \\
Source attribution & Executed-slot allocation provenance
(Proposition~\ref{prop:executed-source-attribution}), instantiated through
ordinary hedged ML-DSA
(Theorem~\ref{thm:native-mldsa-reservation-bridge}). \\
\bottomrule
\end{tabularx}
\end{table}

The coefficient of a record comparison follows from the terminal test.
Confidentiality uses a signed score: an output-bit gap of
$\varepsilon$ changes that score by at most $2\varepsilon$, and the
final factor $1/2$ converts it to message advantage. The component
comparison uses both signs, giving $a_j^++a_j^-+2\delta_j$.
Authentication tests one Boolean event. Its record cost is
$\varepsilon_{{\rm rec},j}$ and its component cost is
$a_j^{\rm auth}+\delta_j$, or $a_j^{\rm msg}+\delta_j$ for online
message integrity. Each adapter accounts for its complete execution.
We charge source-attribution failure once in the original game.
Section~\ref{sec:live-core-instantiation} gives the observations for
the live core. The report extension also charges its declared
projection event.

Proposition~\ref{prop:kem-deferred-challenge} bounds each operational
component term by an ordinary ML-KEM IND-CCA advantage and an explicit
prechallenge ciphertext-intersection cost. This substitution retains the
real decapsulation prefix and the original source-selection accounting.

The three native results therefore share the same records and original
sources. Their terminal tests determine the remaining comparisons.
For the first message it is the score in~\eqref{eq:live-core-next-rekey}.
For confirmation and first-message authentication it is the event
in~\eqref{eq:native-selected-authentication}. For a step selected online
we use the score and event in~\eqref{eq:native-online-confidentiality}--%
\eqref{eq:native-online-integrity}. This last pair also uses the native
order of first computations established by
Lemma~\ref{lem:native-online-chain-correspondence}.

We next apply the same source and record reasoning to adaptive descriptor
families for multiple roots and SDK purpose keys. Bootstrap adds the
question of which earlier responder activation a delayed confirmation
authenticates. Its two original-source connections resolve that question
under their respective observation interfaces. The reply-source profile
defines its record-acquisition term in
Section~\ref{sec:bootstrap-authentication}; the initial-source construction
charges the complete joint entry before partitioning its adaptive reply families
(Theorem~\ref{thm:bootstrap-initial-schedule}). Appendix~\ref{app:common-virtual}
retains the common-virtual-game accounting variant.

\subsection{An explicit renderer for private images and intents}
\label{sec:opaque-rendering}

If the protocol asks twice for the head digest of the same saved image,
both answers must agree. In the alias lemma the renderer can look up
the payload, render the image and hash its bytes. The KEM reduction
faces a harder case. Its challenger gives it the public operation outputs and a
shared-secret challenge, while the selected random tape stays private
even when it occurs in a saved image.

We construct the record-management algorithm for this finite image and
intent interface. It keeps immutable handles, the fields it knows and
sealing coins it samples locally. The original-operation adapter supplies
the protocol's public computations through the common interface $F$.
Holding $F$ fixed here lets us account for the image and intent part of
Proposition~\ref{prop:native-local-accounting}.

\paragraph{Typed record interface.}
An opaque block is an immutable fixed-length private value, represented by a
handle and its known length. Blocks may be correlated and may themselves be
retained sealed reservations. A private image recipe contains at least one
such block; an atom-free recipe is fully known. Local image-encryption and
intent-authentication keys are available to the simulator. Fix a public
budget $M$ for all fresh image-sealing invocations, including failed attempts
and atom-free images. Exact copies and exact retries retain their invocation.
The following operations specify the interface.

\begin{enumerate}
  \item \emph{Seal.} Before sampling its nonce, invocation $i$ fixes the
  recipe, its length and the non-nonce header fields. It samples independent
  $N_i\gets\{0,1\}^{n}$ and privately constructs
  \[
    B_i=\mathsf{Hdr}_i\Vert N_i\Vert C_i,
  \]
  where $C_i$ is the specified deterministic local authenticated encryption
  under that nonce. Local encryption and intent authentication use their
  prescribed fixed algorithms without additional shared-SHA3 calls. All
  admitted image formats place the nonce at the same offset. The reply
  consists of the image handle, length and declared non-nonce metadata;
  $N_i$, $C_i$ and the recipe's private blocks remain private.
  \item \emph{Intent.} A complete canonical prefix $P$ fixes the tag,
  binding, predecessor and the image-head value used in this execution.
  Its tag determines its length. The intent bytes are
  \[
    W_{i,P}=P\Vert B_i\Vert
       \mathsf{MAC}_{k_i}(P\Vert B_i),
  \]
  with one fixed authentication key $k_i$ for image $i$. An exact repeat
  retains the same $(i,P)$ entry. Image and intent tags are disjoint.
  \item \emph{Public commitments.} For a designated private image, the
  ideal-commitment experiment returns independent uniform 256-bit values
  cached by $i$ for its head and by $(d,i,P)$ for its intent digests. The
  domain $d$ ranges over the five admitted intent-digest domains. Atom-free
  objects retain their real hashes and ordinary raw-byte outputs. Later
  validation and zero-value checks use the actual replies in that execution.
  \item \emph{Equality and storage.} Equality queries compare complete
  retained image or intent handles. Copy, retain, persist, delete, crash
  and resume update an explicit durable-state directory. Storage choices,
  errors and branch decisions depend on available metadata, lengths,
  commitments and $F$ outputs. A failure after commit preserves that commit.
  \item \emph{Acquisition.} Atom-free objects are rendered from their known
  inputs. For an opaque object, the admitted successful acquisition emits
  a computable event that permanently falsifies the fixed-candidate
  eligibility before its raw bytes are returned. The stopped score is zero,
  as in Corollary~\ref{cor:target-acquisition-stop}. Unavailable-object requests
  retain their ordinary error result.
\end{enumerate}

All other public computation uses these replies, known inputs, $F$, and the
shared public oracles. This is a classical record interface; the adversary
may have quantum private state and quantum access to those public oracles.
The successful-acquisition condition is supplied by the recovery experiment's
existing event rule. For a fixed target session $s$ and cut $c$, acquisition
of an original version $v$ with $\mathsf{subject}(v)=s$ and
$c<\mathsf{save}(v)$ supplies such an event. Its original-record provenance
determines that save time, including when a later aggregate carries the same
version. Lemma~\ref{lem:adaptive-origin-stop} gives a provenance certificate
valid before the final cut is chosen. A separately permitted projection is
described by its own output operation.

\paragraph{Table algorithm.}
The simulator stores an image table indexed by fresh sealing invocation,
an intent table indexed by image and complete prefix, a commitment cache,
and the directory with its provenance and acquisition ledgers. Each new
image entry records the recipe handles, length, known fields and a sampled
nonce. Each new intent entry records $(i,P)$ and its known length. The
simulator samples designated-private commitments on first use and retrieves
them on exact retry; it constructs and hashes atom-free objects normally.
Equality is structural: images compare equal exactly when their indices
match, and intents compare equal exactly when their pairs $(i,P)$ match.
The storage directory determines the state recovered after a crash, including
an acknowledged or uncertain write that actually committed. At the stipulated
permanent-failure acquisition event the computation returns zero. Every other
transition executes the common public code and $F$. Thus the algorithm uses
no selected opaque block bytes.

\begin{lemma}[Structural equality of complete records]
\label{lem:record-structural-equality}
If the nonces of fresh image invocations are pairwise distinct, then
\[
 B_i=B_j\iff i=j,\qquad
 W_{i,P}=W_{j,P'}\iff(i,P)=(j,P').
\]
\end{lemma}
\begin{proof}
One immutable image handle retains one byte string. Two different image
invocations have different nonces at the same explicit offset, so their
complete wires differ. Intent tags first separate different formats. For
one tag the prefix length is fixed: unequal prefixes differ before the
target image, while equal prefixes with different images differ at the
target's nonce. The same image and prefix give the same deterministic MAC
under the fixed key. Disjoint tags also separate images from intents.
This argument permits coincident private block values and ciphertext
components; the explicit nonces distinguish the complete images.
\end{proof}

Let $I$ be the materialized ideal-commitment experiment, which retains real
private bytes and answers equality by comparing them. Let $T$ be the table
algorithm above. Both use the stipulated fixed-candidate stopping rule and
the same public-oracle model and interface $F$.

\begin{theorem}[Opaque image/intent simulation]
\label{thm:opaque-rendering}
For at most $M$ fresh seals with independent $n$-bit nonces, define
\[
 \delta_{\rm nonce}=\min\{1,\tbinom M2\,2^{-n}\}.
\]
For every efficient terminal bit $D$ computed from the public view, available
metadata and event ledger, $F$ outputs and its own independent test coins
(including its retained application challenge bit),
\[
  |\Pr[I:D=1]-\Pr[T:D=1]|\leq\delta_{\rm nonce}.
\]
For a score $Z\in[-1,1]$ computed from the same available inputs,
$|\mathbb E_I Z-\mathbb E_T Z|\leq2\delta_{\rm nonce}$.
\end{theorem}
\begin{proof}
Couple the fresh sealing nonces, private-commitment replies, scheduler choices,
the state of $F$, and the public-oracle transformations. Before the first
nonce collision, maintain agreement of the directories, known fields,
immutable references, public bytes, statuses and semantic ledgers.
Sealing adds the same metadata and nonce to corresponding entries; its
private ciphertext is confined to the materialized entry. Intent creation
adds the same complete prefix and target reference. A fresh private
commitment gets the coupled reply, and a repeat consults the same cache.
Atom-free commitments and acquisitions use identical materialized inputs.
Lemma~\ref{lem:record-structural-equality} gives agreement of equality replies.
The common public computation consequently makes the same transitions and
oracle calls. Copying, committed or failed writes, deletion, crash and
resume preserve the directory invariant. An opaque acquisition emits the
same permanent-failure event before either computation needs to return its
unknown bytes, and both stopped scores are zero.

The next nonce is independent uniform even when the decision to seal is
adaptive. Across at most $M$ invocations, a union bound charges
$\binom M2 2^{-n}$ for any collision. Coupled terminal bits differ only on
that event; bounded scores can differ by at most two. Identical classical
replies and shared-oracle transformations also preserve a quantum
adversary's state at every coupled good prefix. The argument uses no
measurement of a quantum query.
\end{proof}

\paragraph{Connecting real commitments.}
After replacing complete-wire equality by structural equality, let
$\epsilon_{\rm com}$ bound the bit-distinguishing gap between real private
record digests and the independently cached commitments of the typed
interface. Both games run the same public code, ordinary known-record
computations, available metadata and declared stopping rule. The comparison
counts the full distinguishing computation against their shared public oracle.
It is a cryptographic input to the following composition, separate from the
explicit table simulation. Appendix~\ref{app:primitive-inputs} records the
concrete multi-record bound under fixed-input injectivity, per-nonce
confinement and the prescribed local encryption/authentication algorithms.

\begin{corollary}[Real-record to table-algorithm bound]
\label{cor:real-to-opaque}
For the declared fixed-candidate interface and the commitment conditions
above, let $R$ be the real-digest materialized experiment. With $n=192$,
\[
 |\Pr[R:D=1]-\Pr[T:D=1]|
 \leq\min\{1,\epsilon_{\rm com}+\delta_{\rm nonce}\}.
\]
The corresponding $[-1,1]$ score bound is
$\min\{2,2(\epsilon_{\rm com}+\delta_{\rm nonce})\}$.
When the interface exposes no private-wire equality query, its additional
$\delta_{\rm nonce}$ term is zero.
\end{corollary}
\begin{proof}
First replace real byte equality by structural equality, coupling exactly
as above and charging $\delta_{\rm nonce}$. The resulting game has no extra
nonce-dependent equality observation, so apply the stated private-commitment
replacement there. Its ideal form is implemented exactly by the table
algorithm. A triangle inequality proves the bit bound. Applying that bound
to a test that returns one with probability $(1+Z)/2$ proves the score bound.
Without a private-wire equality operation the first replacement is absent.
\end{proof}

The hybrid order starts with constructible real record inputs, replaces
their private commitments, and then uses the table algorithm for the KEM
challenge embedding. Every changed commitment is consumed by that world's
subsequent intent prefixes, checks and outputs. For an adaptive origin-pair
event $E_j$, early stopping is justified only by a computable event making
$E_j$ false for every continuation and every still-admissible final target
choice attributed to $j$. The disjoint-event accounting above retains this
per-pair obligation.

% Independent module; no changes to the frozen manuscript.
% Source: multi_report_theorem.md, reverse-order report composition.
\subsection{Immutable report families and reverse-order composition}
\label{sec:report-composition}

A cleanup report can authenticate a key and a tag from an earlier
report. Replacing the earlier tag alone would leave a later computation
that still needs the private key. We therefore replace protected rows
in reverse order of creation. By the time the reduction reaches a
chosen key, it has already replaced the later dependent rows. The
earlier transcripts were fixed before that key existed.

\paragraph{Finite interface.}
Fix a public budget $N$ for fresh report-key draws and index the actual draws
chronologically.  Failed attempts retain their indices; a new draw after
definite absence receives a new index.  Exact replay of a retained original
report keeps its index, key, and transcript.  Each $K_i$ is an independent
uniform 256-bit value, with no rejection of the zero key.

Let $\mathcal X$ be the source-atom handles whose bytes the subsequent KEM
adapter will not possess.  In this report replacement, their real bytes are
still available to the surrounding source experiment.  Before using $K_i$, the
machine fixes a bounded immutable transcript recipe $m_i$.  It can contain
available literals and protocol outputs, atoms in $\mathcal X$, earlier raw
keys $K_j$ with $j<i$, and previously returned tags.  Its own key and tag do not
occur in $m_i$.  Lengths, dependency sets, and admission errors are computable
from the recipe metadata.  Native code may draw the key before assembling the
transcript, preserving that order and its failures, provided those inputs do
not depend on the draw.  Every direct or computed use of a raw key retains its
origin dependency; returned tag bytes are available values.

Compute an online protection flag before using the fresh key:
\begin{equation}
 P_i=\mathbf 1\!\left[
 \operatorname{deps}_{\mathcal X}(m_i)\ne\varnothing
 \;\lor\;
 \exists j<i:\ P_j=1\ \land\ K_j\in\operatorname{deps}_{K}(m_i)
 \right].
 \label{eq:report-protection}
\end{equation}
This is forward closure under raw-key dependencies.  It need not propagate
through a previously returned tag.  The flag uses origin references, not
actual key equality or the primitive challenger's bit, and is recomputed in
each comparison world's history.

The following operations specify the interface.  Creation computes
$t_i=\operatorname{HMAC}_h(K_i,m_i)$ and retains the actual nonzero-ID check,
encoding errors, and persistence outcomes.  Replay or validation names the
original report and uses that same $m_i$; supplied identifiers can be compared
with its retained tag.  This operation does not submit changed private bytes
as a new transcript under $K_i$.  Acknowledgement, terminal-tag comparison,
copying, and crash recovery preserve their ordinary original-handle semantics.
Private storage uses the explicit recipe tables of the opaque-record
interface, including its already idealized commitments and its separate error
accounting.  Public report bodies and other outputs use available inputs or
their specified application interface.  Protected raw keys have no additional
unlisted computation or observation channel.

The full real experiment $G$ retains each declared fragment/key disclosure.
Let $\mathsf{Bad}_{\rm rep}$ be the first successful outward projection whose
recipe depends on a raw $K_i$ with $P_i=1$, detected by its actual origin map
before returning the dependent bytes.  An unavailable fragment, a terminal
tag alone, or a projection without that dependency does not trigger it.
Define $G^{\rm stop}$ to end with score zero at this event.  This is an
analysis game; a report-key fragment is not thereby declared to be a
selected-session acquisition or a new freshness failure.

\paragraph{One-row bound and hybrids.}
Use one public compression oracle $h$ for all rows and quantum queries.
The one-row experiment draws a uniform hidden 256-bit key $K$, independently
of $h$ and oracle-independent initial auxiliary state. The algorithm performs
its complete counted preparation and chooses one classical transcript $m$ of
length at most $L$ without a raw-key input. Its first row query returns either
$\operatorname{HMAC}_h(K,m)$ or an independent uniform 256-bit tag; at most
$T\ge1$ row calls repeat this same transcript and cached answer. The algorithm
has no additional raw-key, midstate or key-dependent output interface.
It may continue querying the same $h$ and output a bit. The budget $q$ counts
all public compression queries, including preparation, other real rows,
continuation and the terminal test. Independently sampled known keys may
coincide with $K$; their ordinary calls to $h$ remain allowed and counted.

Let $\epsilon_{\rm row}(T,q,L)$ be any upper bound on this experiment's
bit-distinguishing gap. This operational definition is the primitive input to
the composition theorem. Appendix~\ref{app:primitive-inputs} gives the
compression-random-oracle instantiation, its padding convention and query
accounting. A zero answer is retained for the protocol's actual error check.

Define $H_k$, for $0\le k\le N$, by running the stopped experiment with every
unprotected row real, every protected row $i\le k$ real, and every protected
row $i>k$ replaced by its own independent uniform cached tag.  All later
recipes, flags, errors, and selections use that world's actual replies.
Thus $H_N=G^{\rm stop}$ and $H_0=:S^{\rm stop}$.  A zero ideal tag takes the
original error branch; it is never resampled.  An absent invocation makes no
row query.

\begin{theorem}[Reverse-order report composition]
\label{thm:report-reverse-composition}
For each slot $i$, let $L_i,T_i,q_i$ uniformly bound the transcript length,
row calls, and complete ordinary compression work of the reduction below,
including its prefix, future computation, other real rows, and final test.
Set
\[
 \epsilon_{{\rm rep},i}
   =\epsilon_{\rm row}(T_i,q_i,L_i),
 \qquad
 \epsilon_{{\rm rep},\Sigma}
   =\min\!\left\{1,\sum_{i=1}^N\epsilon_{{\rm rep},i}\right\}.
\]
Under the stated interface, every common efficient terminal bit test $D$
satisfies
\[
 \left|\Pr[G^{\rm stop}:D=1]
       -\Pr[S^{\rm stop}:D=1]\right|
 \le\epsilon_{{\rm rep},\Sigma}.
\]
For a terminal score $Z\in[-1,1]$ which is zero on stopping,
\[
 \left|\mathbb E_{G^{\rm stop}}Z
       -\mathbb E_{S^{\rm stop}}Z\right|
 \le2\epsilon_{{\rm rep},\Sigma}.
\]
The tests use the common available view, semantic ledger, protocol-interface
outputs, and independent test coins, including the retained application
challenge bit.  They have no extra access to protected key or transcript bytes.
\end{theorem}

\begin{proof}
Fix $i$ and compare $H_i$ with $H_{i-1}$.  Run the full prefix from setup,
locally sampling all earlier keys and computing their real rows.  Before the
selected row is used, both hybrids have the same prefix rules.  If slot $i$
is absent or unprotected, their behavior agrees and no row challenge is used.
Otherwise install the one-row challenger's key as an opaque
$\mathsf{KeyRef}(i)$.  The transcript $m_i$ is computable: it excludes its own
key, all referenced earlier keys were locally sampled, and the source atoms
are still available at this proof stage.  Submit this actual transcript to
the row oracle when its first MAC call is reached, retain its reply, and run
the original nonzero and persistence checks.

Maintain the original-report, key-origin, recipe, durable-state, and event
tables.  Every available output agrees with the selected hybrid.  Earlier
real rows remain computable because their immutable transcripts predate
$K_i$ and their keys are known.  A future unprotected row cannot depend on
$\mathsf{KeyRef}(i)$: such a dependency would make it protected by
\eqref{eq:report-protection}.  Compute that row with its locally sampled key
and public $h$.  A future protected row has index greater than $i$ and is
already ideal in both hybrids; sample its first tag, cache it by original
invocation, and retain its symbolic recipe without evaluating the unavailable
bytes.  This includes a later report which contains $K_i$ or another earlier
protected key and tag.  Flags and subsequent known fields are recomputed
using the actual returned tags.

For replay and validation, the selected row is called only on its original
$m_i$; preserve each counted recomputation and compare supplied identifiers
with its retained reply.  Copying, failure after a durable commit, and
crash/resume use the same immutable handles and the actual reached storage
state. Each report retains its original key and immutable transcript. Public decisions
use available metadata and common-interface outputs.  Before an outward
protected-key projection, the same ledger event
$\mathsf{Bad}_{\rm rep}$ enters the common stopped state with application
score zero; any allowed terminal-bit postprocessing uses that state.  No
unavailable key bytes are serialized.  Unprotected-key fragments retain their known bytes.
This gives an explicit online simulation throughout the continuation.

The selected real row therefore produces $H_i$, and its ideal counterpart
produces $H_{i-1}$.  All other real keys are sampled as actual independent byte
strings.  If one happens to equal the unknown challenged key, evaluating its
real MAC with public $h$ automatically preserves that equality; those queries
belong in $q_i$.  No comparison with the hidden key or distinct-key conditioning
is used.  Intentional aliases of the same original report instead share its
single row cache.

The reduction executes the adversary's quantum queries to the same $h$ and
introduces no measurement of those queries or oracle-dependent initial advice.
The one-row lemma bounds the adjacent bit difference by
$\epsilon_{{\rm rep},i}$.  Sum over the reverse-order hybrids and cap at one.
For the score, use the allowed randomized terminal test which outputs one
with probability $(1+Z)/2$; its difference is half the score-expectation
difference, giving the displayed factor two.
\end{proof}

\begin{corollary}[Original eligibility and real disclosures]
\label{cor:report-exposure-residual}
For the full real experiment $G$, define
\[
 \eta_{\rm rep}(Z)
 =\mathbb E_G\!\left[|Z|\,
                    \mathbf1_{\mathsf{Bad}_{\rm rep}}\right].
\]
Then
\[
 \left|\mathbb E_GZ-\mathbb E_{S^{\rm stop}}Z\right|
 \le\min\!\left\{2,\eta_{\rm rep}(Z)
                       +2\epsilon_{{\rm rep},\Sigma}\right\},
\]
and hence
\[
 |\mathbb E_GZ|
 \le\min\!\left\{1,\eta_{\rm rep}(Z)
       +2\epsilon_{{\rm rep},\Sigma}
       +|\mathbb E_{S^{\rm stop}}Z|\right\}.
\]
For $Z=\mathbf1_{\mathsf{Eligible}}(-1)^{b\oplus b'}$,
$\eta_{\rm rep}(Z)=
\Pr_G[\mathsf{Eligible}\land\mathsf{Bad}_{\rm rep}]$.
\end{corollary}

\begin{proof}
Couple $G$ and $G^{\rm stop}$ up to the first
$\mathsf{Bad}_{\rm rep}$.  On its complement the runs agree; on that event the
stopped score is zero while $G$ continues with the actual disclosed bytes.
Therefore
\[
 \mathbb E_GZ-\mathbb E_{G^{\rm stop}}Z
   =\mathbb E_G[Z\mathbf1_{\mathsf{Bad}_{\rm rep}}].
\]
Bound its magnitude by $\eta_{\rm rep}(Z)$ and apply
Theorem~\ref{thm:report-reverse-composition}.  No conditioning on eligibility
or independence of later target selection is used.
\end{proof}

\paragraph{Consequences for the record connection.}
If a disclosure permanently rules out the original target, its
residual is zero. A batch report can expose a key fragment without
yielding an intact record from a member session, so we trace what the
fragment actually contains. For example, if the transcript is known
and its key is later disclosed, recomputing HMAC distinguishes the
real tag from an independent one with difference $1-2^{-256}$.
Applying this example to a native batch would require evidence that
the transcript is available as well as the key.

The report replacement precedes KEM embedding.  Afterwards every protected
row needs only its original cached tag, and every unprotected recipe is free
of $\mathcal X$ and protected raw-key dependencies.  The report layer can
therefore use opaque source handles for tag generation and immutable tag
validation.  A native composition first charges the actual eligible-exposure
residual, then uses a private-record correspondence proved for the
explicitly stopped native interface, and finally applies the report hybrids.
This record hop retains its own comparison error and acquisition rules. For
key-only fragments, that correspondence requires an explicit mapping of
their actual provenance and justification of every permanent-ineligibility
premise used by the renderer.
Direct source-atom or application-plaintext outputs retain their own
adapters and acquisition rules. Immutable validation uses the original
transcript and cached tag; arbitrary changed-byte import has a distinct
MAC-query interface. Under these stated interfaces, the explicit reduction
covers dependent immutable report rows.

% Independent manuscript fragment. Requires amsmath, amssymb, amsthm.
\subsection{Candidate-local simulation in the native experiment}
\label{sec:candidate-local-accounting}

Which record can the reduction leave opaque? The answer depends on
which two reservations supplied the accepted send. We therefore localize
the score in the native experiment first: each eligible send with a unique
usable pair is assigned to that pair, ineligible runs contribute zero,
and eligible attribution failures remain in the original game. We can then
compare each candidate using the records, protected report rows
and stopping rule determined by its own reservations.

Let $\mathsf G_{\rm nat}$ be the full native experiment and
$J=[G]\times[C]$ its fixed reservation-slot budget.  Write $E$ for original
eligibility, $A$ for usable unique attribution, and $I\in J$ for the attributed
pair on $A$.  These variables describe the actual accepted logical challenge
send; selection does not transfer its test bit to another send.  For a bounded
application score $S\in[-1,1]$, set
\[
 Z=\mathbf1_E S,\qquad
 Z_j=\mathbf1_{E\land A\land(I=j)}S.
\]

For each $j$, let $U_j$ and $B_j$ be efficient, cumulative prefix events
computed from the available classical provenance ledger.  They are
noninterfering monitors: before stopping they change no native response or
protocol transition and do not inspect an unknown tape or a quantum query.
A \emph{free stop} $U_j$ permanently falsifies $E\land A\land(I=j)$ for every
continuation.  A \emph{charged stop} $B_j$ may instead be a permitted
protected-report-key projection.  Define $\mathsf G_{\rm nat}^{[j]}$ to run
the native code until the first such event, then enter a stopped state with
candidate score zero before returning the monitored bytes.  An event for one
candidate does not stop the original game or another candidate's reduction.

\begin{lemma}[Origin-anchored adaptive stopping]
\label{lem:adaptive-origin-stop}
Fix candidate $j$ and one of its original reservation slots $o$ which has
occurred.  Suppose a successful capture supplies an intact version $v$ with
\[
 \mathsf{subject}(v)=\mathsf{session}(o),\qquad
 \mathsf{res}(o)\le\mathsf{save}(v).
\]
Assume usable attribution to $j$ assigns its member origins to the selected
session, and eligibility places the cut before both reservations.  Then this
capture is a free $U_j$ event, even if the other reservation or challenge send
has not occurred and the final cut has not been selected.
\end{lemma}
\begin{proof}
In any continuation with $E\land A\land(I=j)$, the selected session is
$s=\mathsf{session}(o)$ and the selected cut satisfies
$c<\mathsf{res}(o)\le\mathsf{save}(v)$.  Acquisition has already inserted
$v$ into the cumulative archive. Thus the execution has already acquired a
version of $s$ saved after $c$, contradicting the original acquisition conjunct.
Copies, resealing, deletion, and recovery preserve this contradiction because
the subject, original save event, and acquisition are immutable.  If the other
reservation or challenge send never occurs, the attributed eligible event is
already false.  Decryption and a final cut value are unnecessary for this
certificate.
\end{proof}

Actual retention of $o$ in $v$ supplies the time inequality by reservation/save
causality; an independently recorded later save of the same session also
suffices.  More generally, let $L_s$ be the largest original save time of an
acquired intact version of session $s$, with $L_s=-\infty$ for an empty set.
For an existing pair $(g,a)$, the original temporal eligibility conditions are
exactly
\[
       L_s\le c<\min\{\mathsf{res}(g),\mathsf{res}(a)\}.
\]
Once the interval is empty, the selector can no longer choose an
eligible cut for this pair. Resealing the acquired version preserves
its first-save time and hence $L_s$. A batch-only key fragment or a
public commitment can depend on the target session without containing
an intact version. Neither raises $L_s$ on that account. A cut inside
the interval must also meet the remaining eligibility conditions.

\paragraph{One native residual.}
In the full native run let $B_I$ denote the charged event for its final
attributed pair, used only on $A$, and define
\[
 \eta_{\rm attr}^{\rm nat}=\Pr_{\mathsf G_{\rm nat}}[E\land\neg A],
 \qquad
 \eta_{\rm proj}^{\rm nat}=\Pr_{\mathsf G_{\rm nat}}[E\land A\land B_I].
\]
The two residuals measure disjoint events in one game. In particular,
\[
 \eta_{\rm proj}^{\rm nat}
 =\sum_{j\in J}\Pr_{\mathsf G_{\rm nat}}
       [E\land A\land(I=j)\land B_j].
\]
The projection residual counts the charged event for the selected pair. For
the report construction, $B_j$ is the actual protected-key projection event
for the closure induced by that candidate's source family. A projection
affecting only another candidate therefore contributes nothing to this
residual. The selected projection probability remains explicit in the
native-game bound, and the experiment's original session-acquisition rule
is preserved.

\begin{lemma}[Native score localization with local stops]
\label{lem:native-local-score}
Let $\Psi_j$ be the same terminal code for
$\mathbf1_{E\land A\land(I=j)}S$ in each experiment, with value zero in its
stopped state.  Then
\begin{align}
 \mathbb E_{\mathsf G_{\rm nat}^{[j]}}\Psi_j
   &=\mathbb E_{\mathsf G_{\rm nat}}
           [Z_j\mathbf1_{\neg B_j}], \label{eq:local-stop-identity}\\
 \mathbb E_{\mathsf G_{\rm nat}}Z
   &=\mathbb E_{\mathsf G_{\rm nat}}[Z\mathbf1_{\neg A}]
     +\mathbb E_{\mathsf G_{\rm nat}}[Z\mathbf1_{A\land B_I}]
     +\sum_{j\in J}
       \mathbb E_{\mathsf G_{\rm nat}^{[j]}}\Psi_j.
       \label{eq:native-local-decomposition}
\end{align}
\end{lemma}
\begin{proof}
Until its first local stop, the stopped machine has the native prefix
distribution.  On a run with $B_j$, its candidate score is zero.  On a run
with $U_j$, every native continuation already has $Z_j=0$, whether or not
$B_j$ later occurs.  On all other runs the executions and candidate scores
agree.  This proves~\eqref{eq:local-stop-identity}.  Partition the one native
score by attribution and its unique index; pointwise,
$\sum_j Z_j\mathbf1_{B_j}=Z\mathbf1_{A\land B_I}$.  Substitution proves
\eqref{eq:native-local-decomposition}. This identity uses the final index
only to attribute the completed run. For a quantum adversary the
argument compares identical classical-prefix instruments and their stopped
continuations, retaining the native quantum-state transformations and
oracle-query interface up to the local stop.
\end{proof}

\paragraph{Local cryptographic comparisons.}
For each $j$, let $\mathsf V^{[j]}$ be its specified real-component
continuation after the local source/record/report replacements, and let
$\widetilde{\mathsf V}^{[j]}$ be its corresponding random-component game.
All terminal events are recomputed in their own game by the same score code;
no event from a native transcript is transferred to an altered transcript.
Suppose the actual stopped-interface correspondence gives
\[
 \left|\mathbb E_{\mathsf G_{\rm nat}^{[j]}}\Psi_j
       -\mathbb E_{\mathsf V^{[j]}}\Psi_j\right|\le\rho_j.
\]
This is a local score bound, including its native source correspondence.
A binary record or report bound contributes twice its value to this score
comparison.  The stopping certificate does not itself establish the
correspondence for permitted partial projections.

For the remaining component step, specialize to $S\in\{-1,+1\}$, so that
$\Psi_j\in\{-1,0,+1\}$.  Suppose its actual fixed-pair adapter satisfies,
for $\sigma\in\{-1,+1\}$,
\[
 \left|\Pr_{\mathsf V^{[j]}}[\Psi_j=\sigma]
       -\Pr_{\widetilde{\mathsf V}^{[j]}}[\Psi_j=\sigma]\right|
 \le a_j^\sigma+\delta_j.
\]
Here $a_j^\sigma$ is evaluated for the complete derived distinguisher,
including its shared-oracle work; $\delta_j$ is the selected pair's correctness term. Write
$b_j=\mathbb E_{\widetilde{\mathsf V}^{[j]}}\Psi_j$.

\begin{proposition}[Candidate-local native accounting]
\label{prop:native-local-accounting}
Under the stated local comparisons, the native message-guessing advantage
satisfies
\[
 \operatorname{Adv}_{\rm msg}
 \le\min\!\left\{\frac12,\frac12\!\left(
 \eta_{\rm attr}^{\rm nat}+\eta_{\rm proj}^{\rm nat}
 +\sum_{j\in J}
    \bigl(\rho_j+a_j^++a_j^-+2\delta_j+|b_j|\bigr)
 \right)\right\}.
\]
\end{proposition}
\begin{proof}
The first two terms of~\eqref{eq:native-local-decomposition} have magnitudes
at most the two native residuals.  The local correspondence costs $\rho_j$.
Subtracting the positive- and negative-event comparisons bounds the local
component score change by $a_j^++a_j^-+2\delta_j$.  Sum the candidate terms,
apply the triangle inequality, and use $|Z|\le1$.  Completing invalid runs
by an independent fair guess gives
$\operatorname{Adv}_{\rm msg}=\tfrac12|\mathbb E_{\mathsf G_{\rm nat}}Z|$.
\end{proof}

The ideal continuations may have different adaptive selectors, so each
$|b_j|$ belongs to its own continuation and remains a separate term.
If a common intermediate game and a global
comparison have separately been established, that shared error can instead be
charged once before the common-game accounting. The proposition above uses
the local comparisons available for each candidate.

\subsection{Native constructors and retained-object transitions}
\label{sec:native-correspondence}

To use the record argument for Continuity we must identify what the
implementation retains when a write fails, what it reuses on a retry
and which bytes it shows its caller. The source audit traces those
questions through 37 entries in 27 files. It gives the source ranges,
Git blob identities and SHA-256 digests for the constructors, state
transitions, known-input computations and additional public outputs
discussed here.

\paragraph{Images, intents and commitments.}
\texttt{durable::seal} encodes and bounds the payload before drawing a
24-byte nonce. The four image tags place that nonce at byte offsets
$[80,104)$, followed by the XChaCha20-Poly1305 ciphertext and authentication
tag. The target-bearing \texttt{PendingWrite} formats have a canonical
tag-dependent prefix, the complete original target image, and a deterministic
HMAC under the journal-derived intent key. These layouts supply the fixed
offset, prefix and intact-target premises of
Lemma~\ref{lem:record-structural-equality}. The wrapping key for a reserved
operation and the intent-authentication key use their respective derivation
domains; the simulator can retain these local storage keys independently of
the selected KEM random tape.

The source uses one image-head domain and five target-bearing intent-digest
domains. Their callers feed the actual returned digests into fences,
checkpoints and subsequent validation. Cancellation intents carry known
descriptors instead of a private target image; selection rows and public
archive descriptors likewise use their actual known-input computations.
This variant distinction supplies the real-versus-designated-private split
used by the commitment table.

\paragraph{Reservation and durable identity.}
Both generation and encapsulation paths persist the reserved operation token
before invoking the corresponding KEM computation. An original retry opens
that token and consumes its retained entropy. The write-intent path first
retains the complete pending wire. Applying it compares the retained original
wire, checks the expected and next heads, installs its original target bytes,
and removes the pending entry. An already installed exact target is handled
as the same operation. A sibling-session write can instead seal a new image
containing unchanged target payload bytes, preserving the target's source
while creating a new physical sealing invocation.

\paragraph{Uncertain commit and recovery.}
A caller can learn that a commit is uncertain even when the new
image is already durable. On \texttt{CommitUncertain} the owner closes
and later reopens against the retained image and pending target. The
storage machine therefore keeps the actual durable outcome and resumes
the operation found there. If a later aggregate contains that version,
acquisition still refers to the version's first save.

\paragraph{Report dependencies.}
Table~\ref{tab:native-report-families} identifies the three report families.
The resolution key and pending tag are serialized in an epoch's traffic
entry. A later closure or fanout report can authenticate that complete entry.
This gives the causal raw-key edge used by the reverse-order report argument.
The code's exact retry paths retain an original key and tag; a genuinely new
attempt draws a new key. The nonzero check applies to the returned identifier,
with a real zero-ID failure branch. Pending keys use the full 256-bit space.

\begin{table}[t]
\centering\small
\begin{tabularx}{\linewidth}{@{}p{0.17\linewidth}X p{0.25\linewidth}@{}}
\toprule
Family & Immutable transcript and key placement & Ordinary report result \\
\midrule
Session closure & The report MAC covers the session payload with its own
64-byte report/key tail zeroed. The pending key is retained in that tail.
Earlier resolution keys can occur in the preceding traffic entries.
& Session and epoch progress, counts and loss metadata. \\
Fanout abandonment & The MAC covers batch metadata and complete member
payloads. Its own key is held in a separate batch tail, outside that metadata.
Member payloads can include earlier resolution keys.
& Per-member progress and loss metadata. \\
Closed-epoch resolution & The MAC covers the old epoch's counters, message
identifiers, retained digests and skipped indices, excluding its own key.
Its pending key is retained in the traffic entry.
& Includes unconsumed plaintext for the authorized application. \\
\bottomrule
\end{tabularx}
\caption{Source correspondence of the report families. The application's
authorized plaintext output retains its own audience in the security
experiment. Report-key fragments use their actual record content and
origin mapping.}
\label{tab:native-report-families}
\end{table}

The normal retained-object path reopens a registered image and revalidates
its original semantic state. The interface mapping assigns those validation
calls to known state, the protocol interface $F$, or the immutable report
row. Arbitrary modified storage inputs and separately returned projections
receive their corresponding parse, authentication and output operations in
an extended experiment. In particular, the batch-key record and a member's
session payload have different contents: their acquisition events are
determined separately. These typed boundaries preserve the existing target
predicate while making each native call's role in the reduction explicit.

% Joint continuation increment, v0.4, 2026-11-07.
% Standalone module: requires amsmath, amssymb, amsthm and the usual theorem
% environments. It does not redefine the report-row error parameter.
% Complete proof record: joint_continuation.md, same directory.
\subsection{Joint continuation from an epoch root to a message}
\label{sec:joint-continuation}

Before an application uses its first message key, the same root has
already supplied confirmation keys and traffic material. It also feeds
the next epoch. We follow these uses together so that a comparison
that changes the root also changes every output derived from it. The
current root stays private. The observer may learn later contributions,
successor roots and the next chain seed. We ask what it can learn
about the selected message under those disclosures.

\paragraph{Primitive experiments.}
Let $F(k,m)$ be HMAC-SHA256 over one public quantum compression oracle
$h:\{0,1\}^{512}\times\{0,1\}^{256}\to\{0,1\}^{256}$, in message-block,
chaining-state order. All keys below have 32 bytes. Put
\[
 \Gamma(a)=\bigl(h(I(a),\mathrm{IV}),h(O(a),\mathrm{IV})\bigr),
\]
where $I(a)=(a\Vert0^{32})\oplus\mathtt{36}^{64}$ and
$O(a)=(a\Vert0^{32})\oplus\mathtt{5c}^{64}$.
Initial auxiliary state is independent of $h$ and the designated hidden
uniform key. All oracle-dependent preparation is executed and counted,
together with the complete continuation.

Write $\varepsilon_{\rm Hjoint}(C,T,q,L)$ for an output-bit gap bound between
\[
 \operatorname{Col}(a)=F(a,R),\quad
 \operatorname{Row}(m)=F(R,m)
 \quad\hbox{and}\quad
 V(\Gamma(a)),\ U(m),
\]
where $R$ is hidden uniform, $U,V$ are independent private random functions,
there are at most $C$ classical column calls, $T$ classical row calls of
length at most $L$, and $q$ ordinary compression queries, which may be
quantum. Equal inputs are cached. Distinct salts with equal $\Gamma$ share
their column reply. The ideal here is the independent-table ideal, without
a separately imposed cross-pair repair.

Write $\varepsilon_{\rm Hrow}(T,q,L)$ for the corresponding real-HMAC-row
versus private-random-function bound for one uniform hidden key and
\emph{multiple adaptive messages}. This differs from the immutable report
row experiment. Finally, $\varepsilon_{\rm Hcol}(q)$ bounds one distinct
zero-salt evaluation $F(0,S)$, for hidden uniform $S$, versus an independent
uniform reply; exact repeats reuse it. The retained compression-QROM
arguments give admissible bounds for these operational experiments.

\paragraph{The finite continuation.}
The initial root $R$ is sampled independently of $h$ and the initial
auxiliary state, before any use. It may be presampled and withheld during
counted preparation, which chooses the public session identifier $s$ and
transcript inputs. The previous root and old wrapping keys may be known.
The current root has only the uses
\begin{equation}
 P_0=F(0,R),\qquad P_s=F(s,R),\qquad P_x^+=F(R,S_x^+).
 \label{eq:joint-root-uses}
\end{equation}
There are at most $t_+$ distinct 32-byte successor contributions $S_x^+$.
They are supplied by known local computation, may be disclosed, and may
depend on the available view. The controller may compute and disclose
\[
 R_x^+=\operatorname{Expand}(P_x^+,I_x^+,32)
\]
and run arbitrary known subsequent computation from these values with the
same public oracle. Repeated arguments remain consistent.

Let $D$ and $D_M$ denote the actual 40-byte rekey and 43-byte message domains.
The three confirmation information strings are
\[
\begin{aligned}
 I_r&=D\Vert\texttt{responder-confirmation},\\
 I_f&=D\Vert\texttt{final-confirmation},\\
 I_a&=D\Vert\texttt{receipt-confirmation}.
\end{aligned}
\]
Their keys are $\operatorname{Expand}(P_0,I_d,32)$. The controller uses
these keys in the real MAC, verification and signing algorithms.
At most $M\ge1$ distinct traffic information strings have the native form
\begin{equation}
\begin{split}
 I_T={}&D\Vert\texttt{epoch-traffic/HKDF-SHA256/ChaCha20Poly1305/}\\
 &\Vert\mathit{context}_{32}\Vert\operatorname{BE64}(e)
 \Vert\mathit{transcriptDigest}_{32}.
\end{split}
 \label{eq:joint-traffic-info}
\end{equation}
It may be chosen after the response confirmation and its signature: the
native transcript digest hashes the complete offer and response wires.
Every repeated complete $I_T$ reuses its entry, regardless of endpoint or
wire representation. At the first allocation of one new string $I^*$,
mark that entry and a direction $d^*\in\{1,2\}$ before any of its blocks
are produced. The choice may use the available prefix; a previously
allocated entry cannot be reclassified as fresh. Later equal-info
aliases inherit the same protection. This is a profile for a designated
fresh traffic allocation, without a cross-endpoint wire-uniqueness premise.

The four 32-byte output blocks are
\[
 (s_{I,1},s_{I,2},a_{I,1},a_{I,2})
 =\operatorname{Expand}(P_s,I,128).
\]
The selected seed $s^*=s_{I^*,d^*}$, designated before allocation,
has only the selected first-step use
\begin{equation}
 \begin{split}
 I_M&=D_M\Vert\texttt{chain}\Vert\operatorname{BE64}(0),\\
 P_M&=F(0,s^*),\qquad
 (s',K)=\operatorname{Expand}(P_M,I_M,64).
 \end{split}
 \label{eq:joint-chain-step}
\end{equation}
The opposite seed, confirmation keys, ACK keys, blocks at every distinct
nonmarked information string and next seed $s'$ may be disclosed.
A repeated $I^*$ remains the same protected entry.
Real ACK computation and arbitrary candidate-ACK verification
remain available. Subsequent chain computations use $s'$ honestly.
The selected key $K$ has only its declared application encryption/decryption
use. Current-root, PRK, old selected-seed and message-key bytes have no
separate public return, comparison or hashing interface. Private persistence
uses an already justified opaque interface; original candidate-local stops
and projection accounting remain in force.

\begin{lemma}[Adaptive information and fresh traffic blocks]
\label{lem:joint-expand-blocks}
Replace each PRK-class HMAC row by an independent private random function,
available only through HKDF-Expand. The resulting interface is exactly an
independent block table indexed by the class, complete information string
and block index, with consistent prefixes. This holds for adaptively
chosen information strings. In particular, the first four traffic blocks
are fresh at installation despite the confirmation dependence of $I_T$.
\end{lemma}
\begin{proof}
For one private row $W$, the inputs are
$I\Vert1$ and $T_{I,j-1}\Vert I\Vert j$ for $j\ge2$, where each preceding
block has exactly 32 bytes and the counter has one byte. Different counters
separate inputs. For the same counter, deleting the counter and, when
present, the fixed-width preceding block shows that equality requires equal
$I$. This is true for every preceding output assignment. Equal $(I,j)$
pairs retrieve equal inputs and replies inductively. New pairs therefore
receive independent table blocks, even when $I$ was chosen from prior
outputs and public-oracle computation.

The confirmation information lengths are $62,58,60$ bytes and each traffic
information length is $155$, so their table entries differ even if the
zero and session salts share a PRK. The marked $I^*$ is also distinct from
every earlier traffic entry. Importantly, this exact rewrite removes
the internal parent-row inputs containing a preceding secret block.
Because that row has no other interface, they leave no extra use of the
selected traffic seed. Public quantum-oracle computation is unchanged.
\end{proof}

An available test uses the public view, event ledger, declared locally
computable outputs and its independent test randomness. It cannot inspect
the designated private roots, PRKs, seeds, message key or tables. For the
application experiment below we can fix either challenge bit and then
average; the test is not given a hidden primitive key.

Let $G_R$ be this real schedule with uniform root. In $G_{\rm leaf}$ replace
\eqref{eq:joint-root-uses} by $V(\Gamma(\cdot))$ and $U(\cdot)$, replace each
column-PRK Expand interface by Lemma~\ref{lem:joint-expand-blocks}'s block
table, and replace \eqref{eq:joint-chain-step} by two independent cached
blocks $s',K$. Real confirmation/ACK operations and successor-root
derivations continue from the known replies; all signatures, checks and
later outputs are regenerated within the current game.

\begin{theorem}[A selected message leaf with joint continuation]
\label{thm:joint-message-leaf}
Let $q_R,q_0,q_1,q_s,q_M$ bound the entire ordinary compression-query work
of the five reductions below. Set
\begin{equation}
\begin{split}
 \Delta_{\rm sched}=\min\{1,\;&
 \varepsilon_{\rm Hjoint}(2,t_+,q_R,32)\\
 &+\varepsilon_{\rm Hrow}(3+4M,q_0,188)
  +\varepsilon_{\rm Hrow}(4M,q_1,188)\\
 &+\varepsilon_{\rm Hcol}(q_s)
  +\varepsilon_{\rm Hrow}(2,q_M,89)\}.
\end{split}
 \label{eq:joint-schedule-error}
\end{equation}
For every available terminal bit, the probability gap between $G_R$ and
$G_{\rm leaf}$ is at most $\Delta_{\rm sched}$.
For a common available score in $[-1,1]$, the expectation gap is at most
$2\Delta_{\rm sched}$. The second PRK-row term is omitted when the interface
itself rules out a second class.
\end{theorem}
\begin{proof}
First give every current-root column and row call to the joint primitive
experiment. Their replies suffice to compute all Expands, confirmations,
signatures, traffic and successor roots. In particular, disclosed successor
contributions and roots require no raw current root. This gives the first
term and retains the same public oracle throughout.

Next replace the at most two column-PRK rows in class-allocation order.
Preallocate the zero-salt class privately if necessary; a second class
exists only when $\Gamma(s)\ne\Gamma(0)$. For a designated class, install
the row primitive's hidden uniform key as its first $V$ reply, withheld
until allocation. The identifier is computed from known salts and $h$;
it does not inspect the private reply. Every equal $\Gamma$ reuses this
handle. Answer Expands by querying the challenge row on the actual
successive messages. Their output blocks are known, so real MACs,
verifications, signatures and continuation can all be performed.
Other class keys are sampled locally. All their real HMAC computations,
and all preparation and future work, count as ordinary queries.

Three one-block confirmation Expands and at most $M$ four-block traffic
Expands give input lengths $63,59,61$ and $156,188,188,188$. Thus the
first class has at most $3+4M$ distinct row calls and a second class at most
$4M$. A never-allocated class causes no challenge call. Their total is at
most $3+4M$; the separate worst-case budgets avoid conditioning a
primitive bound on the oracle-dependent class-equality event.

Independent $V$ replies may coincide numerically. No inequality is assumed:
known-key computations are performed through the same $h$, including
accidental overlap with the challenge key, and are counted by the
primitive experiment. An ideal row is the declared hybrid replacement.
Intentional $\Gamma$ aliases share that row in every game.

Apply Lemma~\ref{lem:joint-expand-blocks} exactly. The selected traffic seed
is now a fresh independent block, and its parent no longer exposes an
input containing it. Presample the single-column challenge seed and
install it at the newly marked entry before its blocks are produced.
There is no earlier equal-info allocation; later aliases retain this
protected handle.
Its only raw use is $F(0,s^*)$; every other block is sampled locally.
The next hop therefore makes $P_M$ independently uniform at cost
$\varepsilon_{\rm Hcol}(q_s)$.

Finally install the row challenge key at that fresh $P_M$. The two actual
Expand inputs have lengths $57$ and $89$, giving the final term.
The same exact block-table rewrite produces independent $s',K$.
The next seed may be disclosed and used in all later chain computation;
the reduction has every value it needs except the key reserved to the
application interface. These five comparisons prove the bit bound.
Randomizing a terminal bit with probability $(Z+1)/2$ proves the score
bound. Equivalently one may fix an independent application challenge bit,
apply this bound, and average over that bit.
\end{proof}

The comparison installs an independent key and regenerates the complete
continuation from it. Subsequent ciphertexts retain their dependence on that
key, and the intermediate roots keep the joint distributions specified by
the corresponding hybrids.

\paragraph{Application confidentiality.}
Use a one-key, one-encryption left-or-right AEAD experiment with nonce
$0^{96}$, equal-length messages, common associated data and the same
additional public oracle $h$. Its classical decryption interface returns
real plaintext or failure, except that an exact challenge query returns
an authenticated marker. Let $\varepsilon_{\rm AEAD}$ be its absolute
guessing advantage over $1/2$. An encryption-only profile uses the
corresponding one-use privacy experiment.

The honest application creates the selected send. The network adversary
receives ciphertext, metadata, status, ACKs and real authentication/error
outcomes, and can request plaintext-free $\operatorname{Resume}(id)$.
The actual payload and its retained send-intent digest are opaque
application handles. Delivery and any payload-bearing resolution output
are private to the honest application; the public continuation has no
additional payload-dependent return. An exact challenge reception can
therefore update public state using its authenticated marker and known
metadata. Other ciphertexts use the declared decryption interface.

The send-origin normalization caches the encryption reply at its first
\emph{private computation}, even if that computation precedes durable
commit or network release. Native crash recovery may repeat encryption,
but always on the same original tuple; the reduction retrieves its cache.
The cache includes the marked material's original allocation and
authoritative owner/history, direction and index. Distinct material
allocations are not merged merely because their epoch numbers agree.

\begin{corollary}[First message of a recovered direction]
\label{cor:joint-first-message}
For an event $A$ computable from the available view and ledger without
reading the hidden AEAD challenge bit $b$, let
$Z=\mathbf1_A(2\mathbf1[g=b]-1)$.
Under the application interface above,
\[
 |\mathbb E_{G_R}Z|
 \le 2\Delta_{\rm sched}+2\varepsilon_{\rm AEAD}.
\]
\end{corollary}
\begin{proof}
Theorem~\ref{thm:joint-message-leaf}, applied for each fixed $b$ and
averaged, changes this score by at most $2\Delta_{\rm sched}$.
In $G_{\rm leaf}$ give the selected encryption/decryption operations
to the AEAD challenger and sample every other block locally.
Confirmation, ACK, successor-root and subsequent-chain computation uses
known material and the same public oracle. Cache the first encryption
response; exact challenge reception needs only the marker and metadata.
At the end return $g$ when $A$ holds and a fresh fair guess otherwise.
Its centered success is exactly $\mathbb E Z$, bounded by
$2\varepsilon_{\rm AEAD}$.
\end{proof}

The public application boundary is substantive. For two challenge
plaintexts with different native send-intent digests, publishing the
result of $\operatorname{send\_message}(id,m_0,ad)$ after the challenge
would distinguish a retained result from an intent conflict. No
record acquisition is needed for that test. The network profile uses
Resume while the honest application retains the plaintext-bearing API;
original eligibility is unchanged.

\paragraph{Entry from a recovered contribution.}
The uniform value above is the epoch root \emph{after} root derivation.
For a separately justified uniform combined contribution $S$, known
previous root $r$ and public response core, the actual derivation is
\[
\begin{aligned}
 P&=F(r,S),\qquad R=F(P,I_R\Vert\mathtt{01}),\\
 I_R&=D\Vert\texttt{pending-root/HKDF-SHA256/}\\
    &\quad\Vert\operatorname{hash}(\texttt{response-core},\mathit{core}).
\end{aligned}
\]
Its information string has 97 bytes and its sole Expand message has 98.
Suppose this is one original Extract/Expand, with exact repeats cached,
and $S,P$ have no additional raw-use or disclosure channel.
The retained root-initialization lemma applies with the entire controller
above as continuation, at bit cost $\Delta_{\rm init}$. Its ideal reply
can be presampled and withheld, providing precisely the joint uniform
root initialization required here. It regenerates confirmations,
signatures, traffic and future epochs; it does not replace a root inside
a fixed public transcript.

Let $\mathsf W_{\rm cont}^{[j]}$ denote this uniform-combined-contribution
continuation, with the same candidate-local score code $\Psi_j$, and put
$b_j^{\rm cont}=\mathbb E_{\mathsf W_{\rm cont}^{[j]}}\Psi_j$.
Consequently,
\[
 |b_j^{\rm cont}|\le
 2\Delta_{{\rm init},j}
 +2\Delta_{{\rm sched},j}
 +2\varepsilon_{{\rm AEAD},j}.
\]
The native accounting term
$b_j=\mathbb E_{\widetilde{\mathsf V}^{[j]}}\Psi_j$ belongs instead to
the random-component game, which may still use the real ContextBound
combiner. If its actual complete comparison to
$\mathsf W_{\rm cont}^{[j]}$ has bit gap at most
$\varepsilon_{{\rm CB},j}$, then
\[
 |b_j-b_j^{\rm cont}|\le2\varepsilon_{{\rm CB},j},
 \qquad
 |b_j|\le2\varepsilon_{{\rm CB},j}+|b_j^{\rm cont}|.
\]
This comparison retains the controller, query interfaces and score/event
code, counts all shared-oracle preparation and continuation, and recomputes
events within each game's execution. Its actual tuple and disclosure
conditions must hold. The randomized test with success probability
$(\Psi_j+1)/2$ gives the first inequality; the second is the triangle
inequality. Native correspondence $\rho_j$, component and authentication
terms, and original-game projection residuals remain in their established
locations. This supplies the continuation bound for the specified profile,
including known old roots and allowed later disclosure.

\subsection{Online selection within a finite message chain}
\label{sec:chain-prefix}

Earlier message keys can help an observer decide which step to mark.
Once either endpoint first computes that step, the choice is fixed.
The observer may then learn the paired successor seed and compute the
rest of the chain. We ask whether the marked message remains protected
with both its earlier keys and its suffix available. We use the joint
root, confirmation and traffic interfaces of
Section~\ref{sec:joint-continuation} and bound the number of steps
available for selection.

\paragraph{Prefix interface and selection.}
At the first allocation of a fresh traffic information string, designate
its material identity and a direction as before. Fix a deterministic
lifetime bound $1\le L\le 2^{64}-1$. With $s_0$ the designated traffic
seed, the prefix has the native equations
\begin{equation}
\begin{split}
 I_i&=D_M\Vert\texttt{chain}\Vert\operatorname{BE64}(i),\\
 P_i&=F(0,s_i),\qquad
 (s_{i+1},K_i)=\operatorname{Expand}(P_i,I_i,64),
 \quad 0\le i<L.
\end{split}
\label{eq:prefix-chain}
\end{equation}
Complete material identity, direction and index identify one original
step. Matching endpoint copies, failed computations and resumptions
retrieve that step's original result. Distinct traffic materials keep
their own histories, including when their epoch and index agree.

Immediately before each index's first private step computation, the
controller either continues the prefix or marks that index. This choice
uses its available history, including earlier disclosed message keys,
messages, authentication outcomes and counted public-oracle computation.
The mark is permanent and occurs at most once. Every computation at an
alias participates in the same first-computation ledger; this includes
transaction-private receive-side advancement through skipped indices.
An index already computed at any such alias has passed its selection point.

Before a mark, the seeds $s_i$ and Extract values $P_i$ are private
handles with only the uses in~\eqref{eq:prefix-chain}. The message keys
of steps passed without a mark may be returned and used arbitrarily
within the finite counted computation. At a marked step $\ell$, $K_\ell$
has the single-original-encryption, classical-decryption and exact-retry
interface of Corollary~\ref{cor:joint-first-message}. Its paired successor
$s_{\ell+1}$ becomes an available value, from which subsequent real chain
computation and disclosure may proceed. Private persistence uses the
same justified opaque interface. If the execution ends without a mark,
the selected application event and score are zero.

This timing separates observing earlier keys from choosing an already
observed target. It also fixes when each protected seed can first be
installed in a reduction. The root's other traffic entries, the opposite
direction, real confirmation and ACK checks, and successor-root disclosure
remain available as in the preceding joint continuation.

Let $\Delta_{\rm traffic}$ denote the first three terms of
\eqref{eq:joint-schedule-error}, capped at one, with $q_R,q_0,q_1$
bounding the expanded controller's entire execution. They compare the joint
root uses and the at most two traffic/confirmation PRK rows. For each
$i<L$, let $q_{s,i}$ and $q_{M,i}$ be deterministic upper bounds on
all ordinary compression queries in the corresponding complete
column and row reductions. In particular, these budgets include the
other real chain steps, known earlier keys, verification, subsequent
chain computation, preparation and terminal testing. Define
\begin{equation}
\begin{split}
 \Delta_{\rm prefix}(L)=\min\biggl\{1,\;&\Delta_{\rm traffic}
  +\sum_{i=0}^{L-1}\Bigl(
       \varepsilon_{\rm Hcol}(q_{s,i})\\
 &\hspace{30mm}+\varepsilon_{\rm Hrow}(2,q_{M,i},89)
                         \Bigr)\biggr\}.
\end{split}
\label{eq:prefix-error}
\end{equation}
The unused second traffic-row term is omitted under the same
one-class condition as in Theorem~\ref{thm:joint-message-leaf}.

\begin{theorem}[Online selected chain leaf]
\label{thm:online-chain-leaf}
For the prefix interface above, compare the real uniform-root schedule
with the experiment that replaces the joint root and traffic rows by
their preceding independent tables and gives each prefix step through
the mark two independent cached blocks $(s_{i+1},K_i)$. Computation
strictly after the marked step is real and starts from its available
successor seed. Every available terminal bit has probability gap at
most $\Delta_{\rm prefix}(L)$ between these experiments. Every common
available score in $[-1,1]$ has expectation gap at most
$2\Delta_{\rm prefix}(L)$.
\end{theorem}
\begin{proof}
The first three comparisons of
Theorem~\ref{thm:joint-message-leaf} retain the new controller as their
complete continuation. Their replies suffice to compute every chain
locally, including keys returned before the mark. They therefore cost
$\Delta_{\rm traffic}$. Apply
Lemma~\ref{lem:joint-expand-blocks} exactly to the traffic Expand table.
The designated $s_0$ is an independent private block; its parent row
has no remaining internal input containing that block.

Proceed in fixed index order. Before comparison $i$, every preceding
step reached before or at the mark has its independent block-pair
representation. If an earlier index was marked, step $i$ is already
in the real suffix and this comparison leaves that branch unchanged.
Otherwise the preceding block table, or the traffic table for $i=0$,
allocates a private uniform $s_i$. At that allocation the controller
has already decided whether the predecessor was marked. Hence the
column reduction can install its presampled hidden seed exactly on
the branch that continues the private prefix. On other branches it
uses locally sampled values. It never needs to predict the later mark.

Until step $i$ is first computed, its seed has only an opaque handle.
All available observations, including $K_{i-1}$ when $i>0$, can be
generated without its bytes. At the first computation, give $F(0,s_i)$
to the single-column challenger. Repeated computations use the same
answer. Its returned $P_i$ lets the reduction perform the real Expand,
all later steps, marking decisions and the complete public continuation.
If the execution never reaches this call, the challenge remains unused.
Thus this comparison costs $\varepsilon_{\rm Hcol}(q_{s,i})$ without
conditioning a primitive bound on either reachability or selection.

The fresh independent $P_i$ is then installed as the hidden key of
the row comparison. Its two Expand messages have lengths 57 and 89
bytes for every permitted index. The first reply supplies the input
prefix of the second query, so the reduction can generate both outputs
and the subsequent real continuation. This costs
$\varepsilon_{\rm Hrow}(2,q_{M,i},89)$. Exact block-table normalization
then removes the parent-row input containing $s_{i+1}$, giving an
independent pair. If step $i$ was marked, its successor is locally
available and all later steps remain real. Otherwise that successor
is the private handle for comparison $i+1$.

This construction includes skipped-key computation and failed private
attempts: their first step results are recorded before later publication
or commitment. All aliases reuse the same table entry. Independent
samples may coincide numerically; the comparisons preserve known-key
computations through the common public oracle and charge their queries.
No distinct-value conditioning is introduced.

The $2L$ comparisons use fixed worst-case budgets and recompute the
mark and terminal test within each game. Their triangle inequality,
together with the initial three comparisons, gives the bit bound.
The randomized test with success probability $(Z+1)/2$ gives the
stated score bound.
\end{proof}

\begin{corollary}[Selected message confidentiality after prefix disclosure]
\label{cor:online-chain-message}
Use the private-application and public-Resume boundary of
Corollary~\ref{cor:joint-first-message}, and let $A$ be an available
event independent of direct access to the hidden challenge bit $b$.
Include existence of a mark in $A$, and put
$Z=\mathbf1_A(2\mathbf1[g=b]-1)$. For the declared one-key,
one-encryption AEAD experiment,
\[
 |\mathbb E Z|\le
  2\Delta_{\rm prefix}(L)+2\varepsilon_{\rm AEAD}.
\]
For fresh authenticated installation under that marked message key,
\[
 \Pr[\mathsf{Bad}_{\rm selected\ msg}]
 \le\min\{1,\Delta_{\rm prefix}(L)+\varepsilon_{\rm INT}\}.
\]
Both primitive budgets include the entire adaptive selection and
subsequent continuation.
\end{corollary}
\begin{proof}
In the independent-pair experiment, a step is marked before its pair
is first generated. All earlier keys and marking decisions are
therefore available without $K_\ell$. Install the AEAD challenger's
key at that pair's message-key handle, sample its successor locally,
and compute the real suffix from that successor. Previous message
keys, confirmation keys and ACK keys are locally known. Exact retries
reuse the original encryption response recorded before its durable
outcome. The selected-key decryption interface supplies the same
authenticated marker or real error/plaintext handling as in
Corollary~\ref{cor:joint-first-message}.

For confidentiality, return $g$ on $A$ and a fair guess otherwise;
the centered success is the ideal experiment's $\mathbb E Z$.
For integrity, a new selected-key inbox installation absent from
the encryption ledger supplies the integrity experiment's winning
decryption pair. Apply Theorem~\ref{thm:online-chain-leaf} to the
respective score and Boolean event. This yields the two bounds.
\end{proof}

For $L=1$, the comparison has the same five primitive terms as
the selected first-message theorem. A fixed target index $\ell$
can use $L=\ell+1$; online selection uses a lifetime bound fixed
before execution. The native chain source uses the stated 56-byte
information strings, constant-width index encoding and zero AEAD
nonce, advances a send index once, and reuses exact outbox bytes on
Resume. Transaction-private receive advancement can compute skipped
keys before commitment; failed-transaction resumption retains their
first-computation identity in the proof ledger.
Section~\ref{sec:native-online-chain} uses this comparison to carry the
same choice of message into the native next-rekey experiment.

% New module; existing joint-continuation and primitive appendices are unchanged.
\subsection{From an accepted ContextBound operation to the message continuation}
\label{sec:contextbound-connection}

To enter the joint schedule, we need to justify replacing the combined
contribution by a uniform value. Continuity derives that contribution
from the complete ContextBound tuple, while
the component challenger supplies a real or uniform 32-byte ML-KEM-768
shared secret. We therefore compare the complete ContextBound operation,
retaining the native X25519 computation and its rejection path.
Write $H$ for the whole SHA3-256 random oracle and $h$ for
the independent SHA-256 compression oracle of
Section~\ref{sec:joint-continuation}. Both are available in every
comparison. The query budget includes all oracle-dependent preparation,
including construction of the established starting state, and the
entire continuation.

\paragraph{Complete effective tuples.}
For a byte string $x$, let
$\operatorname{LP}(x)=\operatorname{BE64}(|x|)\Vert x$.
The admitted native ContextBound computation is $S=H(\operatorname{enc}(T))$,
where $\operatorname{enc}$ concatenates the LP encodings of the ten fields
\[
 T=(D_{\rm CB},\mathit{suite},\operatorname{BE32}(v),
 s_M,s_X,c_M,p_M,c_X,p_X,C_{\rm policy}).
\]
Here
\[
\begin{aligned}
 C_{\rm policy}
 &=\operatorname{LP}(D_{\rm policy})
   \Vert\operatorname{LP}(\mathit{policyDigest})
   \Vert\operatorname{LP}(C_{\rm app}),\\
 C_{\rm app}
 &=H\bigl(\operatorname{LP}(D\Vert\texttt{response-kem})
                \Vert\operatorname{LP}(O)\bigr),
\end{aligned}
\]
Here $O$ is the entire retained signed offer, and the policy digest is that of
the verified admitted policy. In the pinned suite the complete encoded input
has 2639 bytes, and the 32-byte $s_M$ field occupies offsets 76--107.
The reservation's trust-root binding, original-slot identifiers, operation
identifiers and journal identities are kept in the provenance ledger.
The hash input consists of the ten fields above; the retained signed offer
contributes through $C_{\rm app}$.

\begin{lemma}[Complete-tuple hidden-component comparison]
\label{lem:cb-hidden-component}
Let $U$ be uniform on $\{0,1\}^{256}$, independent of the initially sampled
oracles and initial auxiliary state. In a protected family replace one
32-byte component field by $U$, and let $m$ contain all nine other actual
fields, including the other component secret and the complete policy
context. The private interface returns
$H(f(U,m))$, where $f$ is the actual ten-field encoding. Exact repeats are
cached by $m$. Assume that the rest of the experiment accesses $U$ only
through this interface or opaque handles and respects its stated stops.
For at most $d$ distinct private inputs and $Q_H$ ordinary queries to $H$,
replacing the private interface by a random function $\mathcal T(m)$
changes every efficient output-bit probability by at most
\begin{equation}
 \varepsilon_{\rm CB}=\Phi_{256}(d,Q_H).
 \label{eq:cb-hidden-cost}
\end{equation}
The inputs $m$ may be adaptive. $Q_H$ counts the complete derived
experiment, including preparation, all ordinary KEM/hash work, continuation
and terminal test. The same independent oracle $h$ is retained.
\end{lemma}
\begin{proof}
Length-prefix parsing recovers all ten actual fields. For fixed $m$ and any
complete string $w$, at most one $u\in\{0,1\}^{256}$ satisfies
$f(u,m)=w$, so its point probability is at most $2^{-256}$. If $m\ne m'$
as effective field tuples, then $f(u,m)\ne f(u,m')$ for every $u$.
Cache repeats before invoking the private interface. These are precisely
the point-probability and all-secret freshness conditions of the
hidden-input theorem stated in Appendix~\ref{app:primitive-inputs}
\cite[Theorem~2]{dfh2022}. It gives the displayed bound for adaptive
private queries. The independent $h$ interface is part of the unchanged
auxiliary computation. Ordinary $H$ queries, including an accidentally
matching known computation, retain their actual oracle responses and are
counted in $Q_H$. The comparison covers the unconditional experiment,
including collisions and executions that end without acceptance.
\end{proof}

Two original calls can supply the same effective tuple. Hashing that
tuple twice gives the same answer with certainty. Assigning independent
ideal answers to the two calls would give agreement with probability
only $2^{-256}$. Conversely, omitting the other shared secret or another
effective field can merge distinct inputs. We therefore use the original
call to track provenance and the complete tuple to cache the hash
answer.

\paragraph{Accepted original component calls.}
Fix $j=(g,a)$ in the original-reservation ledger. The component
adapter assigns its ML-KEM challenge public key only to generation origin
$g$, retaining the locally sampled X25519 scalar, runtime admission and
owner metadata. It obtains its challenge pair $(c_M^*,K_{\rm ch})$ at the
first actual private encapsulation computation of $a$ on the original
retained offer from $g$. It computes the X25519 leg locally and caches the
whole result before subsequent persistence. Exact retries reuse that
challenge; a fresh entropy reservation, including one after an unsuccessful
attempt, receives a new original slot. Copies and reopenings do
not create a new challenge invocation.

The operational component game has two stages. It first supplies its
generated public key and permits ordinary prechallenge decapsulation.
A single challenge is requested at the specified first computation of
$a$; it returns a fresh encapsulation ciphertext and either its real key
or an independent uniform 32-byte key. Thereafter its decapsulation
oracle rejects that challenge ciphertext. Before challenge creation,
selected-key decapsulation follows the prechallenge interface. Any answer
returned before that point remains in the transcript, including one for a
query that happens to equal the later sampled $c_M^*$.

After challenge creation, a PQ decapsulation under the selected generation
owner with ciphertext $c_M^*$ returns the cached $K_{\rm ch}$, regardless
of accompanying traditional ciphertext or context. All other PQ
ciphertexts are routed to the component decapsulation interface. The
traditional decapsulation and its rejection path remain native; a rejected
traditional leg produces no ContextBound call. The challenged PQ
ciphertext is never submitted to the postchallenge CCA oracle. Decapsulation
under any other original key-pair handle uses its locally known secret key,
even when public-key bytes happen to coincide. The component costs below
refer to this two-stage game.

The immutable admitted binding must show that $a$ consumes the original
offer from $g$. An immutable contrary binding makes attribution to $j$
permanently impossible and gives a free stop $U_j$ under
Section~\ref{sec:candidate-local-accounting}. Provenance follows that
original binding. Missing or unaccepted operations otherwise follow their
native behavior. Reveal and projection cases use the established free
stops, charged stops and local record correspondence. Thus the existing
component costs $a_j^+,a_j^-,\delta_j$ retain exactly their meanings.
Correctness and permitted-projection costs remain at their established
comparisons.
More explicitly, the derived distinguisher $D_j^\sigma$ runs the entire
adapter and outputs $\mathbf1[\Psi_j=\sigma]$ for
$\sigma\in\{-1,+1\}$. Its operational advantage is
\[
 a_j^\sigma=\left|\Pr[D_j^\sigma=1\mid\text{real challenge}]
                    -\Pr[D_j^\sigma=1\mid\text{uniform challenge}]\right|.
\]
The real branch matches $\mathsf V^{[j]}$ up to the selected-pair
correctness cost $\delta_j$; the uniform branch is exactly the declared
$\widetilde{\mathsf V}^{[j]}$. This gives the two output-event
inequalities used in Proposition~\ref{prop:native-local-accounting}.

In the resulting random-component game,
$\widetilde{\mathsf V}^{[j]}$, $K_{\rm ch}=U$ is independent uniform.
The reduction for Lemma~\ref{lem:cb-hidden-component} can materialize all
actual native KEM tapes locally, execute their algorithms, and replace
only the designated protected component returns by $U$. These locally
generated tapes implement the adjacent pair of experiments. Ordinary SHA3 work
inside those algorithms is included in $Q_H$; opaque records and the same
stopping code retain their declared interfaces.

\paragraph{A concrete single-root source profile.}
Use the next rekey of one established authoritative session. The two
endpoints initially have the same established root $r$, confirmed epoch
and control parent, with no pending selected rekey. The old root and
wrapping keys may be known. The supported native call surface includes
failed reservations and writes, process loss and exact resumption,
adversarial public packets, old-epoch traffic, asynchronous final/receipt
processing, the marked first-message interface of
Corollary~\ref{cor:joint-first-message}, and subsequent epochs. Selected
private owners are used exclusively through this call surface; read-only
copies remain passive.
This profile specifies the starting state and permitted calls; eligibility
remains the original recovery predicate.

\begin{lemma}[Next-rekey root-use correspondence]
\label{lem:cb-native-root-use}
In this profile, fix $j=(g,a)$ whose immutable original encapsulation
binding consumes the retained original offer from $g$ for this next rekey.
Let $m^*$ be that original encapsulation's effective tuple: its challenged
$c_M^*$, locally generated $c_X^*$, traditional secret, retained offer and
admitted policy determine all nine fields. Every protected
use of that same tuple feeds one pair $(r,I_R)$, where $I_R$ is the complete
97-byte native pending-root information string. They are one original
Extract/Expand and exact repetitions. No raw combined-secret channel is
added by retry, crash or completion. This invariant holds on every
continuation of the matching binding, independently of later eligibility
or attribution.
\end{lemma}
\begin{proof}
The retained generation token proceeds to a fixed signed offer. Its
reconstruction checks the regenerated public material against the saved
offer. At the responder the encapsulation token is bound to the complete
saved offer; a different offer conflicts, and completing or committed
plans return their retained flights. While these tokens are pending,
ordinary message operations do not change the established rekey root.
The response prefix fixes the session, context, epochs, role, control
parent and retained-offer digest. Equality of $m^*$ fixes both complete
ciphertext fields, hence the response core and $I_R$. Both endpoints
therefore use their common old $r$ and this same $I_R$.

Before exposing its response, the responder replaces the active
encapsulation-token stage by the stage retaining the derived root.
The proposer replaces the generation-token stage before root installation.
The final/receipt installation sites are the live updates of the rekey
root. A crash before a transition recovers the same token and old root; a
committed transition recovers the retained result. In either case, recovery
does not resume an old selected token under a newly installed root. This invariant is preserved
by every supported call, authoritative write outcome and recovery step.
The source artifact records the corresponding pinned call sites and the
complete induction. The statement concerns uses in the protected family
and permits independently sampled secrets to coincide numerically.
\end{proof}

\paragraph{Fresh allocation and score semantics.}
Mark this original target tuple $m^*$ before its first actual private hash
computation, at any protected alias call site and including a computation
whose subsequent write fails. If a crash intervenes after component
challenge creation, its cached ciphertext and locally determined traditional
leg still identify $m^*$ before any such first use. An earlier private use
at any alias site already allocates this entry and cannot be ignored at a
later durable acceptance or network release. Other protected tuples do not
become the target. If the selected original $a$ never supplies a usable
target, the original candidate score is zero; no substitute tuple is chosen.
In the private-random-function game,
$S^*=\mathcal T(m^*)$ is a freshly allocated uniform 32-byte value;
replays and aliases retrieve that same entry. Its uniformity holds at
this online allocation, before the future transcript is generated.
Equivalently, the value can be presampled and withheld until allocation.
Other entries of $\mathcal T$ can be sampled by the controller. By
Lemma~\ref{lem:cb-native-root-use}, all selected uses of $S^*$ are the
one Extract/Expand required by Appendix~\ref{app:joint-root-entry}.
The actual continuation is then the controller of
Section~\ref{sec:joint-continuation}, including its fresh traffic-entry
and direction mark, confirmations, ACKs, successors, shared $h$ work and
private retained state.

Each comparison runs the same efficient terminal code $\Psi_j$ on its own
trace. Its value is zero after the established local stops. The qualifying
predicate in the message game must be computable from the available
classical view and provenance ledger, as required by
Corollary~\ref{cor:joint-first-message}. A hidden challenge bit is not an
additional input to that predicate. This condition matters: if the guess
is always zero and a fictitious predicate accepts exactly when the hidden
bit is zero, the qualified signed score has expectation $1/2$ even for an
otherwise bit-independent view. For the admitted predicate, an independent
fair guess on disqualified runs realizes the full qualified score in the
actual AEAD game.

\begin{theorem}[ContextBound connection for the next native rekey]
\label{thm:cb-next-rekey}
Use the concrete profile and adapters above, the local comparisons of
Proposition~\ref{prop:native-local-accounting}, and the message interface
of Corollary~\ref{cor:joint-first-message}. Fix finite budgets before the
experiment. For candidate $j$, let $d_j$ bound its distinct protected
effective tuples and $Q_{H,j}$ its complete ordinary $H$ queries. Let
$\Delta_{{\rm init},j}$, $\Delta_{{\rm sched},j}$ and
$\varepsilon_{{\rm AEAD},j}$ have the existing whole-continuation meanings
of Appendix~\ref{app:joint-root-entry},
Theorem~\ref{thm:joint-message-leaf} and
Corollary~\ref{cor:joint-first-message}. Define
\begin{equation}
 B_j^{\rm cont}=\min\left\{1,
 2\Phi_{256}(d_j,Q_{H,j})+2\Delta_{{\rm init},j}
 +2\Delta_{{\rm sched},j}+2\varepsilon_{{\rm AEAD},j}\right\}.
 \label{eq:cb-continuation-bound}
\end{equation}
Then $|b_j|\le B_j^{\rm cont}$ and
\begin{equation}
\begin{split}
 \operatorname{Adv}_{\rm msg}\le
 \min\!\left\{\frac12,\frac12\left(
 \eta_{\rm attr}^{\rm nat}+\eta_{\rm proj}^{\rm nat}\right.\right.\\[-2pt]
 \left.\left.\qquad+
 \sum_{j\in J}\bigl(\rho_j+a_j^++a_j^-+2\delta_j+B_j^{\rm cont}\bigr)
 \right)\right\}.
\end{split}
 \label{eq:cb-native-bound}
\end{equation}
The residuals are still evaluated in the one original native experiment.
Each candidate's ideal continuation is compared separately.
\end{theorem}
\begin{proof}
Apply Lemma~\ref{lem:cb-hidden-component} to
$\widetilde{\mathsf V}^{[j]}$. Fresh allocation of $m^*$ gives precisely
the uniform-combined-contribution game
$\mathsf W_{\rm cont}^{[j]}$ from
Section~\ref{sec:joint-continuation}; all other private table entries and
the entire native controller remain present. For any
$\Psi_j\in[-1,1]$, a terminal randomized test that outputs one with
probability $(\Psi_j+1)/2$ converts a bit gap $\varepsilon$ into a signed
score gap at most $2\varepsilon$. Thus
\[
 |b_j-b_j^{\rm cont}|\le2\Phi_{256}(d_j,Q_{H,j}).
\]
The native root-use lemma supplies the single-entry hypothesis for the
existing $\Delta_{{\rm init},j}$ comparison. The joint schedule and AEAD
continuation then give
$|b_j^{\rm cont}|\le2\Delta_{{\rm init},j}
 +2\Delta_{{\rm sched},j}+2\varepsilon_{{\rm AEAD},j}$.
All derived distinguishers include the complete preparation and continuation;
no oracle budget resets between comparisons. Combine these inequalities, use
$|b_j|\le1$, and substitute in
Proposition~\ref{prop:native-local-accounting}.
\end{proof}

The theorem completes the next-rekey continuation bound by supplying the
ContextBound comparison from the original component game. Its
source-specific step is the one-entry root-use invariant. Extensions can
reuse the complete-tuple comparison once their corresponding root-entry
interface is established.

\subsection{Ordinary KEM security with a deferred challenge}
\label{sec:kem-ordinary-bridge}

By the time the native adapter reaches its chosen encapsulation, it
may already have answered decapsulation requests. In the ordinary
single-challenge ML-KEM game the challenge ciphertext is fixed from
the outset. Our comparison keeps the real answers to those earlier
requests and charges the event that one of them names the ciphertext
that the adapter will later receive as its challenge.

In the ordinary game, an honest key pair and a fresh encapsulation
$(c^*,K^*)$ are generated, and the adversary receives
$(pk,c^*,K_\beta)$, where $K_0=K^*$ and $K_1$ is independent uniform.
Its classical decapsulation oracle answers every admissible ciphertext
except $c^*$. Write $\operatorname{Adv}^{\rm cca}(B)$ for the absolute
difference of $B$'s output-one probabilities. All public oracles are the
same as in the two-stage game of
Section~\ref{sec:contextbound-connection}, with the complete preparation,
continuation and terminal test included in the adversary's resources.

ML-KEM encapsulation samples $m^*\gets\{0,1\}^{256}$ and encrypts it
with K-PKE using the prescribed derived coins; K-PKE decryption is
deterministic and returns a 32-byte message
\cite[Algorithms 14--17 and 20]{fips203}. Define the message-recovery
failure probability on this honest challenge distribution by
\[
 \delta_{\rm mr}
 =\Pr[\mathsf{K\mbox{-}PKE.Dec}(sk_{\rm PKE},c^*)\ne m^*].
\]
Here failure means that K-PKE decryption does not recover $m^*$.
The component adapters also have a shared-key correctness term, which
accounts for agreement on the KEM key. We keep both events in the
bound.

\begin{proposition}[Deferred-challenge bridge]
\label{prop:kem-deferred-challenge}
Let $D$ be a complete two-stage component distinguisher, making at most
$q^-$ distinct admissible selected-key decapsulation queries before its
challenge request, or before termination if it makes no such request.
Its challenge request may be adaptive. There is an ordinary-game
adversary $B$ such that
\begin{equation}
 \operatorname{Adv}^{\rm op}(D)
 \le \min\{1,\operatorname{Adv}^{\rm cca}(B)
                 +\zeta_{\rm pre}(q^-)\},
 \qquad
 \zeta_{\rm pre}(q^-)
 :=\min\{1,q^-2^{-256}+\delta_{\rm mr}\}.
 \label{eq:kem-deferred-bridge}
\end{equation}
The reduction makes no additional decapsulation or public-oracle queries.
It performs a ciphertext comparison for each decapsulation request and
retains the complete computation of $D$. The statement holds for a
quantum adversary making classical component calls and quantum queries
to the same static public oracles.
\end{proposition}
\begin{proof}
Adversary $B$ retains its challenge ciphertext and key privately, supplies
only $pk$ to $D$, and forwards the prechallenge decapsulation requests.
If a request equals $c^*$, it stops with output zero before submitting
that forbidden request. At $D$'s challenge request it releases the retained
$(c^*,K_\beta)$. Thereafter it implements the stated challenge-ciphertext
exclusion and forwards all other requests. The challenge key is unused
before release. If $D$ terminates without requesting a challenge, $B$
returns that output unless it has already stopped.

For the comparison, run the full, real prechallenge decapsulation process
of $D$, including the answers following a possible hit. The hidden honest
encapsulation can be sampled in advance without changing this process:
fix the static public-oracle functions and key-generation coins, and
observe that neither this prefix nor its quantum-state transformations
uses $m^*$. Evaluating the challenger's encapsulation earlier changes no
oracle function or adversary register. Thus the prefix's classical query
list $\mathcal C^-$ and private quantum state remain independent of
$m^*$ conditional on those fixed functions and key. This couples the
ordinary game's withheld challenge to the two-stage game's freshly
requested challenge; when there is no request, the unused sample has no
public effect.

Let $V$ be the event $c^*\in\mathcal C^-$. Outside message-recovery
failure, this event implies
\[
 m^*\in
 \{\mathsf{K\mbox{-}PKE.Dec}(sk_{\rm PKE},c):
                                      c\in\mathcal C^-\}.
\]
For every fixed full prefix this set has at most $q^-$ elements. Uniformity
of $m^*$ therefore gives
$p:=\Pr[V]\le q^-2^{-256}+\delta_{\rm mr}$ after averaging over the
oracles, key, and prefix. The argument uses the complete real prefix,
not a query distribution conditioned on avoiding earlier hits.
Adversarial timing and repeated queries do not enlarge this set bound.

Until the first hit, $B$ and the two-stage game give identical responses
and apply identical transformations to the adversary's state. On
$\neg V$ their complete terminal bits agree. Before challenge release
the two values of $\beta$ have the same experiment, so $\Pr[V]=p$ in
both branches. Put
$\gamma_\beta=\Pr[D_\beta=1\land V]$. Since $B$ returns zero on $V$,
\[
 \Pr[D_\beta=1]=\Pr[B_\beta=1]+\gamma_\beta,
 \qquad 0\le\gamma_\beta\le p.
\]
Subtract the two identities and use
$|\gamma_0-\gamma_1|\le p$. This proves the single
$\zeta_{\rm pre}$ cost in~\eqref{eq:kem-deferred-bridge}.
No earlier decapsulation answer is rewritten. The reduction tests only
classical ciphertexts and neither measures nor reprograms a public-oracle
query.
\end{proof}

\paragraph{Substitution in the original-source bounds.}
For each already defined component test $D_j^\sigma$, the term
$a_j^\sigma$ is bounded by
$\operatorname{Adv}^{\rm cca}(B_j^\sigma)
 +\zeta_{{\rm pre},j}^{\sigma}$ with that complete test's budget.
The two signed-event tests consequently retain one such cost each.
A Boolean authentication test $a_j^{\rm auth}$ or $a_j^{\rm msg}$
receives one cost, as does each fixed $(r,j,b)$ bootstrap component test
in either source profile. Equivalently, converting a bounded score
$Z\in[-1,1]$ to the randomized bit $(1+Z)/2$ gives the usual factor two.
The fixed original candidate indices and their existing lifetime sums
are unchanged; there is no additional source or message-index guess.

The adapters continue to bind the challenge to the original generation
and encapsulation slots, cache exact retries, and stop before a selected
unknown tape or private owner would be returned. These source and record
conditions are supplied by their preceding correspondence arguments.
The bridge compares their primitive interfaces without changing the
original reservation events, raw-acquisition rules, or later selected
message. Bootstrap's fresh honest-encapsulation and restored-owner
validation monitor remains a separate part of that correspondence.

\subsection{A native live-core instantiation}
\label{sec:live-core-instantiation}

The next-rekey theorem leaves the native record-comparison cost to
instantiate. We now bound that cost for an authoritative live session.
The application supplies the plaintext
to Send and receives the delivered payload. The network can request an
exact Resume, process control and message traffic, obtain ACKs and
inspect statuses and public commitments. The journal follows the
native image and pending intent through write failure, process loss
and reopening.

An adversary can still acquire old or unrelated records. An intact
target-session version first saved after the cut triggers the original
eligibility stop on acquisition. The only partial secret-field replies
are the selected successor root and next chain seed, subject to the
ownership and phase checks of the joint continuation. Other acquisitions
return their declared complete objects. These operations define the
observation interface used here. Appendix~\ref{app:joint-artifacts}
maps each to the native source.

This core obtains all its replies through the record, primitive and
application interfaces already constructed. Its ordinary calls expose
public metadata or declared available values. Consequently its charged
projection event is empty: $B_j=\varnothing$ for each candidate and
$\eta_{\rm proj}^{\rm nat}=0$. The report composition theorem supplies a
separate extension for cleanup and report interfaces with their declared
projection events.

For candidate $j$, let $M_j^{\rm seal}$ bound fresh outer sealing
invocations, including attempts whose subsequent commit fails. Let
$d_{j,i}$ bound the distinct private image/intent digest inputs in row $i$,
and let $Q_{j,i}$ bound the entire ordinary-$H$ computation of its derived
distinguisher. These are deterministic bounds fixed before execution and
cover every corresponding comparison world. Define
\begin{equation}
\begin{split}
 \varepsilon_{{\rm rec},j}
 &=\min\left\{1,
     \sum_{i=1}^{M_j^{\rm seal}}\Phi_{192}(d_{j,i},Q_{j,i})
       +\delta_{{\rm seal},j}\right\},\\
 \delta_{{\rm seal},j}
 &=\min\left\{1,\binom{M_j^{\rm seal}}{2}2^{-192}\right\}.
\end{split}
\label{eq:live-core-record-cost}
\end{equation}
The second term is the private-wire structural-equality comparison of
Corollary~\ref{cor:real-to-opaque}; it is zero when that equality interface
is absent. Copies and exact pending-target application retain their
original sealing handles. Atom-free rows have $d_{j,i}=0$.

\begin{corollary}[Next-rekey confidentiality on the live core]
\label{cor:live-core-next-rekey}
Use the native next-rekey and marked first-message interfaces of
Theorem~\ref{thm:cb-next-rekey} with the live core above, its recorded source
correspondence and finite whole-execution budgets. Then
\begin{equation}
\begin{split}
 \operatorname{Adv}_{\rm msg}\le\min\left\{\frac12,
 \frac12\left(\eta_{\rm attr}^{\rm nat}\right.\right.\\[-2pt]
 \left.\left.\qquad+
 \sum_{j\in J}\left(2\varepsilon_{{\rm rec},j}
      +a_j^++a_j^-+2\delta_j+B_j^{\rm cont}\right)
 \right)\right\}.
\end{split}
\label{eq:live-core-next-rekey}
\end{equation}
Here $B_j^{\rm cont}$ is the explicit ContextBound, root-initialization,
joint-schedule and AEAD score cost in~\eqref{eq:cb-continuation-bound}.
Attribution is evaluated in the original native experiment.
\end{corollary}
\begin{proof}
The live-profile representation invariant follows the original reservation,
saved-version and sealing handles through each native transition. On its
stopped interface, the materialized record machine has exactly the native
outputs and storage outcomes. Corollary~\ref{cor:real-to-opaque} and
Appendix~\ref{app:primitive-inputs} therefore give
$\rho_j\le2\varepsilon_{{\rm rec},j}$. The core's complete captures use
the existing acquisition stop, and its remaining replies use the declared
interfaces, so its charged projection event is empty. Substitute these two
facts into~\eqref{eq:cb-native-bound}.
\end{proof}

\paragraph{Source selection and attribution.}
The attribution term distinguishes two actual interfaces. In the general
experiment, the selector may name any original reservation handles with
the required session, epoch and role metadata. For example, a generation
reservation $g_1$ may be installed in a transaction-private
\texttt{Plan::Key} whose first persistence fails before a complete pending
target or image becomes durable. Reopening the prior
authoritative state permits a new reservation $g_2$, followed by an honest
accepted exchange using $(g_2,a_2)$. Choose
$c<\mathsf{res}(g_1)<\mathsf{res}(g_2)$ and
$c<\mathsf{res}(a_2)$, with the target signing keys and designated entropy
unexposed and with no acquisition of a post-cut target-session version.
A selector naming $(g_1,a_2)$ can pass
the reservation-event and recovery-cut checks while failing usable
attribution. Every control signature and transcript hash in this trace is
honest. This execution is already charged by
$\eta_{\rm attr}^{\rm nat}$ in the general theorem.

To bind selection to the source that actually produced the traffic material,
Section~\ref{sec:executed-source-attribution} defines an executed-slot
selector that registers $g,a$ before their first private KEM computations
return. A retained-plan induction preserves those original slots, and the
complete signed-body and wire bindings associate an accepted allocation
with them. Its signature experiment retains original signing randomizers
and exact resumptions. That selector receives the explicit signature/hash
bound in~\eqref{eq:executed-attribution-cost}; the general metadata selector
in~\eqref{eq:live-core-next-rekey} retains its original attribution term.

\subsection{Confirmation and first-message authentication}
\label{sec:confirmation-traffic}

The recovered root supplies three confirmation keys as well as the
directional traffic material in Theorem~\ref{thm:joint-message-leaf}.
We now use that joint comparison to bound acceptance of a confirmation
core or first-message ciphertext with no earlier honest origin. For
confirmation, the keys remain behind the protocol's actual MAC and
verification calls.
The adversary may query the independent public oracles $H$ and $h$ in
superposition. The protocol requests, MAC inputs, verification requests
and replies are classical.

\paragraph{Native transcripts and application interface.}
Write $\mathcal H_L(x)=H(\operatorname{LP}(D\Vert L)
\Vert\operatorname{LP}(x))$, with $\operatorname{LP}$ the source's
eight-byte length prefix. For $d\in\{r,f,a\}$, let $C_d$ be the responder,
final or receipt confirmation key. The tag is
\begin{equation}
\begin{aligned}
 t_d&=F(C_d,\mathcal H_{L_d}(X_d)),\\
 (L_r,L_f,L_a)&=(\texttt{response-core},\texttt{final-core},
                         \texttt{receipt-core}).
\end{aligned}
\label{eq:native-confirmation-tag}
\end{equation}
The complete responder core contains the control tag, profile, context,
session, current and target epochs, sender role, control parent, digest of
the complete signed offer, and hybrid ciphertext. The final core contains
its corresponding control prefix, the offer and response wire digests,
and the old sending-chain length. The receipt additionally contains the
complete final-wire digest. Their respective source lengths are 1305,
225 and 257 bytes. Every MAC input in~\eqref{eq:native-confirmation-tag}
has 32 bytes.

The native verifier checks the complete prefix, body width and identity
signature before the relevant MAC verification. The public result follows
the actual success, failure, persistence and status transitions. The MAC
comparison may therefore give the reduction the expected tag internally;
the verification caller obtains the native result. Every honest tag
construction enters a ghost origin ledger at its first private
computation, before a subsequent write or release. Verification
computations have their own query ledger. Exact repetitions and matching
endpoint aliases use the same role-key identity and the same original
construction entry. A different signature encoding of an already logged
core retains that core's origin.

For the marked traffic allocation and direction, use the first-message
interface of Corollary~\ref{cor:joint-first-message}. The encryption nonce
is $0^{96}$, and the associated data is
\[
 D_M\Vert\texttt{aead}\Vert\mathit{header}_{93}
       \Vert\operatorname{BE16}(|ad|)\Vert ad.
\]
The header binds session, direction, epoch, index, message identifier and
plaintext length. The first private encryption is recorded by its exact
$(\mathit{AAD},\mathit{ciphertext})$ pair and original send identity.
Retries retrieve its cached result through Resume. Plaintext delivery
belongs to the private application. A fresh authenticated receive means
a successful AEAD check followed by installation of the corresponding
native inbox entry. The source's \texttt{PreviouslyConsumed} result is
a historical consumption status. An existing inbox entry retains its
original authenticated receive-intent digest and plaintext; a cached
return refers to that earlier receive.

Let $\mathsf{Bad}_{\rm conf}$ mean that, before the declared local stop,
a successful native confirmation verification accepts a complete core
with no earlier honest private tag construction under that role-key
identity. Let $\mathsf{Bad}_{\rm msg}$ mean that a fresh authenticated
receive under the selected first-message key installs a pair absent from
its honest encryption ledger. These definitions include failures and
unaccepted executions in the probability space.

\paragraph{Operational bounds.}
Fix a role order $r,f,a$. After exact-input caching, let $T_d$ bound the
distinct 32-byte messages evaluated under $C_d$ by both tag construction
and verification, and let $V_d$ bound verification attempts. Repeated
guesses count in $V_d$. All budgets are fixed before execution. Define
\begin{equation}
 \Delta_{\rm conf}=\min\left\{1,
       \sum_{d\in\{r,f,a\}}
          \varepsilon_{\rm Hrow}(T_d,q_d,32)\right\},
 \qquad V=V_r+V_f+V_a.
\label{eq:confirmation-row-cost}
\end{equation}
Here $q_d$ counts every ordinary $h$ query of that row's complete
distinguisher. In particular, it includes preparation, the target
application, all other real primitives and terminal testing. If $B_c$
bounds all such work outside these three confirmation MAC rows, the
conservative choices
\[
 q_d=B_c+4\sum_{e>d}T_e
\]
are valid: earlier rows are already ideal, later rows remain real, and a
32-byte HMAC evaluation uses at most four compression calls including
its two pads. The row lemma separately charges the replaced row's
private evaluations.

Let $\chi_H(Q)$ bound production of two distinct complete strings with
equal $H$ output by a $Q$-query algorithm. It counts the entire derived
experiment's ordinary and honest $H$ evaluations, including preparation
and the final collision test. For the random oracle with 256-bit output,
an admissible bound is
\[
 \chi_H(Q)=\min\{1,C_{\rm col}(Q+2)^3 2^{-256}\},
\]
with the universal constant of the quantum collision theorem
\cite[Theorem~7]{zhandryCollision2015}. The operational notation retains
that constant.

Finally, let $\varepsilon_{\rm INT}$ bound a quantum adversary's success
in the one-key, at-most-one-encryption ciphertext-integrity experiment
with nonce $0^{96}$, classical decryption and the same independent public
oracles. A successful decryption of a pair absent from the encryption
ledger wins. All repeated decryption attempts are included in its query
budget. The encryption ledger begins at the first private computation.

\begin{proposition}[Joint confirmation and message authentication]
\label{prop:confirmation-message-auth}
Use the controller of Theorem~\ref{thm:joint-message-leaf} and the above
native interfaces. Each $C_d$ is accessed through its tag and verification
service; the selected message key is accessed through its single original
encryption, classical decryption and exact-retry service. Then
\begin{equation}
\begin{split}
 \Pr_{G_R}[\mathsf{Bad}_{\rm conf}\lor\mathsf{Bad}_{\rm msg}]
 \le\min\{1,\;&\Delta_{\rm sched}+\Delta_{\rm conf}
       +V2^{-256}\\
       &+\chi_H(Q_H^{\rm conf})+\varepsilon_{\rm INT}\}.
\end{split}
\label{eq:confirmation-message-auth}
\end{equation}
The same controller retains the qualified confidentiality bound
$|\mathbb E_{G_R}Z|\le2\Delta_{\rm sched}
 +2\varepsilon_{\rm AEAD}$ of
Corollary~\ref{cor:joint-first-message}. Its confirmation and verification
operations are part of that complete continuation.
\end{proposition}
\begin{proof}
Apply Theorem~\ref{thm:joint-message-leaf} to the union event's terminal
bit. This costs $\Delta_{\rm sched}$ once and provides the joint
independent-key experiment $G_{\rm leaf}$. All transcript strings and
subsequent information strings are generated in the current world.

For confirmation, replace the three HMAC rows in the fixed role order by
independent private random functions $W_d$. The adaptive classical row
lemma applies to every tag and verification evaluation; its ordinary
query budgets include the remaining real rows. This costs
$\Delta_{\rm conf}$. Independent key samples may coincide numerically;
the real rows retain the shared compression oracle, and the row
comparisons account for those coincidences.

Consider a second execution that returns failure on every verification
whose digest has no earlier honest tag-construction entry under that
role. The coupling stops at the first internal MAC equality on such an
unissued digest, before subsequent persistence or a public response.
Up to this first hit, the two executions have
the same classical replies and induce the same transformation of the
adversary's quantum private state. In the second execution, the unissued
value $W_d(m)$ is independent of the available prefix: its verification
result is fixed to failure and its table value remains private. Each
of at most $V$ attempted tags equals this value with probability
$2^{-256}$. A union bound therefore bounds the first fresh-digest
success by $V2^{-256}$. It also covers repeated guesses for the same
digest. If a fresh complete core instead has a digest already used by
an honest construction under that role, the two different framed core
strings give a collision for $H$. The reduction records these classical
inputs and outputs the pair; its probability is bounded by
$\chi_H(Q_H^{\rm conf})$.

For message integrity in $G_{\rm leaf}$, sample every key except the
selected message key locally, preserving the actual confirmation, ACK,
successor-root and subsequent-chain computation. Forward the selected
original encryption and every selected-key decryption to the integrity
challenger. Exact retries use the first private encryption cache. A new
authenticated inbox installation without an encryption origin entails a
successful new decryption pair, so its probability is at most
$\varepsilon_{\rm INT}$. Existing inbox entries and historical consumed
statuses use their original metadata. The union bound in
$G_{\rm leaf}$ proves~\eqref{eq:confirmation-message-auth}.

For confidentiality, the AEAD reduction samples the confirmation keys
locally and executes their actual verification services. Thus
Corollary~\ref{cor:joint-first-message} applies with its unchanged
complete-continuation budget. For an independent application challenge
bit, both arguments may first fix that bit and then average.
\end{proof}

To carry this bound into the original recovery experiment,
Section~\ref{sec:native-authentication} places the record, component,
complete-tuple and root-initialization comparisons before the uniform
epoch-root interface. The selected authentication event supplies the
terminal bit for each comparison, whose budget covers its complete
continuation.

\subsection{Executed sources and signing reservations}
\label{sec:executed-source-attribution}

An honestly signed offer identifies a session and epoch, yet those
coordinates can describe two draws if the first save failed before a
token became durable. Dating the contribution therefore requires one
more step: follow the retained plan to the draw whose token reached
the private KEM call. The same plan identifies that draw after a restart.

We assume one authoritative history per endpoint for the selected session,
keyed by the actual signer, session and epoch, with the admitted context
and control parent. Passive copies, exact pending-target reconciliation
and reopening preserve this history. The journal identity is part of the
operation scope, but the signed control prefix omits it. Thus two independently
active journals at the same coordinates would require an additional
ownership rule.

For each local epoch, register its original generation or encapsulation
slot immediately before the first actual private KEM computation. The
registers $\mathsf{ExecG}$ and $\mathsf{ExecA}$ contain the existing
reservation indices in $[G]$ and $[C]$; the sampling events and their cut
tests are unchanged. Exact computation retries use the registered slot.
The \emph{executed-slot selector} chooses the logical challenge send and
cut as before, then uses these two registered handles. A missing register
fails the existing reservation-event test. Registration uses the saved
plan's identity and public bindings, independently of future acceptance.

\begin{lemma}[Stability of executed source slots]
\label{lem:executed-slot-stability}
Each authoritative local epoch has at most one executed generation origin
and one executed encapsulation origin. Failed later writes, process loss
and exact resumptions preserve these origins and their complete bindings.
\end{lemma}
\begin{proof}
Generation runs only after \texttt{Plan::Key} has been saved successfully;
encapsulation likewise requires a successful save of the exact offer and
\texttt{Stage::Encapsulation}. An error on this first save
exits before the KEM call. If the target became durable, recovery obtains
that token; otherwise, the unused reservation can be followed by a new
draw. Note that the reservation token is already authoritative when the
KEM call begins. Subsequent phases store that token or its result and
root. After a failed phase update, recovery reopens an earlier phase
with the same source. The responder also rejects a different complete
offer while its saved plan for that local epoch exists. Completion clears
the plan while advancing the epoch; terminal retirement ends the session.
Induction over these transitions
establishes both invariants without requiring distinct public-key or
ciphertext bytes.
\end{proof}

An honest signing ledger records the actual key, purpose, complete body
and the plan that originated the call at the first private signing invocation,
before any subsequent write or release. The signed string is
\[
 M=\mathit{signatureContext}\Vert\mathit{purpose}_{1}
       \Vert\operatorname{BE32}(|B|)\Vert B.
\]
The ledger includes every honest invocation under the actual key across
public purposes and exact owner aliases. Absence from this actual-key
ledger implies absence from its projected owner-index query ledger.
Let $\mathsf{Bad}_{\rm sig}^{0}$ be a successful target-key verification
of a string absent from this ledger, before the first target signing-key
or signer-file disclosure. An alternative valid signature of an issued
body is an issued-message alias. Let $\mathsf{Bad}_H$ be a collision
between distinct classical framed SHA3 inputs used by the association
checks. Write $\chi_H(Q_H)$ for the complete-query collision bound from
Section~\ref{sec:confirmation-traffic}, counting preparation, honest
computation and witness extraction.

\begin{proposition}[Source association for the executed-slot selector]
\label{prop:executed-source-attribution}
For the controller above and the unchanged recovery eligibility $E$,
usable source association $A$ satisfies
\[
 E\land\neg A\ \Longrightarrow\
       \mathsf{Bad}_{\rm sig}^{0}\lor\mathsf{Bad}_H.
\]
If $\varepsilon_{\rm sig}$ bounds the fresh-body event in the complete
signing execution stopped before returning a target key or signer-file
capture, then
\begin{equation}
 \eta_{\rm attr}^{\rm exec}
 \le\Delta_{\rm attr}
 :=\min\{1,\varepsilon_{\rm sig}+\chi_H(Q_H)\}.
 \label{eq:executed-attribution-cost}
\end{equation}
\end{proposition}
\begin{proof}
Eligibility excludes both target-identity disclosures throughout the
execution. A comparison stopped immediately before the first such
return has the same prefix, including the adversary's quantum-state
transformation. Recording whether a fresh-body verification has already
succeeded gives exactly the event $\mathsf{Bad}_{\rm sig}^{0}$.

Consider a first proposer traffic installation before either bad event.
Its saved offer $O$ and regenerated private owner carry the unique
registered generation $g$. A verified response body has an honest
signing invocation. Its complete context, session, epoch, role and parent
locate the responder's same local epoch, hence its unique encapsulation
$a$. The checked offer-wire digest equals the digest of the offer
retained by $a$. Collision freedom for these two actual framed inputs
makes the offered wires equal. The response therefore carries $a$'s
ciphertext on $g$'s public key with the exact KEM context used by the
proposer. Native regeneration checks the public key against the saved
offer before decapsulation and root confirmation. Source attribution
identifies these original calls and bindings; their component-output
correctness is charged separately through $\delta_j$.

A different valid outer response signature can preserve the same body.
It keeps the source pair and root inputs, while the presented complete
response supplies its actual traffic information and final-flight
binding. The joint schedule already accounts for these information aliases.

A responder installs traffic after checking a final against both its
saved offer and ready response. The final's honest signing invocation
comes from a proposer transition of the preceding kind. Its two checked
wire digests give the same complete offer and response in the absence
of $\mathsf{Bad}_H$, and hence the same $g,a$. Receipt acceptance makes
no additional root allocation. Cached returns and reopened plans preserve
their earlier source association by
Lemma~\ref{lem:executed-slot-stability}. All association tests use available
ledger identities and public strings, so the component adapter can
evaluate them with an unknown selected tape. This proves the implication;
the probability bound follows by a union bound.
\end{proof}

In the failed-save trace of Section~\ref{sec:live-core-instantiation},
both draws stay in the history. The executed generation register resolves
their roles: it names $g_2$ because only $g_2$ reached computation.
A selector using the metadata alone still carries the mismatch into its
attribution accounting.

\paragraph{The signing-coin interface.}
Fix a signature scheme whose internal signing procedure is deterministic
given key, message and randomizer. A reservation-aware fresh-message
experiment provides classical Reserve, Open, Execute and Verify calls.
Reserve samples the randomizer and binds key, operation, purpose
and complete body. Open returns that binding and randomizer at any time.
Execute uses the original tuple, records the message at its first private
computation and caches the exact signature; retries reuse it. Verify wins
on a valid message absent from the corresponding signing-query ledger.
The controlled native calls use exactly the tuples bound by their
reservations; the proposer also checks that its recomputed body matches
the saved body. The hashed reservation binding exposes no arbitrary
tuple-retarget operation.

The owner interface starts with provisioned keys and supplies public keys,
public identity and lifecycle metadata, signing reservations and signatures.
It stops before returning a target key or signer file and provides no earlier
byte-dependent observation of that secret owner file. Note that opening
a signing reservation is a different operation: its record is represented
exactly, without requiring the simulator to know the owner's serialized
key seed. Additional owner-file observations require their own
correspondence cost.

For $U$ independently generated honest signing owners, use the usual
per-owner-index message ledgers. Reopening or reusing an owner preserves
its index. Accidental public-key coincidences remain in the distribution,
and the native admission checks are performed normally. All query and
time bounds cover the complete wrapper. The native pair reserves only
ML-DSA's 32-byte randomizer; its P-256 component uses deterministic
RFC~6979 signing and low-$s$ canonicalization
\cite{rfc6979,fips204}. The pinned implementation routes
\texttt{try\_sign} to RFC~6979 with empty additional data
\cite{rustcryptoEcdsa017}; the ordinary
P-256 game below uses exactly this native signing distribution.

\begin{proposition}[Signature component adapters with retained randomizers]
\label{prop:reservation-signature-components}
For this owner interface, the native pair's reservation-aware fresh-message
success probability satisfies
\[
 \varepsilon_{\rm pair,res}
 \le\varepsilon_{\rm MLDSA,res}^{\rm mu}.
\]
For classical adversaries it also satisfies
\[
 \varepsilon_{\rm pair,res}
 \le\varepsilon_{\rm P256,EUF}^{\rm mu}
 \le\min\{1,U\varepsilon_{\rm P256,EUF}^{\rm su}\},
\]
with the complete simulation time and at most the original Execute-call
budget of signing queries in the single-user reduction.
\end{proposition}
\begin{proof}
For the first inequality generate the P-256 keys locally and route each
ML-DSA reservation and execution to its reservation-aware challenger.
Privately Open each original reservation before materializing its record,
counting these calls, and disclose the known randomizer through the
native observation interface when requested. Fixed-context ML-DSA
signatures come from the challenger; deterministic P-256 signatures and
canonicalization are local. A fresh accepted pair projects to a fresh
ML-DSA message forgery.

For the second inequality generate ML-DSA keys and every reserved
randomizer locally. Obtain the P-256 component from the native
deterministic signing challenger at first private computation, and cache
the pair by original reservation. All reservation records and openings
use known bytes. A repeated body with new randomness can have a new
ML-DSA signature while its deterministic component repeats. A fresh
accepted pair yields a P-256 forgery on the same unqueried message.
Both simulations are exact. Guessing one of the $U$ owner indices gives
the displayed single-user bound; other owners are simulated locally.
\end{proof}

FIPS~204, Algorithm~7, derives internal masking randomness from a secret
key seed, the external randomizer and the message representative
\cite{fips204}. The reservation-aware ML-DSA game records this
external-randomizer interface. Section~\ref{sec:reservation-ordinary-bridge}
connects it to ordinary native unforgeability in the shared random-XOF
model. The ordinary P-256 inequality gives the classical instantiation.
The following separation shows what the extra observation of reserved
coins can change in a generic signature scheme.

\begin{proposition}[A separation induced by signing-coin disclosure]
\label{prop:signature-coin-separation}
Let $\Sigma_0$ be a deterministic EUF-CMA signature scheme whose secret
key encoding $s$ has fixed length $L$. Define
\[
 \operatorname{Sign}_{+}(s,M;r)
   =\bigl(\operatorname{Sign}_0(s,M),r\oplus s\bigr),
 \qquad r\leftarrow\{0,1\}^{L},
\]
and verify the first component and the appendix's fixed length.
Its ordinary fresh-message game has exact transcript-preserving
reductions to and from that of $\Sigma_0$, with their simulation time counted.
One opened signing reservation and its signature reveal $s$.
\end{proposition}
\begin{proof}
For every fixed secret and signing transcript, the appendix of each
fresh call is independent uniform. An ordinary reduction simulates it
with fresh uniform bytes and projects a fresh forgery to $\Sigma_0$;
conversely a base forger ignores the appendices. In the reservation
experiment, XOR of the original randomizer and its appendix recovers
$s$, enabling a fresh signature with the base scheme's correctness
probability.
\end{proof}

The constructed scheme preserves ordinary unforgeability. Its signature
and opened reservation together reveal the signing secret, establishing
a separation between the two experiments. The native ML-DSA result uses
the declared randomizer interface and the bridge to ordinary unforgeability.
For the executed-slot controller with the declared owner interface,
\eqref{eq:executed-attribution-cost} replaces the attribution residual
in Corollary~\ref{cor:live-core-next-rekey}, using the appropriate component
advantage from Proposition~\ref{prop:reservation-signature-components}.
The record, KEM-correctness and joint-continuation terms remain in their
existing positions. The source map, full proof and finite fault/alias
checks are in the supplement indexed in Appendix~\ref{app:source-attribution}.

% Full proof, owner routing, and exact source identities accompany this module.
\subsection{From signing reservations to ordinary unforgeability}
\label{sec:reservation-ordinary-bridge}

The native signer can open a reserved randomizer and later reuse it on
a retry. An ordinary signing query returns a signature, while the
reservation interface also permits the randomizer to be opened. The
bridge must therefore reproduce the opened values and preserve the
dependence between repeated executions. We recognize those repeats by
the pair consisting of that randomizer and the complete public
message representative. For ML-DSA, the argument also accounts for the
oracle shared by key generation and hedging.

Use the owner interface of
Proposition~\ref{prop:reservation-signature-components}. A fresh owner
generation receives a family index; exact references and imported copies
use the same index. Independently generated owners remain separate when
their public keys coincide. The family ledger records each original
message at its first private invocation, including failed publication.
It is contained in the native actual-key ledger. Stopping before a
target-family key or signer-file return therefore gives an auxiliary
fresh-message event containing the native stopped event. All routing and
stopping decisions use declared lineage and public metadata, without
testing hidden secret equality.

\paragraph{Independent hedging.}
Suppose signing factors as
\[
 \operatorname{Sign}(s,K,M;r)
 =S\bigl(s,d_{pk}(M);F_K(r,d_{pk}(M))\bigr),
\]
where the complete computation after the public descriptor $d_{pk}(M)$
is deterministic given its internal tape, and $K$ is generated
independently of the remaining owner state. The ordinary signer draws a
fresh hidden $r\in\{0,1\}^n$ per call. The reservation interface reveals
original randomizers before or after execution and retains the complete
original-message ledger. The joint PRF experiment supplies the same public
oracle access as the signature game, including quantum queries when
present; its private function calls are classical and its ideal random
function is independent of the public companion oracle.

\begin{proposition}[An ordinary-signature bridge for independent hedging]
\label{prop:independent-hedging-bridge}
For one protected family and at most $q$ distinct executed
$(r,d_{pk}(M))$ pairs, there are an ordinary wrapper-EUF adversary $B$ and
joint-PRF distinguishers $D_0,D_1$ such that
\begin{equation}
 p_{\rm res}(A)\le p_{\rm EUF}(B)
   +\operatorname{Adv}^{\rm jprf}_F(D_0)
   +\operatorname{Adv}^{\rm jprf}_F(D_1)
   +\beta_n(q),\qquad
 \beta_n(q)=\min\{1,\tbinom q2 2^{-n}\}.
 \label{eq:independent-hedging-bridge}
\end{equation}
Each primitive interface receives at most $q$ calls. The bounds include
the complete simulation, public descriptor work and owner observations.
\end{proposition}
\begin{proof}
Replace the reservation signer's hedging function by a private random
function, using independence to simulate the remaining key state. This
costs the first joint-PRF advantage. Simulator $B$ samples the visible
reservation randomizers and caches signatures by the full pair
$(r,d_{pk}(M))$. For a new pair, it queries the ordinary signer on
$M$; for a repeated pair, it returns the cached signature. Note that
every slot records its message in the local ledger at its first private
execution, even if publication subsequently fails.

Replace the ordinary challenger's hedge by a private random function.
Until its fresh hidden input repeats, its internal tapes are independent
uniform values and can be coupled with the tapes of the distinct visible
pairs. The $t$-th hidden draw repeats a prior input with probability at
most $(t-1)2^{-n}$, giving $\beta_n(q)$. This charge concerns the ordinary
challenger's draws. Arbitrarily many opened reservations and adaptively
selected visible collisions are represented exactly by the pair cache.
Returning the ordinary challenger to its real hedge costs the second
joint-PRF term. Its signing queries are a subset of the executed-message
ledger, so a fresh output is an ordinary forgery. Descriptor aliases need
no separate charge: factorization gives the same signature, while the
ledger records every original message.
\end{proof}

For several protected families with a global stop before any of their
key/file disclosures, the same argument uses indexed multiuser primitive
games and the birthday term
$\min\{1,2^{-n}\sum_h\binom{q_h}{2}\}$, with no extra user factor.
Permitting other-owner disclosures instead uses the selected-family
construction below. All displayed query caps hold across the relevant
complete hybrids; an Execute-call cap is always a valid conservative
choice for $q$.

\paragraph{The shared ML-DSA oracle.}
FIPS~204 key generation expands a uniform 32-byte seed into consecutive
blocks $(\rho,\rho',K)$ of lengths $(32,64,32)$ bytes. Signing computes
\[
 \mu=H(\mathit{tr}\Vert M',64),\qquad
 c=H(K\Vert r\Vert\mu,64).
\]
The remaining signing computation uses the other key components, $\mu$
and $c$~\cite{fips204}. The pinned native backend passes the reserved
32-byte randomizer to \texttt{try\_sign\_with\_seed}; the checked implementation
uses $K$ only for this hedge in the signature computation. The full source
map records the fixed crate archive and its upstream commit.

Model SHAKE256 as a public random XOF $H$, with one independent infinite
output string per complete input and consistent replies to every requested
prefix. Other public algorithms retain the same interfaces in the
comparison and are independent of the sampled random-XOF table. For a
selected ML-DSA-65 family, let $q_*^{\rm res}$ and
$q_*^{\rm ord}$ bound all remaining calls that can address a 34-byte or
128-byte input in the two complete experiments, after rewriting every
designated target seed expansion and hedge as in the proof. Count public
adversarial queries, preparation, verification and all other-owner work.
Define
\[
 \delta_{\rm C}(v)=\min\{1,v2^{-256}\},\qquad
 \delta_{\rm Q}(v)=\min\{1,2v2^{-128}\}.
\]

\begin{theorem}[Native ML-DSA reservation bridge]
\label{thm:native-mldsa-reservation-bridge}
For the declared owner interface and a selected native ML-DSA-65 family,
let $q$ bound the distinct executed target $(r,\mu)$ pairs in every
relevant experiment. There is an ordinary hedged-ML-DSA EUF adversary $B$ with at most $q$
signing queries such that
\begin{equation}
 p_{\rm res}^{H}(A)
 \le p_{\rm EUF}^{H}(B)+\beta_{256}(q)
       +\delta_{\mathsf X}(q_*^{\rm res})
       +\delta_{\mathsf X}(q_*^{\rm ord}).
 \label{eq:native-mldsa-reservation-bridge}
\end{equation}
Here $\mathsf X=\mathsf C$ denotes the classical random-oracle model and
$\mathsf X=\mathsf Q$ the quantum random-oracle model. Reservation and
signing calls are classical in both. Non-target owners may be generated
locally and exposed; the target interface stops before key or signer-file
return and has no earlier secret-byte observation. Both sides use the
same ordinary native algorithm and public-oracle model, with their full
simulation time and oracle budgets.
\end{theorem}
\begin{proof}
Sample independent uniform $\xi,\rho,\rho',K$ of the specified lengths,
a random XOF $H_0$, and a private random function
$R:(r,\mu)\mapsto\{0,1\}^{512}$. Let
\[
 g=\xi\Vert\operatorname{byte}(6)\Vert\operatorname{byte}(5),\qquad
 \mathcal S_K=\{K\Vert r\Vert\mu:\ |r|=32,\ |\mu|=64\},
\]
where these lengths are in bytes. Define $H_*$ by replacing the first
128 output bytes at $g$ with $\rho\Vert\rho'\Vert K$, and the first
64 bytes at each $K\Vert r\Vert\mu$ with $R(r,\mu)$. Preserve every
remaining output bit from $H_0$. The two input sets have disjoint lengths
34 and 128. Conditional on $\xi$ and the keygen output, all other full XOF strings
remain independent uniform, including those in the selected slice.
Thus $H_*$ and the keygen blocks have exactly the native joint distribution.

Rewrite every target-lineage seed expansion, including reopened owners,
to use those sampled blocks, and every target hedge, including retries,
to use $R$. These are exact substitutions under $H_*$. Change every
remaining oracle call to $H_0$. In this ideal experiment the query-register
distribution is independent of the now-unused $\xi,K$. The remaining
base state may depend on $\rho,\rho',H_0$, and every non-target owner
runs its unchanged algorithm with the same shared-oracle calls.

For any ideal query, its expected mass on
$\mathcal D=\{g\}\cup\mathcal S_K$ is at most $2^{-256}$: each relevant
length sector has that hidden-point or hidden-prefix probability, and the
two sector masses sum to at most one. A classical identical-until-hit
argument gives $q_*2^{-256}$. For a quantum query, the difference between
the two XOR oracle actions has norm at most $2\sqrt{p_t(\mathcal D)}$.
Telescoping using the ideal prefix states and applying Jensen's inequality
gives expected final-state distance at most $2q_*2^{-128}$, which also
bounds the difference in a final event probability
\cite{bennett1997strengths}. This comparison includes verification.

Apply it first to the reservation experiment. In the resulting ideal
experiment $R$ is independent of the public oracle, so the pair-cache
coupling from Proposition~\ref{prop:independent-hedging-bridge} gives the
ordinary ideal experiment with loss $\beta_{256}(q)$. Apply the same
oracle comparison in the reverse direction to the complete ordinary
experiment. Its signer is now the native ordinary challenger, giving
\eqref{eq:native-mldsa-reservation-bridge}. The actual reduction $B$, however,
uses only that signing oracle and public interfaces; the statistical comparison
does not require it to know or program $\xi$ or $K$.
\end{proof}

The source inventory makes the query budget concrete. After the designated
rewrites, native ML-DSA-65 uses SHAKE256 input lengths 66 for secret-vector
and mask expansion, 48 for challenge-to-ball sampling, 832 for commitment
hashes, and 1952 for the public-key hash. Matrix expansion uses SHAKE128.
The protocol's fixed context is 42 bytes; its complete descriptor input
has length $155+|B|$ bytes. These target computations therefore agree
exactly under $H_*$ and $H_0$, including every rejection-loop iteration.
Arbitrary public SHAKE256 calls and non-target seed/hedge calls remain in
$q_*$. Safe-domain calls are common operations in the query hybrid. A
finite computation cutoff followed by a limit proves the same statistical
bound for unbounded safe-domain rejection loops; full running time remains
part of the ordinary-EUF term.

\paragraph{Adaptive owners and the pair signature.}
For at most $U$ fresh owner families, guess the target's creation index
uniformly and use that index across aliases. Embed the single-user
challenger there and generate all other owners locally, answering their
permitted disclosures. This works for adaptive creation and gives
\begin{equation}
 \varepsilon_{\rm pair,res}
 \le\min\left\{1,U\left(
       \varepsilon_{\rm MLDSA,EUF}^{\rm su}
       +\beta_{256}(q)
       +\delta_{\mathsf X}(q_*^{\rm res})
       +\delta_{\mathsf X}(q_*^{\rm ord})\right)\right\}.
 \label{eq:native-pair-ordinary-mldsa}
\end{equation}
The ML-DSA oracle ledger is a subset of the selected native pair owner's
ledger, hence of its actual-key ledger. No aggregation of distinct pair
owners with coincident ML-DSA components is assumed. Their generation
and disclosure follow the same distribution as before. Cross-owner seed or
hedge-prefix coincidences are among the remaining shared-oracle queries
already counted in the selected-target comparison.

Equation~\eqref{eq:native-pair-ordinary-mldsa} supplies an ordinary-ML-DSA
instantiation of the signature term in
Proposition~\ref{prop:executed-source-attribution}, preserving its record,
KEM-correctness and joint-continuation accounting.

\subsection{Native authentication for the selected recovered root}
\label{sec:native-authentication}

Suppose the selected send identifies a recovered root. Which earlier
confirmation or first-message installation can we attribute to honest
work under that root? We record the authentication event in the native
history and compare its Boolean outcome through the same sequence of
games. This gives an authentication bound for the root whose first
message the confidentiality experiment tests.

\paragraph{Target event.}
Fix the live-core, single-root next-rekey profile and an admitted marked
controller of Theorem~\ref{thm:cb-next-rekey} and
Proposition~\ref{prop:confirmation-message-auth}.  In particular, for each
fixed candidate its target contribution and marked traffic entry are
designated at the allocation points required by those interfaces; the
terminal selector corresponds to those designated entries whenever
it returns that eligible candidate.  As in
Definition~\ref{def:recovery-game} and its adaptive version, the experiment
has one application challenge.  Its accepted logical send fixes
the target root schedule through the send's stored provenance handle.  The
terminal selector may choose coordinates, reservation handles, and the cut
adaptively; eligibility requires those coordinates to identify that same
accepted send.  Let $E$ be its native eligibility event, $A$ its
source-attribution event, and $I\in J=[G]\times[C]$ the attributed pair of
original generation and encapsulation reservations on $A$.

The target root and its event history are defined by the actual accepted
challenge send, including on executions where $A$ fails.  A challenge without
an accepted send has $E=0$ and we set the target event flag to zero.  On $A$,
the root handle is the schedule entry generated from the effective
ContextBound tuple $m^*$ of the attributed pair.  The handle is carried
through replay, recovery, exact tuple aliases, and the resulting role keys
and marked first-message key.  Ledger aliasing follows
Proposition~\ref{prop:confirmation-message-auth}.

For this target, let $\mathsf{Bad}_{\rm auth}$ be the union of the following
events: a successful native response, final, or receipt confirmation whose
complete core is absent from the earlier honest private-tag ledger under the
corresponding target role key; or a fresh authenticated first-message receive
and inbox installation whose $(\mathrm{AAD},\mathrm{ciphertext})$ pair is absent
from the earlier honest private-encryption ledger under the marked target
first-message key.  These are the events of the uniform-root authentication
proposition restricted to this selected root schedule.  The honest ledgers
are populated at the first private computation, before persistence or public
release.  Repeated deliveries and previously consumed inbox results refer
back to their original ledger entries.  Genuine aliases of the selected
schedule entry share its honest ledger, while independently allocated
entries have their own provenance.

The experiment records these event flags as execution proceeds.  Its eventual
evaluation reads the accepted challenge send's root history, which may
include confirmation events before that send.  For a fixed pair $j=(g,a)$,
$a$ determines $m^*$ before its first private hash computation
at any protected alias, including a computation followed by a failed write,
as specified in Section~\ref{sec:contextbound-connection}.  Other protected
tuples have their own entries.  If $a$ never supplies a usable
target, the candidate event is zero.  The prescribed online marks for the
root entry, traffic entry, and direction remain fixed.  These rules connect the
recorded history to the admitted controller throughout the comparisons.
The probability bounded below is the joint probability
$\Pr[E\wedge\mathsf{Bad}_{\rm auth}]$ in the original native game.  The
attribution residual bounds its part on $E\wedge\neg A$.

\paragraph{Authentication-specific comparisons.}
For each fixed $j\in J$, use the native stopped continuation
$G_{\rm nat}^{[j]}$, its real-component record idealization $V^{[j]}$, and the
uniform-component continuation $\widetilde V^{[j]}$ from the candidate-local
accounting.  Apply to each of these worlds the same terminal Boolean code
\[
 W_j=\mathbf1[
 E\wedge A\wedge I=j\wedge\mathsf{Bad}_{\rm auth}],
 \qquad W_j=0\text{ after a declared local stop}.
\]
All eligibility, attribution, tuple selection, ledger membership, and event
flags are evaluated on the execution of that world.  They are computable from
the classical execution and provenance ledgers available to the existing
reduction interface; no test of an unknown challenge bit is used.

Let $a_j^{\rm auth}$ be the real-versus-uniform advantage of the operational
ML-KEM adapter of Section~\ref{sec:contextbound-connection} when its output
bit is $W_j$.  Its public key is installed at the original generation's
first actual private invocation, its encapsulation challenge at the original
encapsulation's first actual private invocation, and every post-challenge decapsulation of
the selected PQ ciphertext under the selected generation owner returns the
challenge secret, regardless of the traditional ciphertext or context.
The adapter preserves the native traditional-component validation and
rejections, the prechallenge history, and the exact-retry caches.
The selected-pair correctness allowance is
$\delta_j$.  Consequently this adapter gives
\begin{equation}
 \left|\Pr_{V^{[j]}}[W_j=1]
       -\Pr_{\widetilde V^{[j]}}[W_j=1]\right|
 \le a_j^{\rm auth}+\delta_j.
 \label{eq:native-auth-component}
\end{equation}
This uses the same two-stage component experiment as the existing adapter,
with the authentication event as its output bit.

Use the whole-continuation budgets of the ContextBound, root-initialization,
joint-schedule, and confirmation comparisons.  Define
\begin{align}
 \Omega_j^{\rm auth}
 :=\min\bigl\{1,{}&
 \Phi_{256}(d_j,Q_{H,j})+\Delta_{{\rm init},j}
       +\Delta_{{\rm sched},j}+\Delta_{{\rm conf},j}
       +V_j2^{-256}\nonumber\\[-1mm]
       &+\chi_H(Q_{H,j}^{\rm conf})
       +\varepsilon_{{\rm INT},j}\bigr\}.
 \label{eq:native-auth-continuation}
\end{align}
Here $d_j$ counts distinct effective protected tuples in the ContextBound
comparison; $\Delta_{{\rm init},j}$ is the event-bit cost of replacing the
selected contribution's single native root-initialization entry by a fresh
root; $\Delta_{{\rm sched},j}$ is the shared-compression joint-schedule cost;
$\Delta_{{\rm conf},j}$ replaces the three confirmation HMAC rows;
$V_j$ counts the native confirmation verification attempts;
$\chi_H$ is the framed-input collision bound; and
$\varepsilon_{{\rm INT},j}$ is the one-key, at-most-one-encryption
first-message AEAD integrity advantage with the native decryption interface.
Each query budget includes preparation, ordinary primitive computations,
honest protocol work, later continuation, and terminal event evaluation of
its derived experiment.  Budgets are fixed before interaction and are not
reset at a challenge or at a hybrid boundary.

\begin{theorem}[Native selected-root authentication]
\label{thm:native-selected-authentication}
In the live-core single-root next-rekey profile, with the admitted marked
controller just specified and the origin-preservation interfaces of
Section~\ref{sec:native-correspondence}, Lemma~\ref{lem:cb-native-root-use},
and Theorem~\ref{thm:joint-message-leaf},
\begin{equation}
 \Pr_{G_{\rm nat}}[E\wedge\mathsf{Bad}_{\rm auth}]
 \le
 \min\!\left\{1,\;
 \eta_{\rm attr}^{\rm nat}
 +\sum_{j\in J}
  \min\{1,\varepsilon_{{\rm rec},j}
           +a_j^{\rm auth}+\delta_j+\Omega_j^{\rm auth}\}
 \right\}.
 \label{eq:native-selected-authentication}
\end{equation}
The record cost $\varepsilon_{{\rm rec},j}$ is the live-core record event
cost~\eqref{eq:live-core-record-cost}, including its original fresh-seal
equality allowance.  The attribution
residual is measured once in the original native game.  For the executed-slot
selector, it may be instantiated by
Proposition~\ref{prop:executed-source-attribution} and the established
reservation-signature bridge.  An explicitly modeled projection extension adds
its original selected-candidate residual $\eta_{\rm proj}^{\rm nat}$ and uses
the corresponding event-bit record/report comparison in place of
$\varepsilon_{{\rm rec},j}$.
\end{theorem}

\begin{proof}
First localize the native event.  Write
$F=E\wedge\mathsf{Bad}_{\rm auth}$.  On attributed native executions the
events $I=j$ partition $F\wedge A$.  For a general declared projection
extension, split off $F\wedge A\wedge B_I$.  Its probability is at most
$\eta_{\rm proj}^{\rm nat}=\Pr[E\wedge A\wedge B_I]$; likewise
$\Pr[F\wedge\neg A]\le\eta_{\rm attr}^{\rm nat}$.
For each fixed $j$, couple the original execution and $G_{\rm nat}^{[j]}$
until the first free or charged local stop.  The free-stop certificate
permanently falsifies $E\wedge A\wedge I=j$ on every continuation.  A charged
stop is absent on the summand restricted to $\neg B_j$.  Therefore the same
identical-prefix argument as Lemma~\ref{lem:native-local-score}, now applied to a
Boolean variable, gives the exact identity
\[
 \Pr_{G_{\rm nat}}[F\wedge A\wedge I=j\wedge\neg B_j]
 =\Pr_{G_{\rm nat}^{[j]}}[W_j=1].
\]
Note that this identity remains valid when a recorded authentication event occurs
before the challenge send or terminal selection: the stopped predicate is the conjunction
with final eligibility and attribution, and its free-stop certificate
concerns every possible continuation.  Consequently
\begin{equation}
 \Pr_{G_{\rm nat}}[F]
 \le\eta_{\rm attr}^{\rm nat}+\eta_{\rm proj}^{\rm nat}
    +\sum_j\Pr_{G_{\rm nat}^{[j]}}[W_j=1].
 \label{eq:native-auth-localization}
\end{equation}
The live-core interface has $B_j=\varnothing$ and
$\eta_{\rm proj}^{\rm nat}=0$.

Next compare a fixed continuation.  The native representation invariant and
immutable-history simulator preserve the selected event's classical
execution, role-key provenance, private-tag and private-encryption ledgers,
and original allocation indices.  The real-to-opaque record comparison
therefore applies with output $W_j$.  Because it is a bit, its probability
changes by at most $\varepsilon_{{\rm rec},j}$.  This gives
\[
 \Pr_{G_{\rm nat}^{[j]}}[W_j=1]
 \le\varepsilon_{{\rm rec},j}
       +\Pr_{V^{[j]}}[W_j=1]
 \le\varepsilon_{{\rm rec},j}+a_j^{\rm auth}+\delta_j
       +\Pr_{\widetilde V^{[j]}}[W_j=1],
\]
where the second inequality is the operational component comparison
\eqref{eq:native-auth-component}.  All role-key and honest-origin ledgers
remain private reduction bookkeeping.  Membership is checked by recording
the actual complete core or encryption pair at its native invocation; it
requires no inspection of quantum queries or unknown component tape.

In $\widetilde V^{[j]}$, the protected PQ component is an independent uniform
256-bit value.  Apply Lemma~\ref{lem:cb-hidden-component} to the same Boolean
continuation, with the complete effective tuple cache.  It replaces each
protected ContextBound answer by a private random-function entry at cost
$\Phi_{256}(d_j,Q_{H,j})$.  Exact aliases, including aliases occurring before
the eventual selected send, share their entry.  The target entry is marked
at its first actual private use as prescribed by the existing controller
interface, so its fresh value is sampled without conditioning on a future
selection.  If this target is never allocated or never becomes eligible,
$W_j=0$.

Lemma~\ref{lem:cb-native-root-use} places every use of this target contribution
in one old-root and complete-info initialization entry.  Its event-bit
comparison~\eqref{eq:appendix-root-initialization} costs
$\Delta_{{\rm init},j}$ and leaves the uniform selected root controller of
Proposition~\ref{prop:confirmation-message-auth}.  All other tuple
entries, numerical coincidences, asynchronous confirmations, failed writes,
recoveries, and permitted successor computations are simulated as in that
controller.  The event $W_j=1$ implies an instance of its selected-root
confirmation or marked first-message integrity event.  The existing
uniform-root authentication bound hence gives
\[
 \Pr_{\widetilde V^{[j]}}[W_j=1]
 \le\Omega_j^{\rm auth}.
\]
In particular the shared-compression schedule, the confirmation rows,
fresh-tag guesses, framed-input collisions, and first-message integrity
contribute exactly the event costs displayed in
\eqref{eq:native-auth-continuation}.  Additional eligibility and source
requirements can only restrict this event in the final world.

Each fixed candidate is thus bounded by
$\min\{1,\varepsilon_{{\rm rec},j}+a_j^{\rm auth}+\delta_j
+\Omega_j^{\rm auth}\}$.  Substitute these bounds in
\eqref{eq:native-auth-localization} and cap the result by one.  The partition
was used only in the original native game.  Each ideal continuation runs its
own selector and records its own event flags, so no partition of different
hybrid worlds is required.
\end{proof}

\paragraph{Interpreting the authentication bound.}
The theorem establishes honest-origin authentication for the root
identified by the accepted send: it bounds the probability of a confirmation
or marked first-message installation with no matching honest origin.
Testing the same Boolean event at each stage makes each record or primitive comparison
enter with coefficient one. In particular, $a_j^{\rm auth}$ is the
operational component advantage for the distinguisher that tests this
event, with the component assumption evaluated at that distinguisher's
budget for the complete continuation.

\subsection{Native recovery with an online selected message}
\label{sec:native-online-chain}

The receiver may advance the chain and then reject the incoming
packet. If it computed a step for the first time during that attempt,
the observer had to decide whether to mark the step just beforehand.
Waiting for the sender would let it choose a step that had already
been computed. Earlier disclosed keys may guide the choice. We now
tie that choice to the accepted send by following its root, traffic
material and step from their first allocations, together with the
original reservations and the session's acquisition history.

\paragraph{The online native controller.}
Retain the established-session live core of
Section~\ref{sec:live-core-instantiation}, with one authoritative history
per endpoint and passive copies, and fix $1\le L\le2^{64}-1$ before
interaction. For each fixed candidate $j=(g,a)$, designate its original
effective tuple $m^*$ immediately before its first private ContextBound
hash computation at any protected alias, as in
Section~\ref{sec:contextbound-connection}. Designate the complete
traffic entry $I^*$ and direction before that entry's first block allocation.
These designations are retained through later acceptance, recovery and
equal-input aliases.

The controller records each computed step under $(\mathcal M,d,i)$.
Here $\mathcal M$ identifies the complete original traffic material,
$d$ the direction and $i$ the chain index. Equal-info aliases of one
root schedule share $\mathcal M$. Independently allocated material has
its own handle even if the epoch, index or key bytes coincide.

Immediately before a new step in the selected direction, the controller
uses the history so far and the step's public descriptor to choose
Continue or Mark. It can mark at most once, at an index $i<L$. An exact retry after a
failed write finds the entry for its original step. This table belongs
to the reduction. The endpoints follow their ordinary transitions and
return their ordinary responses.

For each step passed with Continue before Mark, the observer may obtain
its message key; the prefix seed and Extract value retain private
handles. At a marked
step $\ell$, the observer may obtain the paired successor seed
$s_{\ell+1}$ after its derivation. This named projection is bound to
$(\mathcal M,d,\ell)$, irrespective of subsequent progress at either
endpoint. Later chain steps are real computations from that seed.
The selected current root, private prefix seeds, Extract values and
marked message key have only the uses in the preceding interfaces.
The already admitted successor-root, opposite-direction and other
traffic-entry observations remain available.

These named projections are the only extra observations. Asking for
the successor seed returns that one value. Acquiring a session version
that contains the very same seed gives the observer the record too.
The acquisition keeps its full scope and the version's original
first-save time. Public commitments and failed captures keep their
ordinary behavior.

There is one left-or-right application operation, tied to the marked
step's original send. Its messages have equal length and common
associated data. A native rejection, uncertain outcome or exact retry
does not move that operation to another step. The terminal selector
chooses the cut and original reservation handles as before and must
identify this same accepted logical send. If no mark or accepted send
is reached, the original valid-target condition is false. Application
payloads and payload-bearing delivery remain private; the adversary
has the existing plaintext-free Resume and public protocol interfaces.
Thus $E$ is exactly the original eligibility predicate evaluated on
this controller's execution.

A mark may first occur during receive-side advancement, before the
target send is reached. Accordingly the one-key AEAD experiments below
allow classical decryption queries before their at-most-one encryption
call. Such earlier queries retain their actual replies. After the
left-or-right call, an exact challenge decryption returns the existing
authenticated marker. This operational convention preserves the
native receive-before-send order and is included in the stated AEAD
advantages.

\begin{lemma}[Native online-chain correspondence]
\label{lem:native-online-chain-correspondence}
For this controller and a fixed candidate $j$, the stopped native live
core realizes the prefix interface of Section~\ref{sec:chain-prefix}
after the existing record, component, ContextBound and root-entry
comparisons. Its original send and receive computations use one cached
step per $(\mathcal M,d,i)$. The marked step has at most one original
encryption tuple, and the live-core charged projection event remains
empty.
\end{lemma}
\begin{proof}
At traffic installation, the exact root schedule and complete traffic
information string determine the material handle. The role-dependent
assignment in the native traffic constructor makes one endpoint's
send seed the matching endpoint's receive seed; the opposite direction
has its other handle. Copying, reopening and exact pending-target
reconciliation preserve this association. The source-attribution
induction fixes the KEM origins of each installed schedule independently
of which message index is later marked.

Induct on actual private chain calls across both endpoints. For index
zero, the constructor supplies the corresponding traffic block. A
later call uses the predecessor's successor with the fixed native
domain and big-endian index. A repeated call under the same material,
direction and index therefore repeats the same Extract/Expand inputs.
Record its result at first computation, including in transaction-private
receive state. If receive authentication or a later durable update
fails, the native state may return to an earlier counter, while this
proof ledger retains the original result. A retried advance obtains
that result again. Taking an already retained skipped key likewise
introduces no new computation or marking opportunity. This is an
exact representation of the deterministic native computation in
each comparison world.

The sender stores its pending message plan before reaching the chain
step and encryption. A first-retention failure exits before encryption;
once encryption is reached, the authoritative plan binds the message
identifier, plaintext and associated data. A failed later write or
process loss recovers that same plan or its committed outbox. Exact
resumption uses its original encryption tuple, and successful commitment
advances the send index. Hence the first encryption answer can be
cached before its durable outcome, with at most one original tuple
under the marked material and index. This also covers a mark whose
first step computation occurred at the receiver.

The representation invariant for images and intents carries the same
material and step handles through their immutable versions. Acquiring
an intact version of an occurred candidate reservation's session,
saved at or after that reservation, certifies the existing free stop.
Lemma~\ref{lem:adaptive-origin-stop} applies before a final cut is chosen:
every eligible cut precedes both reservations and therefore that
acquired version's save. This certificate is independent of message
index or marking time. Complete old or unrelated versions
retain their ordinary correspondence. The new individual projections
return exactly earlier message blocks and the marked pair's successor,
which the prefix reduction already supplies. They expose no additional
protected record or report key. Thus the live core retains
$B_j=\varnothing$ and its established real-to-opaque record comparison.

After the ContextBound and root-entry comparisons, the complete native
continuation therefore has the three online designations and the
private prefix uses of Theorem~\ref{thm:online-chain-leaf}. Real control
verification, ACKs, other materials and later epochs remain in that
continuation. The whole-execution budgets include their computations
and all newly permitted prefix and suffix observations.
\end{proof}

\paragraph{Original-game costs.}
Use the original native $E,A,I$ and its residual
$\eta_{\rm attr}^{\rm nat}=\Pr[E\wedge\neg A]$. Let
$\varepsilon_{{\rm rec},j}$ be the record event cost in
\eqref{eq:live-core-record-cost}, now with the complete online controller's
budgets. The operational component adapters retain their definitions
from Section~\ref{sec:contextbound-connection}, with their output events
evaluated in this expanded experiment. For candidate $j$, set
\begin{equation}
 B_j^{\rm chain}=\min\bigl\{1,\;
  2\Phi_{256}(d_j,Q_{H,j})+2\Delta_{{\rm init},j}
  +2\Delta_{{\rm prefix},j}(L)+2\varepsilon_{{\rm AEAD},j}\bigr\}.
 \label{eq:native-chain-continuation}
\end{equation}
Every budget is fixed before interaction and covers preparation,
failed attempts, all endpoint aliases, the adaptive mark, real suffix
computation and the terminal test in its corresponding comparison.
The prefix sum in \eqref{eq:prefix-error} gives the explicit index-wise
comparison cost. No additional combinatorial message-index guessing
factor is introduced; all other whole-execution budgets may also grow
with $L$.

Let $\mathsf{Bad}_{\rm msg}^{\rm chain}$ record a fresh authenticated
receive and inbox installation under the marked step whose complete
$(\mathrm{AAD},\mathrm{ciphertext})$ pair is absent from the earlier
honest encryption ledger for that material, direction and index.
The flag is recorded at that installation, including before the target
send, and later encryption does not erase it. Accepted-send provenance
identifies the event at the end of the original experiment. For the
component adapter with Boolean output
$W_j=\mathbf1[E\wedge A\wedge I=j\wedge
\mathsf{Bad}_{\rm msg}^{\rm chain}]$, write its operational advantage
as $a_j^{\rm msg}$ and define
\begin{equation}
 \Omega_j^{\rm chain}=\min\bigl\{1,\;
  \Phi_{256}(d_j,Q_{H,j})+\Delta_{{\rm init},j}
  +\Delta_{{\rm prefix},j}(L)+\varepsilon_{{\rm INT},j}\bigr\}.
 \label{eq:native-chain-integrity-cost}
\end{equation}
The confidentiality test and the installation event each have their
own distinguisher. We use the complete query budget of the
distinguisher for the test being evaluated.

\begin{theorem}[Native online selected-message security]
\label{thm:native-online-message}
In the online native live-core experiment above,
\begin{equation}
\begin{split}
 \operatorname{Adv}_{\rm msg}^{\rm chain}
 \le\min\biggl\{\frac12,\frac12\Bigl[
    \eta_{\rm attr}^{\rm nat}\\[-2pt]
 \qquad+
 \sum_{j\in J}\bigl(2\varepsilon_{{\rm rec},j}
       +a_j^++a_j^-+2\delta_j+B_j^{\rm chain}\bigr)
                       \Bigr]\biggr\},
\end{split}
\label{eq:native-online-confidentiality}
\end{equation}
and
\begin{equation}
\begin{split}
 \Pr_{G_{\rm nat}}[E\wedge\mathsf{Bad}_{\rm msg}^{\rm chain}]
 \le\min\biggl\{1,\;&\eta_{\rm attr}^{\rm nat}\\[-2pt]
 &+\sum_{j\in J}\min\{1,\varepsilon_{{\rm rec},j}
       +a_j^{\rm msg}+\delta_j+\Omega_j^{\rm chain}\}\biggr\}.
\end{split}
\label{eq:native-online-integrity}
\end{equation}
For the executed-slot selector, the corresponding original-game
residual in each bound is at most $\Delta_{\rm attr}$ from
\eqref{eq:executed-attribution-cost}, using the complete expanded
signing experiment and its established reservation-signature bridge.
\end{theorem}
\begin{proof}
For confidentiality, put
$S=2\mathbf1[g=b]-1$ and run the same available candidate terminal
code $\Psi_j=\mathbf1[E\wedge A\wedge I=j]S$ in every world.
The identity in Lemma~\ref{lem:native-local-score} localizes the
original score before any candidate-specific comparison. The native
free stops remain valid by
Lemma~\ref{lem:native-online-chain-correspondence}. Its record
correspondence gives $\rho_j\le2\varepsilon_{{\rm rec},j}$;
the two operational component-event comparisons cost
$a_j^++a_j^-+2\delta_j$.

In the random-component world, apply
Lemma~\ref{lem:cb-hidden-component} with the unchanged complete tuple
cache. The target contribution is selected at its own first private
allocation, before the later traffic and message designations. Its
single root-use entry is still that of
Lemma~\ref{lem:cb-native-root-use}, since later chain work does not
reuse the combined contribution or alter its original pending-root
inputs. The two score comparisons therefore cost
$2\Phi_{256}(d_j,Q_{H,j})+2\Delta_{{\rm init},j}$.
The resulting uniform-root continuation is exactly the online prefix
interface by the new correspondence lemma.
Corollary~\ref{cor:online-chain-message} bounds its qualified signed
score by $2\Delta_{{\rm prefix},j}(L)+2\varepsilon_{{\rm AEAD},j}$.
Every selector, disclosure, eligibility check and terminal event is
recomputed within that world. In particular, qualifying a run uses
the available provenance and view rather than direct access to $b$.
Thus the random-component score is bounded by $B_j^{\rm chain}$.
Substitution in Proposition~\ref{prop:native-local-accounting}
proves~\eqref{eq:native-online-confidentiality}.

For integrity, apply the Boolean localization identity from
\eqref{eq:native-auth-localization} to the selected-message event
$W_j$. The same free-stop certificate makes a stopped candidate's
final conjunction false, even if its recorded inbox event preceded
the challenge send. The record comparison costs
$\varepsilon_{{\rm rec},j}$, and the operational component comparison
costs $a_j^{\rm msg}+\delta_j$. Apply the complete-tuple and root-entry
comparisons to this Boolean continuation. The remaining event is a
subset of the selected-key integrity event in
Corollary~\ref{cor:online-chain-message}, so its probability is bounded
by $\Omega_j^{\rm chain}$. A fresh successful decryption before the
honest encryption is already a valid integrity witness; the reduction
may return it immediately. This justifies the same bound when Mark
was first reached by the receiver. Summing the native candidate
partition and capping by one proves~\eqref{eq:native-online-integrity}.

Finally, the executed-slot association proof is an induction over
traffic installation and its originating KEM plans. Its signed-body
and full-wire checks precede the message-index choice. Advancing,
retaining or reopening the installed traffic entry preserves those
origins, while the earlier-key projections only enlarge the complete
adversarial continuation. Proposition~\ref{prop:executed-source-attribution}
therefore gives the stated original-game attribution bound with
its updated whole-execution resources. This term is charged once,
and the root, record and component costs retain their displayed
coefficients; no additional guess of the marked message index is used.
\end{proof}

\paragraph{Reading the bound in ordinary primitives.}
Table~\ref{tab:native-online-substitutions} expresses
\eqref{eq:native-online-confidentiality}--\eqref{eq:native-online-integrity}
in ordinary primitives for the executed-slot selector. Read each
$\sigma\in\{+,-,\mathrm{msg}\}$ term with the complete resources of
that terminal test. The source pair and the send it identifies are the
same throughout the substitution.

\begin{table}[htbp]
\caption{Primitive substitutions for native online selected-message security.
The attribution cost is charged once outside the candidate sum.}
\label{tab:native-online-substitutions}
\small
\begin{tabularx}{\textwidth}{@{}L{0.19\textwidth}Y@{}}
\toprule
Term & Established substitution and its role \\
\midrule
Record cost $\varepsilon_{{\rm rec},j}$ &
$\displaystyle\min\!\left\{1,
\sum_{i=1}^{M_j^{\rm seal}}\Phi_{192}(d_{j,i},Q_{j,i})
+\delta_{{\rm seal},j}\right\}$.
\par Here $\delta_{{\rm seal},j}=\min\{1,\binom{M_j^{\rm seal}}2 2^{-192}\}$.
Equation~\eqref{eq:live-core-record-cost} counts fresh sealing attempts,
including failed commits; copies retain their original handles. \\
\addlinespace[3pt]
Component cost $a_j^\sigma$ &
At most $\min\{1,\operatorname{Adv}^{\rm cca}(B_j^\sigma)
+\zeta_{{\rm pre},j}^{\sigma}\}$.
\par Here
$\zeta_{{\rm pre},j}^{\sigma}
=\min\{1,q_{j,\sigma}^{-}2^{-256}+\delta_{{\rm mr},j}^{\sigma}\}$
by Proposition~\ref{prop:kem-deferred-challenge}.
The bound retains real prechallenge decapsulation. \\
\addlinespace[3pt]
Correctness $\delta_j$ &
The selected pair's shared-key correctness cost from
Section~\ref{sec:contextbound-connection}. It remains separate from the
message-recovery event $\delta_{{\rm mr},j}^{\sigma}$ inside the component
bridge. \\
\addlinespace[3pt]
Continuation $B_j^{\rm chain}$,
$\Omega_j^{\rm chain}$ &
Equations~\eqref{eq:native-chain-continuation} and
\eqref{eq:native-chain-integrity-cost} combine ContextBound replacement,
root initialization, the finite prefix sum in~\eqref{eq:prefix-error},
and one-key AEAD confidentiality or integrity. Each comparison costs twice
its event bound for the signed score and once for the Boolean event. \\
\addlinespace[3pt]
Attribution $\Delta_{\rm attr}$ &
$\min\{1,\varepsilon_{\rm sig}+\chi_H(Q_H)\}$ by
Proposition~\ref{prop:executed-source-attribution}, using the ordinary
ML-DSA instantiation below. This is the original-game association cost
for the executed-slot selector. \\
\bottomrule
\end{tabularx}
\end{table}

Specifically, for at most $U$ fresh signing-owner families, the established
ordinary-signature bridge~\eqref{eq:native-pair-ordinary-mldsa} permits
\[
\begin{split}
 \varepsilon_{\rm sig}\le\min\bigl\{1,U\bigl(&
 \varepsilon_{\rm MLDSA,EUF}^{\rm su}+\beta_{256}(q)\\[-2pt]
 &+\delta_{\mathsf X}(q_*^{\rm res})
  +\delta_{\mathsf X}(q_*^{\rm ord})\bigr)\bigr\}.
\end{split}
\]
Here $q$ bounds distinct executed target $(r,\mu)$ pairs,
$\beta_{256}(q)=\min\{1,\binom q2 2^{-256}\}$, and the two $q_*$ budgets
count remaining shared-XOF calls in the complete reservation and ordinary
experiments. The established oracle terms are
$\delta_{\mathsf C}(v)=\min\{1,v2^{-256}\}$ in the classical model and
$\delta_{\mathsf Q}(v)=\min\{1,2v2^{-128}\}$ in the quantum model.

Confidentiality uses the two component tests $+$ and $-$, each with its
own prechallenge term, before the outer factor $1/2$. Integrity uses the
single $\mathrm{msg}$ test. The existing $2\delta_j$ and $\delta_j$
coefficients therefore remain in their respective bounds. The live core
has zero charged projection residual. Online marking incurs the displayed
$L$-term prefix sum; its complete-query budgets include preparation,
failed attempts, endpoint aliases, real suffix computation and terminal
testing. These are the same interfaces and budgets as in the theorem.

Setting $L=1$ recovers the first-message schedule cost. For a longer
chain the observer can see earlier keys and mark a later step while
it is still uncomputed. That choice adds the prefix cost above. The
reservations and acquisitions are still those of the accepted send.
If the interface also permits partial record or report projections,
its $\eta_{\rm proj}^{\rm nat}$ and local comparison costs enter as
in the preceding native theorems.


% Finite multiple-root entry. Uses the existing joint-primitive bounds.
\subsection{Several root entries from one recovered contribution}
\label{sec:multiple-root-entry}

The same recovered contribution may be used with different old roots
or different root-information strings. An observer who learns one
result may use it to select the inputs for another. We compare the
whole family so that a repeated request still gets its original root.
Let $S$ be uniform on 256 bits, initially independent of the public compression
oracle $h$ and auxiliary state, and retain $F$ and $\Gamma$ from
Section~\ref{sec:joint-continuation}. Each root request supplies a 32-byte
salt $r$ and a complete native 97-byte information string $I$. Its real value is
\[
 P_r=F(r,S),\qquad R_{r,I}=F(P_r,I\Vert\mathtt{01}).
\]
The parent contribution is used through these Extract calls; each resulting
PRK is used through its root-Expand interface. The controller retains the
original source ledger and any established local stops.

\paragraph{Descriptors and available data.}
The root-use descriptor is
$e(r,I)=(\Gamma(r),I)$, with exact byte equality. Equal descriptors retrieve
the same root entry across callers and resumptions. Indeed,
\eqref{eq:appendix-column-compression} shows that equal $\Gamma$ values give
equal Extract values for every $S$, and equal complete information strings
then give equal Expand values. Assign each newly encountered $\Gamma$ class
its next lifetime index $i$, preserving this index after a restart.
Numerical output coincidences between distinct classes keep their ordinary
shared-oracle distribution.

The controller computes salts, information strings and its selection from
its available state: executed preparation, independent coins, public-oracle
responses, unmarked roots, and the declared outputs of the marked schedule.
At most one fresh descriptor is marked before its first root evaluation.
Its root is installed in one retained schedule handle; every descriptor alias
retrieves that handle. All its uses belong to the same interface of
Section~\ref{sec:joint-continuation}, including its session salt, confirmation
and successor calls, complete traffic-information table and budgets. Other
root entries may be disclosed and used in later available computations,
including as salts. The retained handle supplies exactly that schedule's
operations. This is a rule about available computation; a locally known salt
can numerically coincide with a protected value.

Fix finite bounds $D,C,N_1,\ldots,N_C$ before execution: at most $D$ distinct
salt strings, at most $C\le D$ classes, and at most $N_i$ distinct complete
information strings in the $i$th class. These bounds hold in every comparison
experiment. All oracle-dependent preparation is executed and counted. Exact
repeats are cached. Write $q_S$ for the complete ordinary compression-query
budget in the column reduction, and $q_i$ for that in the $i$th row reduction.
These budgets include the whole controller, all other real primitives,
continuation and terminal test. An additional independent whole-hash oracle
may remain available throughout.

\begin{theorem}[Finite multiple-root entry]
\label{thm:multiple-root-entry}
Replace the real root interface by a private table
$\mathcal W(\Gamma(r),I)$ of independent uniform 32-byte entries, preserving
all descriptor aliases and the controller above. Every efficient available
output-bit test changes by at most
\begin{equation}
 \Delta_{\rm mr}=\min\!\left\{1,
 \varepsilon_{\rm Hjoint}(C,0,q_S,32)
 +\sum_{i=1}^{C}\varepsilon_{\rm Hrow}(N_i,q_i,98)\right\}.
 \label{eq:multiple-root-error}
\end{equation}
In the table experiment, an online marked fresh descriptor supplies the
independent uniform root required by
Theorem~\ref{thm:joint-message-leaf}, together with the complete interactive
controller. The marked descriptor may belong to a previously used class.
\end{theorem}
\begin{proof}
Compute $\Gamma(r)$ from available salt bytes and retain its class index.
Use one column call for the first salt in each new class; reuse its reply for
later equal classes. The descriptor identity makes this an exact rendering
of the real Extract interface. Lemma~\ref{lem:appendix-joint-hmac} replaces
these calls by $P_i=V(\Gamma(r))$ at the first cost in
\eqref{eq:multiple-root-error}. The remaining root computation consists of
HMAC rows under the independent private keys $P_i$.

Replace these rows in first-class-allocation order. At step $i$, presample
the hidden uniform row-challenge key and assign it to the $i$th new class
when that class is first encountered. Before that event the key has no use.
Classes already replaced use their private tables; later classes have locally
sampled independent keys and are evaluated with the same $h$. The controller
accesses the selected PRK through its row alone. Thus the real row branch
matches the preceding hybrid and the ideal row branch matches the next one.
Lemma~\ref{lem:appendix-hmac-instances} bounds this step by
$\varepsilon_{\rm Hrow}(N_i,q_i,98)$. If the class is never allocated, the
hybrids agree. Summing over the fixed class budget establishes the bound.

Distinct classes remain distinct reduction slots even when their sampled
PRK bytes coincide. In a real row branch, local evaluation under another
class's key uses the same $h$ and automatically preserves any resulting
coincidence; these ordinary queries are included in $q_i$. Thus the argument
uses the full independent-key distribution.

For each ideal row, $I\mapsto I\Vert\mathtt{01}$ is injective. Its root
interface is therefore exactly the table $\mathcal W(\Gamma(r),I)$.
At a fresh descriptor, its value is independent of all previously used table
entries and the oracle. Presample the designated uniform root, withhold it
during the executed prefix, and assign it when the controller marks that
fresh descriptor. All other entries retain their independent sampling and
all aliases retrieve the retained handle. The selection uses the available
prefix, so this deferred assignment has exactly the table distribution.
It embeds the whole controller in the existing root experiment. The row
hybrids already covered every class, so target selection adds no class guess.
\end{proof}

\paragraph{Explicit query accounting.}
Let $B$ bound all ordinary compression work outside this entry service,
including its entire downstream schedule and test. Cache the two public
$\Gamma$ computations per distinct salt. With
$N=\sum_{k=1}^{C}N_k$, valid conservative choices are
\begin{equation}
 q_S=B+2D+5N,\qquad
 q_i=B+2D+5\sum_{k=i+1}^{C}N_k.
 \label{eq:multiple-root-budgets}
\end{equation}
A 98-byte HMAC evaluation uses five compression calls, including its two key
pads. In the column comparison every root row is real; at row step $i$, the
later rows are real and the earlier rows are tables. External class scheduling
accounts for $2D$, including new salts that share an existing class. The
column primitive's own public pad work remains included in its established
bound. $B$ is a uniform bound over these complete derived experiments.

\begin{corollary}[First message after a marked root entry]
\label{cor:multiple-root-message}
For the available controller above and its marked schedule, use the
first-message interface and available qualification predicate of Corollary~\ref{cor:joint-first-message}. If $Z$
is its qualified signed message score, with zero on runs having no marked
challenge, then
\begin{equation}
 |\mathbb E Z|\le
 \min\{1,\,2\Delta_{\rm mr}+2\Delta_{\rm sched}
                      +2\varepsilon_{\rm AEAD}\}.
 \label{eq:multiple-root-message}
\end{equation}
The schedule and AEAD reductions include the complete entry controller and
its shared-oracle work. If the contribution itself comes from the protected
ContextBound table of Lemma~\ref{lem:cb-hidden-component}, its additional
signed-score cost is $2\Phi_{256}(d,Q_H)$. For a candidate whose stopped execution realizes this entry controller,
it replaces the single-entry $2\Delta_{\rm init}$ term by
$2\Delta_{\rm mr}$ in the corresponding continuation bound; the original accounting terms keep their
established definitions.
\end{corollary}
\begin{proof}
Fix each application challenge bit for the entry comparison and average
afterwards, as in the retained continuation. The randomized output test
with success probability $(Z+1)/2$ converts
Theorem~\ref{thm:multiple-root-entry}'s bit gap into a score gap of at most
$2\Delta_{\rm mr}$. Its ideal marked root initializes the entire controller
of Corollary~\ref{cor:joint-first-message}, which supplies the remaining two
terms. Apply the same conversion to the ContextBound comparison when present.
\end{proof}

% Source-backed owned SDK purpose-key table; prior root modules are retained.
\subsection{An owned SDK purpose-key table}
\label{sec:sdk-purpose-connection}

Repeated SDK derivation raises two questions: which calls share key
bytes, and how each owner's lifetime is managed. Calls with the
same source and complete derivation inputs share the key bytes. Each
successful call creates its own \texttt{DerivedKey} owner, which can be
closed separately. The native
\texttt{SharedSecret::derive\_key} operation takes a purpose, protocol
label and context and allows repeated calls on the source owner. The
source checks cover all five purposes at both peers, including output
agreement and owner lifetimes. We apply the finite entry-table
argument to the SDK's actual information encoding.

\begin{lemma}[Information-length form of the entry comparison]
\label{lem:entry-information-length}
For the finite service $R_{r,I}=F(F(r,S),I\Vert\mathtt{01})$, let the
complete available information strings in class $i$ have length at most
$A_i$ bytes. Put $L_i=A_i+1$ and $\lambda_i=\ell(L_i)$, within the stated
SHA-256 padding domain. Use the descriptor classes and finite budgets of
Theorem~\ref{thm:multiple-root-entry}. With the parent contribution used
through the Extract interface and each PRK through its Expand row, the
whole interactive service is within output-bit gap
\begin{equation}
 \Delta_{\rm ent}=\min\!\left\{1,
 \varepsilon_{\rm Hjoint}(C,0,q_S,32)
 +\sum_{i=1}^{C}\varepsilon_{\rm Hrow}(N_i,q_i,L_i)\right\}
 \label{eq:entry-information-length}
\end{equation}
of the consistent independent table indexed by $(\Gamma(r),I)$.
The continuation may process each returned value through its declared
interface. Conservative complete ordinary-query budgets are
\[
 q_S=B+2D+\sum_i(2+\lambda_i)N_i,\qquad
 q_i=B+2D+\sum_{k>i}(2+\lambda_k)N_k.
\]
\end{lemma}
\begin{proof}
Use the column replacement and first-class-allocation row hybrids of
Theorem~\ref{thm:multiple-root-entry}. At row $i$, every queried message
has length at most $L_i$. The map $I\mapsto I\Vert\mathtt{01}$ is injective
also for variable-length strings. Hence the final private row gives exactly
the full-information table. Each actual Expand costs at most
$2+\ell(L_i)$ compression calls, yielding the displayed budgets.
The reductions receive the root-row replies and run the entire declared
continuation from them. This proves the table comparison for that
continuation. The native 97-byte case has $L_i=98$ and $\lambda_i=3$.
\end{proof}

\paragraph{The actual SDK descriptor.}
Let $\sigma=\texttt{QPeriapt-SDK-HKDF-SHA256-v1}$, a 27-byte string, and
$r_0=\sigma\Vert\mathtt{00}^5$, its 32-byte zero extension. HMAC pads both $\sigma$ and $r_0$ to the same
64-byte key block~\cite{rfc2104}, so $F(\sigma,S)=F(r_0,S)$ for every $S$ and every $h$.
The complete SDK information string is
\begin{equation}
\begin{split}
 I={}&D_{\rm sdk}\Vert\operatorname{BE32}(1)
 \Vert\operatorname{BE32}(2)\Vert\mathit{policyState}_{36}
 \Vert\mathit{rootDigest}_{32}\\
 &\Vert\operatorname{BE32}(p)
 \Vert\operatorname{BE16}(|a|)\Vert a
 \Vert\operatorname{BE32}(|c|)\Vert c
 \Vert\operatorname{BE16}(32),
\end{split}
 \label{eq:sdk-purpose-info}
\end{equation}
where $D_{\rm sdk}=\texttt{QPeriapt-SDK-Key-v1}$ has 19 bytes. The fixed
suite code $1$ denotes ML-KEM-768+X25519 and profile code $2$ denotes
ContextBound. Here $p$ is one
of the five purpose codes, $a$ is the protocol label, and $c$ is the context.
The policy state is its 4-byte version followed by its exact-document
32-byte digest; the root digest is SHA-256 of the runtime's pinned trust
root. Labels contain $1$--$255$ printable non-space ASCII bytes and contexts
have at most $65536$ bytes. Thus
\[
 |I|=107+|a|+|c|\in[108,65898],\qquad |I\Vert\mathtt{01}|\le65899.
\]
Length parsing recovers every effective field. The descriptor is
$(\Gamma(r_0),I)$ across all owners of the same protected contribution.
Runtime identities and owner identifiers serve their separate lifecycle
roles. The effective digest bytes determine their respective info fields.

\paragraph{Owner correspondence.}
Fix a protected contribution family with common value $S$ and preserve the
actual metadata of each source and derived owner. Represent a derived owner
by $(\mathit{runtime},\mathit{live},I)$ and a reference to the common key
entry $K_I$. Every accepted derivation creates an owner, and equal $I$ values refer
to the same key entry. Closing a derived owner changes its own live
flag. Closing a source owner prevents subsequent derivation through
that owner; however, previously returned owners keep their keys. Once a
runtime closes, later open checks fail. A derivation or export that
already passed admission follows its actual completion path. Recipients
keep any copies they have exported.

These transitions follow directly from the source's separate
\texttt{Option} fields, shared runtime state and admission checks. The inspected
FFI derivation path also records reservation, publication and parent-revocation
metadata. A first private KDF computation allocates the descriptor
entry even if its later handle publication fails; a retry reuses that key
entry while obtaining its new owner. The Rust in-flight bound and the
FFI live-handle bound are preserved as metadata. The proof's lifetime bound
$N$ counts distinct complete infos.

\begin{proposition}[Native SDK key-table correspondence]
\label{prop:sdk-purpose-table}
For one initially independent uniform protected contribution, let the
family's source owners use the Rust derivation interface above, together with
derived-owner export, close and runtime lifecycle operations. The parent
contribution and internal Extract value remain within their stated KDF
interfaces. Let $N$ bound distinct infos before execution and $L$ bound their
Expand-message length. Then the entire owned-API experiment differs from
one using independent uniform keys $K_I$, with all aliases and owner states
preserved, by output-bit gap at most
\begin{equation}
\begin{split}
 \Delta_{\rm sdk}=\min\{1,\;&
 \varepsilon_{\rm Hjoint}(1,0,B+2+(2+\ell(L))N,32)\\
 &+\varepsilon_{\rm Hrow}(N,B+2,L)\}.
\end{split}
 \label{eq:sdk-purpose-error}
\end{equation}
The budgets cover complete preparation, other real computations and the
terminal test. At full SDK input limits, $L=65899$ and $\ell(L)=1031$.
\end{proposition}
\begin{proof}
Normalize the fixed salt as above. The first-block HKDF definition~\cite{rfc5869}
gives $P=F(r_0,S)$ followed by $F(P,I\Vert\mathtt{01})$; the source keeps
its PRK internal.
An induction over admitted calls preserves each owner state and every
public status: validation and admission use available metadata; derivation
creates a fresh owner referring to its descriptor's deterministic value;
exports apply that owner's current guards; close and publication operate
on their existing owner identities. Repeated key values can therefore be
memoized while all native owner transitions continue to execute.
Apply Lemma~\ref{lem:entry-information-length} with $D=C=1$.
The class/row comparison changes only these internal key values. All
field-dependent choices, exports and lifecycle responses are recomputed
in their own experiment by the same controller.
\end{proof}

For the existing five-purpose source workload, the label has 13 bytes and
its context 24, giving $|I|=144$, $L=145$ and $\ell(L)=4$. Its two peer
owners use five complete descriptors when their effective policy bindings
agree. Thus $N=5$ gives column ordinary budget $B+32$ and row budget $B+2$.
The table comparison permits disclosure of all derived entries and
establishes their joint pseudorandom separation. In particular, the
probability of a key collision among exports from distinct allocated infos,
with at most $N$ total infos, is at most $\min\{1,\binom N2 2^{-256}+\Delta_{\rm sdk}\}$.

\paragraph{A selected key for one application encryption.}
An application may mark a new complete descriptor before its first private
KDF evaluation, keep its key private, and use it in the one-key,
one-encryption AEAD experiment defined in
Section~\ref{sec:joint-continuation}, with that experiment's nonce,
equal-length messages, associated data and decryption restriction. Every source owner and derived alias of that descriptor routes
to the same key-use interface, with its own native liveness checks. Other
derived entries may be exported. Let $Z$ be the qualified signed score,
with qualification computed from the available classical view and ledger
and value zero when the marked challenge is absent. Fresh table allocation
then supplies the uniform AEAD key, giving
\[
 |\mathbb E Z|\le\min\{1,\,2\Delta_{\rm sdk}
                                  +2\varepsilon_{\rm AEAD}\}.
\]
Fix each application challenge bit for the KDF comparison and average the
two score bounds. In the final AEAD game, use the adversary's guess on
qualified runs and an independent fair guess otherwise. A preceding
complete-tuple ContextBound comparison contributes its established
$2\Phi_{256}(d,Q_H)$ term when the stopped execution realizes that source
family and available controller.

The SDK table thus instantiates the finite entry proof for owned purpose
keys and their declared application use. The durable rekey driver is the
corresponding instance with its own 97-byte pending-root information and
established source invariant.


\subsection{Bootstrap authentication from an original surviving source}
\label{sec:bootstrap-authentication}

When the initiator accepts a bootstrap connection, what can it infer
about the responder? The order in which the peers use the two secrets
matters. The initiator first encapsulates to the responder's prekeys
and uses the resulting $S_0$ for Initial confirmation. Only then does
the responder encapsulate to the initiator's newly reserved reply key
to obtain $S_1$. The joint schedule holds a pending root and activates
message state after Final. We follow that order back from connection
acceptance to the expected responder's earlier activation of the same
signed transcript.

\paragraph{Transcript and schedule.}
Let $c=(\mathit{ctx},x,y)$ contain the admitted context and the complete
canonical signed Initial and Reply wires, including their signature
bytes. With the source's framed SHA3 oracle $H$, define
\[
 \begin{split}
 f_c&=\mathtt{QPBSF001}\Vert\mathit{ctx}
       \Vert H_{\rm initial}(x)\Vert H_{\rm reply}(y),\\
 s_c&=H_{\rm session}(f_c).
 \end{split}
\]
The two complete-wire hashes include their distinct source domains.
Write $F$ for HMAC under the shared SHA-256 compression oracle $h$.
The bootstrap KDF and message domains $D_B,D_M$ have 48 and 43 bytes.
Let $y_{\rm core}$ contain the Reply prefix and ciphertext, ending before
its confirmation tag and outer signature, and put
$t_1=H_{\rm reply\text{-}core}(y_{\rm core})\in\{0,1\}^{256}$ using the
source's distinct framed SHA3 domain. Set
$I_s=D_B\Vert\texttt{handshake-seed}$,
$I_R=D_B\Vert\texttt{confirmation-responder}$ and
$I_I=D_B\Vert\texttt{confirmation-initiator}$.
The schedule is
\begin{equation}
 \begin{aligned}
 P_B&=F(t_1,S_0\Vert S_1),\\
 (\mathit{seed},K_R,K_I)&=
    (F(P_B,I_s\Vert\mathtt{01}),F(P_B,I_R\Vert\mathtt{01}),
       F(P_B,I_I\Vert\mathtt{01})),\\
 I_B&=D_B\Vert\texttt{session-root/}\Vert\mathit{ctx},\\
 K_B&=F(F(s_c,\mathit{seed}),I_B\Vert\mathtt{01}),\\
 I_A&=D_M\Vert\texttt{initial/HKDF-SHA256/ChaCha20Poly1305}
       \Vert\mathit{ctx},\\
 (R_0,C_I,C_R,A_I,A_R)&=
       \operatorname{Expand}(F(s_c,K_B),I_A,160).
 \end{aligned}
 \label{eq:bootstrap-schedule}
\end{equation}
The three bootstrap row inputs have lengths 63, 71 and 71 bytes.
Here $|I_B|=93$, $|I_A|=111$, and the activation row inputs have
lengths $112,144,144,144,144$. The responder confirms its activation
with the direction-1 epoch-zero ACK under $A_I$; the initiator's ordinary
ACK uses $A_R$ and direction 2. The authenticated epoch-zero string is
\[
 D_M\Vert\texttt{acknowledgement}\Vert\mathtt{QPCMACK1}
       \Vert s_c\Vert\mathtt{01}\Vert\operatorname{BE64}(f),
\]
where $f$ is the admitted consumption floor. This string has 107 bytes;
the transmitted prefix and 32-byte tag have 81 bytes in total.

The responder compares the complete 104-byte Final prefix to its own
$f_c$ before MAC verification and durable completion. Hence an accepted
honest Final equates $\mathit{ctx},x,y$ in the absence of collisions
on these framed hash inputs. A valid alternative signature on an issued
Reply body preserves its upstream $t_1,P_B$, while its actual complete
wire remains part of $s_c$ and the root descriptor.

\begin{definition}[Online bootstrap experiment]
\label{def:bootstrap-online}
The admitted operations are context admission, original bootstrap
reservation and resumption, Final, both message activations, epoch-zero
ACK generation and verification, checkpoint restoration, and the
specified source/key and record disclosures. The subsequent application
interface consists of these ACK operations and the specified observations
of the four sibling values. The controller admits concurrency, active
delivery, exact retries and passive copies along one authoritative
durable history for each original operation. Immediately before a class's
first private ACK computation or use at either endpoint, a fixed online rule
may register that complete class, provided its Reply has passed the
initiator's signature, context and confirmation checks or has been formed
by the authoritative responder. It rejects registration when the
cumulative source/guard closure already reconstructs $A_I$. Repeated
classes retain their rank; an unprotected used class stays unprotected.
A later request that would return an ACK recipe for any registered class,
or replace its trusted private state, stops before that return. An already
observed terminal event remains latched. Other observations and classes
follow their actual algorithms.
\end{definition}

The closure includes direct ACK material, $K_B$, the seed, both
contributions, and records and wrapping keys that together recover them.
Acquired versions survive local deletion in this history. Define
$\mathsf{Coll}$ on the relevant context, signed-wire and session-ID
hash inputs. The schedule comparison stops at this event and replaces
only the exported $A_I$ of each registered class by an independent cached
32-byte key. Both endpoints and exact retries retrieve that key.
The actual internal block 4 still feeds Expand's block 5 before this
export replacement; all sibling outputs, confirmations and subsequent
verification remain those of the current experiment.

\paragraph{A concrete surviving-source profile.}
We trace the ACK key back to the original KEM pair that supplies the
surviving entropy. In this profile, every protected class has an
honest authoritative Initial. The first contribution $S_0$ is publicly available
at creation, and the traditional KEM scalars are available at their
original sampling events. The initiator's signing identity may also be
disclosed. The expected responder's signing owner is protected, with its
BootstrapResponder purpose used through the authoritative response
history. The surviving source is the reply PQ generation and response
encapsulation pair. Its private owner provides the actual bootstrap and
checkpoint-restoration interfaces. Fix lifetime bounds $G,A,R$ for
reply-key reservations, response-encapsulation reservations and protected
registrations, counting draws whose subsequent save fails.

\begin{lemma}[Original reply pair and effective family]
\label{lem:bootstrap-reply-origin}
On the source-correctness event and before a fresh-body responder-signature
event, a relevant framed-hash collision, an honest Initial-nonce collision
or the encapsulation event
of Lemma~\ref{lem:bootstrap-validation}, a protected class identifies a
fixed original pair $j=(g,a)\in[G]\times[A]$. At $a$'s first actual
computation that pair fixes its complete ContextBound tuple, $S_1$,
Reply core and bootstrap PRK family. Same-body signature aliases retain
that family and their actual downstream full-wire descriptors.
\end{lemma}
\begin{proof}
The initiator persists its reply-key reservation before its first
generation call. Its 96-byte entropy consists of the 64-byte ML-KEM
$d\Vert z$ and a 32-byte traditional scalar. The responder persists
its response-encapsulation reservation before the first encapsulation;
its 64-byte entropy contains the 32-byte ML-KEM input $m$ and its
traditional scalar. The authenticated token binds the kind, policy,
operation and, for encapsulation, complete peer/context input.
Restoration consumes that same tape. Failed later writes, pending-target
reconciliation and exact retries reuse the original executed slot.

Monitor a fresh accepted purpose-6 body immediately after its signature
verification and before private decapsulation. Otherwise its honest
signing ledger and checked complete-Initial digest identify the response
history. Distinct honest owners publishing an equal complete Initial
require equal independently sampled 32-byte Initial nonces. Charging
that event identifies $g$, and the authoritative saved response
identifies $a$. Its fixed tape, peer key, context and complete Initial
determine both ciphertexts and component results and all ten fields of
the ContextBound tuple. Source correctness supplies the same $S_0$ at
the two endpoints. Thus $P_B=F(t_1,S_0\Vert S_1)$ is fixed before
response signing. The framed initial/reply context separation and
Lemma~\ref{lem:bootstrap-validation}'s monitor preserve that selected
challenge association at every admitted alias.
\end{proof}

For each of the adjacent hybrids below, write $\alpha_{r,b}$ for this
association cost. A valid choice is the sum, clipped at one, of the
protected responder's purpose/body signature bound, one complete-query
framed-hash collision bound, and $\binom G2 2^{-256}$. The ordinary
signing-reservation bridge of
Theorem~\ref{thm:native-mldsa-reservation-bridge} instantiates the
appropriate surviving signer. Its budget includes the complete
verification and continuation. The source pair requires $GA$ slots;
the original response already identifies the relevant PRK request.

\paragraph{Opening and original component embedding.}
Let $H_k$ replace the first $k$ registered ACKs and set
$H_{r,0}=H_{r-1}$, $H_{r,1}=H_r$. In either hybrid let $J_r$ be
the pair of the $r$th registration, or $\bot$ if it is absent. For
an available terminal bit $D$, the $\bot$ terms cancel: the adjacent
machines are identical until that registration. After the association
stop, their remaining difference is the sum of the expectations of
$D\mathbf1[J_r=j]$ over fixed pairs $j$.

Opening the selected source, $S_1$, PRK or seed before registration
permanently excludes every future protected class of that family: the
known $S_0$ and traditional material complete its eventual ACK recipe.
The fixed-candidate reduction can stop with zero on this branch.
After $J_r=j$, the global pre-return stop preserves the available prefix
and latched bit. These certificates use public bindings and immutable
provenance, so they are executable before any unknown bytes are needed.

For fixed $(r,j,b)$, use the two-stage operational ML-KEM game of
Section~\ref{sec:contextbound-connection}. Associate its public key with
the original generation slot $g$; represent its PQ tape by an opaque handle
and generate the traditional leg locally. At $a$'s first private
encapsulation, check its fixed peer binding, request the challenge once,
and cache the ciphertext and challenge key before later persistence.
Selected-owner decapsulations use the prechallenge oracle before that
call. Afterwards its exact challenge ciphertext uses the cached key,
and other ciphertexts use the permitted decapsulation oracle. Other
owners keep their local keys. Native policy, token, body and traditional
rejection checks remain in their actual order.

\begin{lemma}[Fresh encapsulation and restored-owner validation]
\label{lem:bootstrap-validation}
Let $v_j$ bound all other original fresh encapsulation-input slots under
the selected PQ public-key bytes, including unused and failed-retention
draws and restored-owner validation draws. Let
$\delta_{{\rm PKE},j}(v_j+1)$ bound a K-PKE correctness failure among
these encapsulations and the selected one. The probability that another
honest fresh encapsulation has the selected challenge ciphertext is at most
\begin{equation}
 \nu_{{\rm enc},j}=\min\{1,v_j2^{-256}
                   +\delta_{{\rm PKE},j}(v_j+1)\}.
 \label{eq:bootstrap-validation}
\end{equation}
\end{lemma}
\begin{proof}
ML-KEM encrypts the original 32-byte input $m$ using its specified
K-PKE coins~\cite[Algorithms~14, 15, 17 and 18]{fips203}. Before a
correctness failure, equal ciphertexts decrypt to equal values of $m$.
For each other fixed sampling slot, the later of its draw and the
selected draw is independent uniform, giving equality probability
$2^{-256}$. A union bound over the fixed lifetime slots proves the
claim. A sampled but selectively unexecuted reservation remains counted.
\end{proof}

The adapter caches earlier honest fresh ciphertexts and checks them at
challenge creation; it checks later ones before their results are used.
This covers other response contexts and expanded-owner restoration.
The latter performs fresh encapsulation and an expected/recovered-secret
comparison. On a nonmatching ciphertext the ordinary decapsulation oracle
therefore preserves the actual validation; a matching ciphertext incurs
\eqref{eq:bootstrap-validation}. Exact replay of $a$ reuses its cached
challenge. Write $a_{rj,b}$ for this complete component bit-comparison
advantage, and $\delta_j$ for the source correctness cost, including the
original Initial contribution, selected reply pair and imported/generated
owner checks. The random component enters the complete ContextBound
comparison, costing $\Phi_{256}(d_j,Q_{{\rm CB},rj,b})$ with $d_j=1$
for this source-consistent interface. Its distinguished table output
then realizes the hidden-second service of
Theorem~\ref{thm:bootstrap-joint-entry}.

\paragraph{Immutable records in the comparison.}
Mark the conservative union of private record versions containing a
value represented by an opaque handle at any stage of this reduction:
the original tapes, expanded owner, component/contribution, PRK, seed,
and stored root or activation values. Containing images, pending
intents and later source-preserving copies inherit the mark. Public
confirmation tags and fingerprints remain ordinary available outputs.
A marked version stays marked after deletion or later opening.
Let $C_j$ be successful raw-byte acquisition of any such object and put
\begin{equation}
 \eta_{{\rm cap},r}=\sum_{b=0}^{1}
      \Pr_{H_{r,b}}[J_r\ne\bot\ \land\ C_{J_r}].
 \label{eq:bootstrap-capture}
\end{equation}
Each hybrid uses its own history. Zeroing the candidate bit before a
marked return changes its expectation by at most
$\Pr_{H_{r,b}}[J_r=j\land C_j]$; summing pairs gives
\eqref{eq:bootstrap-capture}.

On the remaining branch, apply Corollary~\ref{cor:real-to-opaque}
while real tapes are still known: first structural private-wire equality,
then image/intent commitment replacement, and then the opaque renderer.
Only afterwards embed the unknown component tape. A valid record cost is
\[
 \rho_{rj,b}=\min\left\{1,
   \sum_{i=1}^{M_j}\Phi_{192}(d_i,Q_{i,rj,b})
                  +\binom{M_j}{2}2^{-192}\right\},
 \qquad d_i\le1+5K_i^{\rm hdr}.
\]
All fresh seals, including failed writes, enter $M_j$; exact retries
reuse their original invocation. The nonce term is omitted when the
interface has no private-wire equality query. The per-seal nonce remains
private before the declared raw acquisition stop. This is the actual
acquisition interface of the constructive record theorem.

Thus the branch with no marked raw acquisition has
$\eta_{{\rm cap},r}=0$. The general acquisition interface includes the
explicit residual~\eqref{eq:bootstrap-capture}. Replacing it by a
cryptographic record cost requires a joint comparison for the native
journal key, its directly encrypted images and its derived operation
and signer wrapping keys, preserving ciphertext-dependent commitments,
exact copies and later openings. Guarded readability determines which
observations reconstruct a source; the record comparison supplies the
corresponding hidden-byte simulation.

\paragraph{From effective descriptors to protected classes.}
For an incoming independent opaque PRK family $p$, the root descriptor
is $e=(p,\Gamma(s_c),I_B)$. Its seed's only raw uses are the root-entry
Extracts; each PRK's raw uses are its
three labelled Expand calls. We therefore replace the parent row
interfaces by private tables before embedding their child keys. The
source stores $K_B$ in its pending state before activation, leaving
the interval in which online registration can take place.

\begin{lemma}[Deferred roots and full-wire classes]
\label{lem:bootstrap-deferred}
An independent private root-table entry accessed only through its
identity before first raw return or activation can be sampled at that
first exposure. At a protected first activation it supplies a fresh
hidden challenge root. Before $\mathsf{Coll}$, two distinct complete
classes in one PRK family sharing a root descriptor give a collision
in the public compression oracle. With a complete query bound $Q_h$,
that event has probability at most
$\chi_h(Q_h)=\min\{1,C_{\rm col}(Q_h+2)^3 2^{-256}\}$.
\end{lemma}
\begin{proof}
Before exposure the entry is an unused independent random register;
all available classical and quantum computation is independent of its
value. Its sampling commutes with that prefix. At the first exposure,
draw a local value for a raw return or unprotected activation, or assign
the next hidden challenge for a protected activation. Exact descriptor
aliases retrieve the same value; independent entries may coincide.
This argument applies after the record and full-root-table comparisons,
which supply the value-independent handle interface.

Equal $I_B$ implies equal contexts. If the distinct classes also had
equal $s_c$, the collision-free complete-wire hashing would equate
$x,y$. Their salts therefore differ. The injective pad map $I$ gives
two different inputs $(I(s_c),\mathrm{IV})$ with equal $\alpha$.
The random-function collision bound~\cite[Theorem~7]{zhandryCollision2015}
applies to this witness with all honest work, monitoring and extraction
in $Q_h$. One global finder suffices, without a pairwise transcript factor.
\end{proof}

Use $\Delta_{{\rm ent},p,b}$ from
Lemma~\ref{lem:entry-information-length} with information length 93.
For a bound $N_p$ on family $p$'s distinct root descriptors, define
\begin{equation}
 \begin{split}
 D_b=\min\left\{1,\sum_p\left[
   \varepsilon_{\rm Hrow}(3,q_{B,p,b},71)
     +\Delta_{{\rm ent},p,b}\right.\right.\\
 \left.\left.\hspace{12mm}
  +\sum_{i=1}^{N_p}\bigl(
      \varepsilon_{\rm Hjoint}(1,0,q_{A,pi,b},32)
      +\varepsilon_{\rm Hrow}(5,q_{T,pi,b},144)\bigr)
       \right]\right\}.
 \end{split}
 \label{eq:bootstrap-tail}
\end{equation}
Every query budget covers the complete derived experiment, including
unprotected activations, other families, confirmations, allowed sibling
observations and final testing, uniformly across its hybrids.

\begin{proposition}[ACK-only comparison from opaque PRKs]
\label{prop:bootstrap-ack-table}
For the incoming opaque services and online experiment above, the
ACK-only output-bit gap is at most
$\min\{1,D_0+D_1+\chi_h(Q_{h,0})+\chi_h(Q_{h,1})\}$.
\end{proposition}
\begin{proof}
In each hybrid, stop at the descriptor-collision monitor. Replace its
three bootstrap outputs by a private row and its root entries by their
complete-descriptor table. By Lemma~\ref{lem:bootstrap-deferred}, a
protected first activation now receives the next independent hidden
root. Apply one Extract-column and one five-message-row comparison,
then the exact prefix-consistent Expand-table identity. This yields
five independent cached blocks per protected class at cost $D_b$;
unprotected activations still use real computation under their current roots.

After both hybrids have been converted to full tables, couple the fourth
block to the independent replacement ACK. The other four blocks have the
same joint distribution.
The overwritten table value has no remaining internal use, because the
Expand-table identity already removed the raw block-5 row input.
Class aliases, registration, stops and continuation agree. The triangle
inequality through these common-table experiments gives the stated bound.
\end{proof}

\begin{theorem}[Original reply-source schedule connection]
\label{thm:bootstrap-source-schedule}
For the surviving-source profile above, set
\[
 \begin{split}
 L_{rj,b}={}&\rho_{rj,b}+a_{rj,b}+\delta_j
       +\Phi_{256}(d_j,Q_{{\rm CB},rj,b})\\
       &+\Delta_{{\rm bs\text{-}entry},rj,b}
        +D_{{\rm tail},rj,b}+\chi_h(Q_{h,rj,b}).
 \end{split}
\]
Here $D_{{\rm tail},rj,b}$ is the appropriate cost for one hybrid
of~\eqref{eq:bootstrap-tail}. The complete online ACK schedule satisfies
\begin{equation}
 \begin{split}
 \varepsilon_{\rm bs\text{-}schedule}\le\min\left\{1,
 \sum_{r=1}^{R}\left[
       \alpha_{r,0}+\alpha_{r,1}+\eta_{{\rm cap},r}\right.\right.\\
 \left.\left.\hspace{16mm}
       +\sum_{j\in[G]\times[A]}\sum_{b=0}^{1}
              (\nu_{{\rm enc},j}+L_{rj,b})\right]\right\}.
 \end{split}
 \label{eq:bootstrap-source-schedule}
\end{equation}
Uniform valid bounds give the resource factor $2RGA$ on
$\nu_{\rm enc}+\rho+a+\delta+\Phi_{256}
 +\Delta_{\rm bs\text{-}entry}+D_{\rm tail}+\chi_h$,
plus $2R\alpha+\sum_r\eta_{{\rm cap},r}$, clipped at one.
\end{theorem}
\begin{proof}
Telescope the adjacent ACK hybrids and charge their two association
events. Cancel absent registrations and decompose each remaining
terminal bit by its original pair. The pre-opening certificates add
zero; marked acquisition contributes~\eqref{eq:bootstrap-capture},
and fresh encapsulation collisions contribute
\eqref{eq:bootstrap-validation}. In each fixed-candidate hybrid,
replace constructible records, install the original component adapter,
apply the complete-tuple comparison and the hidden-second entry, and
then use the comparison for the corresponding hybrid in
Proposition~\ref{prop:bootstrap-ack-table}.
Both resulting full-table experiments couple with the same terminal
code $D\mathbf1[J_r=j]$, their own source ledgers and registration
history, and the same used ACK entry. Their candidate expectations
are equal. Summing their two costs and then the fixed pairs and ranks
proves~\eqref{eq:bootstrap-source-schedule}.
\end{proof}

\begin{theorem}[Historical responder activation]
\label{thm:bootstrap-terminal}
Let $\mathsf{Bad}$ be the latched event that a registered initiator
connection succeeds without an earlier durable activation of the same
$c$ under its expected responder identity. Let $q_v<2^{256}$ be a
deterministic bound, valid in every comparison below, on the number of
distinct triples $(c,m,t)$ whose verification reaches the keyed ACK
comparison before the first matching responder activation of $c$.
Here $c$ is the registered complete class, $m$ is the complete 107-byte
authenticated ACK string, and $t$ is the proposed 256-bit tag. Exact
class aliases share $c$, and a repeated identical triple is counted once.
Using the same tag with a different class or authenticated string counts
as a different triple. A keyed comparison followed by a floor rejection
also enters this count. Then
\begin{equation}
 \Pr[\mathsf{Bad}]\le\min\left\{1,
  \Pr[\mathsf{Coll}]+\varepsilon_{\rm bs\text{-}schedule}
  +E_{\rm bs\text{-}mac}
  +\frac{q_v}{2^{256}-q_v+1}\right\}.
 \label{eq:bootstrap-terminal}
\end{equation}
In the ACK-only ideal experiment $H_R$, mark immutable versions containing
any registered replacement ACK key, including containing images, pending
intents and exact copies. Marks retain their original typed provenance
through deletion; public ACK tags are ordinary row outputs. Let
$C_{\rm ACK}$ be raw acquisition of any marked object and
$\eta_{\rm mac}=\Pr_{H_R}[C_{\rm ACK}]$. If $\rho_{\rm mac}$ is the
real-to-opaque image/intent cost for this union of marked versions, then
an explicit one-sided terminal comparison cost is
\begin{equation}
 E_{\rm bs\text{-}mac}=\min\left\{1,
   \eta_{\rm mac}+\rho_{\rm mac}
   +\sum_{r=1}^{R}\varepsilon_{\rm Hrow}(T_r,q_r,107)\right\}.
 \label{eq:bootstrap-mac-records}
\end{equation}
The record cost has the displayed $\Phi_{192}$ and nonce-collision
form, with all fresh seals counted. Each $T_r$ counts honest keyed
and verification evaluations, including a valid tag later rejected by
the floor check; each $q_r$ covers its complete derived experiment.
For an interface admitting no marked raw acquisition,
$\eta_{\rm mac}=0$.
\end{theorem}
\begin{proof}
Charge the framed-hash collision, then apply
Theorem~\ref{thm:bootstrap-source-schedule}. In $H_R$, zero the local
terminal bit before a marked raw return; its expectation changes by at
most $\eta_{\rm mac}$. Start with known sampled replacement keys and
construct real records. Corollary~\ref{cor:real-to-opaque} replaces
private-wire equality and then commitments before installing its opaque
renderer, at cost $\rho_{\rm mac}$. The experiment's global recipe
stop continues to preserve the latched event; the local acquisition
comparison is charged separately.

In this opaque-record experiment, replace registered uniform-key MAC rows
in registration order using Lemma~\ref{lem:appendix-hmac-instances}.
Each selected key is unused until its registration, and its only raw
computation in the defined operation surface is the 107-byte ACK row.
Other keys and their real operations remain in the complete continuation.
Exact retries cache the same message; numerical key coincidences are
handled through the shared oracle as before. The last experiment has opaque
records, independent functions on the authenticated strings and its local
capture-zero rule. This establishes the one-sided cost in
\eqref{eq:bootstrap-mac-records}.

Before the first matching responder activation, that responder has no
admitted role-2 message state that can release the direction-1 tag.
The initiator's ordinary direction-2 ACK has a separately designated key and
direction string. The client checks the exact session, direction, epoch-zero
encoding and tag before connection success. Write $W_c$ for the independent
MAC function of registered class $c$, and consider the first pre-activation
keyed equality $t=W_c(m)$, including one followed by a floor rejection.
Before this equality, an exact repeated triple has the same negative
outcome. For each new triple, prior negative comparisons at $(c,m)$ exclude
only the tags already proposed at that pair; other class/message pairs
have independent table entries. The unreleased value $W_c(m)$ is therefore
uniform among at least $2^{256}-q_v+1$ values conditional on those misses.
Summing the conditional first-hit probabilities over at most $q_v$
distinct triples gives the last term. Connection success before matching
activation entails such a hit, and the local capture-zero rule only removes
executions from this event. Replaying a previously authentic READY proof
already supplies the matching historical activation. Later global recipe
stops preserve the event observed on the earlier prefix.
\end{proof}

A READY proof remains evidence of the expected responder's earlier
durable activation of the same signed transcript under the theorem's
conditions. It may arrive after the initiator restarts; the initiator
checks its current local authority when returning the connection.
An exact replay of READY refers to the same activation event. The supplement
records the schedule and durable-order pins, independent HMAC equations,
short/long alias controls, late registration countermodels and
acquisition accounting in both comparison hybrids.

\subsection{Initial-source authentication with adaptive reply families}
\label{sec:bootstrap-initial-source}

Now suppose the Initial supplies the surviving secret. Before a Reply
arrives, the initiator has already used that contribution for
confirmation. Different replies can supply different second
contributions and thus different inputs to the long Extract. We first
compare the original source's complete service, including both uses,
and only then restrict the private table to the parent family of the
selected registration. In this order the source and entry comparisons
are charged once per original pair, however many reply families arise.

\paragraph{Initial-surviving profile.}
Use Definition~\ref{def:bootstrap-online}, adding the authoritative
inventory's publication and recovery operations. Each protected class
starts from an honest durable Initial whose admitted context selects a
PQ leaf from the expected responder's authoritative inventory. Context
admission and that selected leaf remain fixed under later signing-key
disclosure. The initiator's BootstrapInitiator signing owner is protected
and supplies purpose-5 signatures through this initiation history. The
responder's signing owner, all traditional scalars, and the initiator's
reply-private owners are available from their creation; each actual
second contribution is available at its computation. Thus an adversarial
Reply retains its actual signature, decoding, decapsulation, confirmation
and rejection behavior.

The selected PQ owner's private-use interface consists of inventory
generation/restoration and Initial decapsulation. The observation surface
comprises the bootstrap/ACK outputs, public image heads and the five
pending-intent digest domains, specified source/key and raw-record
disclosures, and the activation siblings $R_0,C_I,C_R,A_R$ with the uses
of Section~\ref{sec:bootstrap-authentication}. Further authenticated
reports over private state enter through
Section~\ref{sec:report-composition}'s additional report interface.
The original recipe rule governs the selected PQ generation and
encapsulation sources, their component key and $S_0$. It retains
acquired versions and prior openings, and preserves an already latched
event at a subsequent global recipe stop. Initial-confirmation material
follows its actual joint-entry interface and may remain available while
$S_0$ is protected.

Fix complete-execution bounds $G,A,N,R$ on inventory PQ-generation
reservations, Initial-encapsulation reservations, Initial-nonce sampling
slots and protected registrations. Count failed-retention and unused
draws. For each original pair $j=(g,a)\in[G]\times[A]$, fix a bound
$L_j$ on its distinct long descriptors, including computations followed
by failed verification or persistence. These bounds hold in every derived
experiment. Exact retries retain their original slots.

\begin{lemma}[Inventory origin and the complete initial service]
\label{lem:bootstrap-initial-origin}
Before the initiator fresh-body signature, relevant framed-hash,
Initial-nonce and fresh-encapsulation collision events, and on source
correctness, a protected class identifies an original pair $j=(g,a)$.
That pair fixes one complete ContextBound tuple and $S_{0,j}$. Its
initial confirmation and all later bootstrap Extracts realize
\begin{equation}
 \begin{aligned}
 P_{0,j}&=F(\sigma_{32},S_{0,j}),\\
 P_{B,j}(t_1,v)&=F(t_1,S_{0,j}\Vert v),
       \qquad v=\text{the actual available }S_1.
 \end{aligned}
 \label{eq:bootstrap-initial-service}
\end{equation}
The complete long descriptor is
$e=(v,\alpha(t_1),\beta(t_1))$. Distinct descriptors are ranked at their
first actual long request, with that rank retained across failed checks,
disclosure, deletion and exact retry.
\end{lemma}
\begin{proof}
Inventory generation commits the reservation before invoking the saved
generation tape and commits the public leaf before publication. Restoration
regenerates that owner from its original sealed seed, checking its embedded
public key and its equality to the selected published component. The
authoritative inventory rejects duplicate published component fingerprints.
It retains the selected PQ and traditional references separately; they may
belong to different hybrid owners. One-time consumption changes a later
record to an operation tombstone and preserves the executed origin.

Durable initiation commits its Initial-encapsulation reservation before
private computation. Its saved nonce, reply key, peer, context and tape
determine the ciphertext and all ten ContextBound fields. Restored
completion requires equality with the saved signing body. The responder
verifies purpose 5 before Initial decapsulation. Outside the stated
association events, its issued-body ledger and checked bindings therefore
identify the same original pair. Source correctness and the
fresh-encapsulation monitor give the same complete tuple at admitted
selected-challenge uses.

The Initial computes its SDK short Extract, initial Expand and MAC before
signing. The later schedule concatenates this same $S_{0,j}$ with the
actual $S_1$ and uses the 32-byte Reply-core digest $t_1$ as its salt, giving
\eqref{eq:bootstrap-initial-service}. A same-body Initial signature alias
preserves the short service while its complete wire affects later Replies.
A Reply signature alias retains its long descriptor and its actual
complete-wire root binding. The joint service therefore preserves both
the first-confirmation history and all later adaptive requests.
\end{proof}

For endpoint $(r,b)$, a valid association cost is
\[
 \alpha^{I}_{r,b}=\min\left\{1,
   \varepsilon_{{\rm sig}I,r,b}+\chi_H(Q_{H,r,b})
       +\binom N2 2^{-256}\right\}.
\]
Here the purpose/body term uses
Theorem~\ref{thm:native-mldsa-reservation-bridge} for the protected
initiator. One complete framed-hash finder covers prekey fingerprints,
contexts and the Initial/Reply domain bindings. The source-correctness
and fresh-encapsulation terms are charged below, separately from this
association cost.

\paragraph{First comparison: the entire original-source service.}
Use the adjacent ACK experiments $H_{r,0},H_{r,1}$ of
Section~\ref{sec:bootstrap-authentication}. After the association stop,
let $J_r$ be the Initial pair of the $r$th registration, or $\bot$.
For an available terminal bit $D$, the absent-registration terms cancel,
and the remaining difference partitions into
$D\mathbf1[J_r=j]$. Opening the selected original source or $S_{0,j}$
before registration excludes every descendant of that source: its
already available traditional material and eventual known $S_1$ complete
the ACK recipe. A fixed-pair reduction then returns zero before those
bytes are needed. Opening one long parent excludes only that parent's
descendants; other long families retain their ordinary interfaces.
After registration the global pre-return rule preserves the latched bit.

Mark the immutable union of private versions containing the selected
tapes and every value represented by an opaque handle during any later
parent comparison, including inventory, plan, checkpoint, image, intent
and containing copies. This fixed typed-dependency closure retains its
marks through every comparison. Let $C_j^I$ be raw acquisition of this union.
The pair partition gives the exact two-endpoint acquisition term
\begin{equation}
 \eta^I_{{\rm cap},r}
   =\sum_{b=0}^{1}\Pr_{H_{r,b}}
          [J_r\ne\bot\ \land\ C^I_{J_r}].
 \label{eq:bootstrap-initial-capture}
\end{equation}
The union covers every later family, so this term is charged before
partitioning long-table allocations. Each endpoint uses its own history.
On the remaining branch, apply the real-to-opaque record comparison
while tapes are known, at the cost $\rho^I_{rj,b}$ of
Section~\ref{sec:bootstrap-authentication} for these marked versions.

At reservation $g$, register an opaque 64-byte PQ generation tape; its
first generation uses the operational ML-KEM game's public key. Install
the challenge at $a$'s first encapsulation, using an opaque handle for
its 32-byte PQ input. Their independent traditional scalars are sampled
and available at the original reservations.
Before first private use, an unopened opaque tape has no value-dependent
observation, so its independent register may be supplied at that first use.
This deferral preserves reservation times and failed-draw counts. Cache
the challenge ciphertext/key before later persistence, and retain the
prechallenge and permitted postchallenge decapsulation interfaces. Other
owners retain their locally generated keys and actual algorithms.
The purpose-5 origin check restricts selected-challenge returns to the
one complete tuple and its retries.

Inventory restoration uses seed regeneration with the checks in
Lemma~\ref{lem:bootstrap-initial-origin}. The distinct reply owner is
locally known and retains its actual expanded-import validation. Any
admitted auxiliary expanded import of the selected owner retains the
fresh encapsulation comparison of
Lemma~\ref{lem:bootstrap-validation}. The same lemma supplies
$\nu^I_{{\rm enc},j}$ for all other original fresh encapsulation-input
slots under the selected public-key bytes. Count abandoned as well as
executed slots. Let $a^I_{rj,b}$ be the complete operational component
advantage and $\delta^I_j$ include Initial-pair correctness,
generated/imported-owner consistency and correctness of honest Reply
pairs needed for a shared endpoint class. Arbitrary submitted Reply
ciphertexts retain their actual decapsulation result.

The random component branch supplies the complete-tuple comparison at
cost $\Phi_{256}(d_j,Q^I_{{\rm CB},rj,b})$, with $d_j=1$ for this
source-consistent interface. Its distinguished output realizes
$S_{0,j}$ in~\eqref{eq:bootstrap-initial-service}. Apply the entire
hidden-first comparison of
Theorem~\ref{thm:bootstrap-joint-entry}. This replaces the short PRK
and every long PRK by an independent short entry $V_0$ and a private
table $V_B(v,\alpha,\beta)$, at cost
$\Delta^I_{{\rm bs\text{-}entry},rj,b}$. It retains the shared raw
point of the short call and a long call with known half $P$, all actual
initial confirmation operations, later family openings and the complete
continuation. Denote this transformed endpoint by $T_{rj,b}$ and put
\begin{equation}
 \begin{split}
 E^I_{rj,b}={}&\nu^I_{{\rm enc},j}+\rho^I_{rj,b}+a^I_{rj,b}
                     +\delta^I_j\\
             &+\Phi_{256}(d_j,Q^I_{{\rm CB},rj,b})
                     +\Delta^I_{{\rm bs\text{-}entry},rj,b}.
 \end{split}
 \label{eq:bootstrap-initial-upstream}
\end{equation}
Each query bound covers the entire experiment in its replacement,
including all later long requests. The source comparison therefore pays
for that complete continuation once, before allocation ranks separate
the possible parent families.

\begin{lemma}[Allocation-local parent comparison]
\label{lem:bootstrap-initial-allocation}
In $T_{rj,b}$, let $K_r$ be the first-allocation rank of the long
descriptor used by registration $r$ when $J_r=j$. The candidate bit
has the pointwise decomposition
\[
 D\mathbf1[J_r=j]=\sum_{\ell=1}^{L_j}
       D\mathbf1[J_r=j\ \land\ K_r=\ell].
\]
For fixed $\ell$, its unused table cell can receive a presampled
independent opaque parent at its first request. Its terminal-tail
comparison costs
$T^I_{rj\ell,b}=D^I_{{\rm tail},rj\ell,b}
                    +\chi_h(Q^I_{h,rj\ell,b})$.
Here $D^I_{{\rm tail},rj\ell,b}$ is the one-parent endpoint
specialization of~\eqref{eq:bootstrap-tail}, including that parent's
complete family of downstream root descriptors.
\end{lemma}
\begin{proof}
The complete descriptors partition the long requests, and immutable
first-allocation ranks partition their possible selected parent. The
private table supplies an unused independent value at each first request,
even when that request depends on all earlier confirmations, the short
entry and other long entries. A fixed rank can therefore receive its
presampled challenge then; a never-allocated rank contributes zero.
The table retains numerical coincidences between independent values.

Opening that parent, its seed or a complete descendant recipe before
registration makes this fixed rank's candidate bit zero under cumulative
eligibility. Another family from the same source may still register at
its different rank. After registration, the global recipe stop retains
the available prefix and latched bit. These certificates allow the
one-family endpoint of Proposition~\ref{prop:bootstrap-ack-table},
retaining every other family and the short entry in the auxiliary
controller. It admits confirmations and pending-root retention before
registration, defers unopened root entries, and preserves all five
activation blocks before the separate ACK export replacement. Its
complete-descriptor collision monitor contributes the stated
$\chi_h$ term.
\end{proof}

\begin{theorem}[Two-level initial-source schedule connection]
\label{thm:bootstrap-initial-schedule}
For the initial-surviving profile, the online ACK schedule satisfies
\begin{equation}
 \begin{split}
 \varepsilon^I_{\rm bs\text{-}schedule}\le\min\Bigl\{1,
  \sum_{r=1}^{R}\Bigl[\alpha^I_{r,0}+\alpha^I_{r,1}
       +\eta^I_{{\rm cap},r}\\
 \hspace{10mm}+\sum_{j\in[G]\times[A]}\sum_{b=0}^{1}
        \Bigl(E^I_{rj,b}+\sum_{\ell=1}^{L_j}T^I_{rj\ell,b}\Bigr)
       \Bigr]\Bigr\}.
 \end{split}
 \label{eq:bootstrap-initial-schedule}
\end{equation}
The budgets and $L_j$ may differ across original pairs. With common
upper bounds and $L_j\le L$, the expression is at most
\begin{equation}
 \begin{split}
 \min\Bigl\{1,\;2R\alpha^I+\sum_r\eta^I_{{\rm cap},r}
   +2RGA\bigl[E^I+L T^I\bigr]\Bigr\}.
 \end{split}
 \label{eq:bootstrap-initial-uniform}
\end{equation}
\end{theorem}
\begin{proof}
Telescope ACK registrations and charge the endpoint association events.
Cancel absent registrations, then localize to each original pair.
The immutable acquisition modification gives
\eqref{eq:bootstrap-initial-capture}; the record, component,
complete-tuple and whole-service entry comparisons cost
\eqref{eq:bootstrap-initial-upstream} once for each endpoint and pair.

Partition the resulting candidate bit in $T_{rj,b}$ by its online
allocation rank. Lemma~\ref{lem:bootstrap-initial-allocation} applies
the terminal-tail comparison at each fixed rank. The two full-table
endpoints now have the same law: the selected ACK is its independent
class entry, other classes retain their ordinary computations or earlier
replacements, and the common terminal code recomputes source/rank flags,
registration and stops from that endpoint's history. Their candidate
expectations agree. The triangle inequality and sums over $\ell,j,r$
give~\eqref{eq:bootstrap-initial-schedule}. The second partition occurs
after the entire source and joint-entry comparison, so only the subsequent
tail costs carry the additional sum over $\ell$. Common upper bounds give
\eqref{eq:bootstrap-initial-uniform}.
\end{proof}

\begin{corollary}[Historical activation from the initial source]
\label{cor:bootstrap-initial-terminal}
For the initial-surviving profile and the latched event of
Theorem~\ref{thm:bootstrap-terminal}, let $q_v<2^{256}$ bound the distinct
pre-activation verification attempts defined in that theorem, with their
registered class, complete authenticated ACK string and proposed tag. Then
\[
 \Pr[\mathsf{Bad}]\le\min\left\{1,
  \Pr[\mathsf{Coll}]+\varepsilon^I_{\rm bs\text{-}schedule}
   +E^I_{\rm bs\text{-}mac}
   +\frac{q_v}{2^{256}-q_v+1}\right\}.
\]
Here $E^I_{\rm bs\text{-}mac}$ is
\eqref{eq:bootstrap-mac-records}, evaluated in this profile's ACK-only
endpoint with its own immutable marks, acquisition probability and
whole-execution budgets.
\end{corollary}
\begin{proof}
Apply Theorem~\ref{thm:bootstrap-initial-schedule} to reach the same
ACK-only interface. The opaque-record and 107-byte-row comparisons in
Theorem~\ref{thm:bootstrap-terminal} then apply with the stated
profile's observations. Before matching durable responder activation,
no admitted responder ACK operation releases that registered direction-1
tag. Its fresh-guess calculation and exact-replay history give the bound.
\end{proof}

The construction is pinned to the inventory, staged initiation and
complete-tuple source at commit \texttt{471900bf}. Its complete source
map, analytical reduction and retained finite controls appear in the
\href{https://github.com/billlza/q-periapt/blob/81483b52a24a4cc81c2276c5a0d0b667c42761cd/docs/continuity/BOOTSTRAP_INITIAL_SOURCE_ADAPTER_V1.md}{initial-source supplement}.
Both source profiles therefore authenticate historical responder
activation under their respective original-operation eligibility
conditions. For the initial source, comparing its complete service first
also preserves adaptive reply families while charging the source and
joint-entry costs once per original pair.

